% 本文件由 LaTeX 排版 Agent 根据有效 translation.json 自动生成。
% author=Justin Martyr; work_id=0128; title=与特黎丰对话（犹斯定殉道者）

\chapter{与特黎丰对话（犹斯定殉道者）}

% structure: p0001 source=h1 target=omitted reason=page-title-already-rendered-by-parent

\begin{itemize}
\item \Emph{第1–9章}
\end{itemize}

\begin{itemize}
\item \Emph{第10–30章}
\end{itemize}

\begin{itemize}
\item \Emph{第31–47章}
\end{itemize}

\begin{itemize}
\item \Emph{第48–54章}
\end{itemize}

\begin{itemize}
\item \Emph{第55–68章}
\end{itemize}

\begin{itemize}
\item \Emph{第69–88章}
\end{itemize}

\begin{itemize}
\item \Emph{第89–108章}
\end{itemize}

\begin{itemize}
\item \Emph{第109–124章}
\end{itemize}

\begin{itemize}
\item \Emph{第125–142章}
\end{itemize}

\SourceRecord{Translated by Marcus Dods and George Reith.；Ante-Nicene Fathers；Vol. 1.；Edited by Alexander Roberts, James Donaldson, and A. Cleveland Coxe.；Buffalo, NY: Christian Literature Publishing Co.,；1885.；<http://www.newadvent.org/fathers/0128.htm>.}

% structure: p0001 source=h1 target=section reason=page-title

\section{与特黎丰对话录（第1-9章）}

% structure: p0002 source=h2 target=subsection reason=relative-heading-rank; base=h2

\subsection{第一章 引言}

一日早晨，我在\OriginalTerm{柱廊}{Xystus}的步道间散步时，有一个人，名叫\Person{特黎丰}，带着几个同伴，遇见了我。\par

\Emph{\Person{特黎丰}：}\Quote{问候你，哲学家！}\par

说完这话，他立即转身与我同行；他的朋友们也跟随着他。\par

\Emph{\Person{尤斯定}：}\Quote{有什么要紧的事吗？}\par

\Emph{\Person{特黎丰}：}\Quote{过去在\Place{阿尔戈斯}，苏格拉底派的\Person{科林图斯}教导我，不可轻看或漠视那些穿着此等装束的人，反而要向他们显示一切善意，并与他们交往，因为或许从这种交往中，无论对这样的人还是对我自己，都会产生某些益处。总之，如果双方中有一方获得益处，那对两者都是好的。因此，每当我看见有人穿着这样的服装，我便欣然趋近他；如今，出于同样的理由，我也乐意前来与你交谈；这些同伴与我同来，是希望亲自从你那里听到一些有益的东西。}\par

\Emph{\Person{尤斯定}：}\Emph{（戏谑地）}\Quote{但是，最卓越的人啊，你是谁呢？}\par

于是他坦率地告诉了我他的名字和家族。\par

\Emph{\Person{特黎丰}：}\Quote{我叫特黎丰，是受割礼的希伯来人；我最近从那里发生的战争中逃出，现在在\Place{希腊}度日，主要在\Place{科林多}。}\par

\Emph{\Person{尤斯定}：}\Quote{然而，您从哲学中获得的益处，又怎能比得上从您自己的立法者和先知那里所得的呢？}\par

\Emph{\Person{特黎丰}：}\Quote{为何不可？哲学家们岂不是将一切议论都转向天主？他们不是持续不断地提出关于天主的惟一性和眷顾的问题？调查神性，这岂不就是哲学真正的职责吗？}\par

\Emph{\Person{尤斯定}：}\Quote{确实如此，我们也是这样相信的。但大多数人并没有考虑过天主是否惟一或有多位，他们是否关心我们每一个人，好像这种知识对幸福毫无贡献；不仅如此，他们甚至试图说服我们，天主照顾宇宙及其种类和物种，却不照顾我和你以及每个个别人，因为否则我们就不必日夜向他祈祷了。但理解这个结论并不难；因为坚持这些意见的人结果变得无所畏惧和言论放肆，任意言行，既不惧怕惩罚，也不期望从天主得到任何益处。他们怎能如此呢？他们断言同样的事情将永远重复发生；而且，我和你将会以同样的方式再次生活，既没有变得更好也没有变得更坏。但还有另一些人，他们以为灵魂是不朽的、非物质的，便相信即便犯了恶也不会受惩罚（因为非物质的东西是无感觉的），而且由于灵魂不朽，它不需要从天主那里得到任何东西。}\par

\Emph{\Person{特黎丰}：}\Emph{（微微一笑）}\Quote{请告诉我们您对这些事情的看法，您对天主持何种观念，以及您的哲学是什么。}\par

% structure: p0015 source=h2 target=subsection reason=relative-heading-rank; base=h2

\subsection{第二章 尤斯定叙述他研究哲学的经过}

\Signature{犹斯定}：我愿告诉你我所想的；因为事实上，哲学是最伟大的财富，在天主面前最受尊崇，它引领我们归于他，且唯独它荐举我们；而那些致力研究哲学的人，确实是圣人。然而，哲学是什么，以及为何被赐予人类，这大多数人都未能认识；因为既不会有柏拉图学派、斯多葛学派、逍遥学派、理论派，也不会有毕达哥拉斯学派，这知识原是\Emph{一}。我愿告诉你它为何变得多头多脑。发生的情况是：最初涉足它（即哲学）的人，因而被视为杰出人物，但后继者却不探求真理，只钦佩前人的坚毅和自律，以及学说的新颖；并且每个人都认为从老师那里学来的就是真理；然后，这些后来者又将此等及类似的东西传给他们的后继者；这体系便以被称为该学说之父的人的名字命名。起初，我渴望亲自与其中一人交谈，便委身于一个斯多葛学派的人；与他相处相当长一段时间后，当我没有获得对天主的进一步认识时（因为他自己并不认识，并说这样的教导是不必要的），我便离开他，转而投靠另一个人，他被称为逍遥学派，且如\Emph{他}自以为的那样精明。此人，在最初几天款待我之后，要求我确定费用，以便我们的交往不致无利可图。因此我也离弃了他，认为他根本不算哲学家。但当我灵魂热切渴望听闻独特而卓越的哲学时，我来到一位毕达哥拉斯学派的人面前，他非常有名，一个非常看重自己智慧的人。然后，当我与他交谈，愿意成为他的听者和门徒时，他说：\Quote{那么怎样？你通晓音乐、天文学和几何学吗？如果你未曾先受那些使灵魂脱离可感知之物、使之适合心智之物的教导，并从而能沉思其本质上的可敬与本质上的美善，你怎能期望领会任何有助于幸福生活的事物呢？}在称赞了这些学科的许多分支，并告诉我它们是必要的之后，当我向他承认我的无知时，他打发我走了。因此，我相当不耐地接受了这事，正如当我希望落空时所能预料的，尤其因为我以为此人有些知识；但再次想到我将不得不花时间在这些学科上徘徊时，我无法忍受更久的耽搁。在我无助的情况下，我想起要与柏拉图学派的人会面，因为他们声名显赫。于是我尽可能多地与一位最近在我们城市定居的人在一起——一个睿智的人，在柏拉图学派中地位很高——我进步了，并且每日都有极大的提高。对非物质事物的领悟完全压倒了我，对理型的思索使我的灵魂长上了翅膀，以至于没过多久，我就以为自己已成为智者；我的愚钝竟至于此：我期望立即能看见天主，因为这就是柏拉图哲学的终极目的。\par

% structure: p0017 source=h2 target=subsection reason=relative-heading-rank; base=h2

\subsection{第三章 犹斯定讲述他归依的经过}

\Signature{犹斯定}：当我作此准备时，有一段时间我希望充满极大的宁静，并避开人迹，便常去一处离海不远的田野。一天，当我靠近那地方——到了那里后我打算独自一人——有一个老人，外表绝不可鄙，表现出温和而可敬的风度，在远处跟随着我。我停下脚步，转向他，目光颇为锐利地注视着他。\par

\Signature{老人}：你认识我吗？\par

\Signature{犹斯定}：不。\par

\Signature{老人}：那么，你为什么这样看我？\par

\Signature{犹斯定}：我惊讶，因为你碰巧和我同在同一个地方；因为我没想到会在这里见到任何人。\par

\Signature{老人}：我牵挂我家里的几个人。他们离开我了，所以我亲自来寻找他们，也许他们会在某处出现。但你为什么在这里？\par

\Signature{犹斯定}：我喜爱这样的散步，在这里我的注意力不会分散，因为与自己的对话不会中断；这样的地方最适合爱智之学。\par

\Signature{老人}：那么，你是个好辩者，却不爱行动或真理？你难道不更以成为实践的人为目标，而是成为诡辩家吗？\par

\Signature{犹斯定}：还有什么比这更伟大的事呢：揭示那治理万有的理性，抓住它，并骑乘其上，俯视他人的错误及其追求？但若没有哲学和正确的理性，就没有任何人能有智德。因此，每个人都必须从事哲学，并将此视为最伟大、最尊荣的工作；而其他事物只是次要或第三等重要的，虽然，确实，若它们依赖于哲学，便具有适度的价值，值得接纳；但若脱离哲学，不与哲学相伴，对于追求它们的人而言，便是粗俗而鄙陋的。\par

\Signature{老人（打断）}：那么，哲学能带来幸福吗？\par

\Signature{犹斯定}：当然，而且只有它。\par

\Signature{老人}：那么，什么是哲学？什么是幸福？请告诉我，除非有什么阻碍你开口。\par

\Signature{犹斯定}：那么，哲学就是对真实存在之物的认识，以及对真理的清晰洞察；幸福则是这种知识与智慧的酬报。\par

\Signature{老人}：但你称什么为天主？\par

\Signature{犹斯定}：那始终持守同一本性、同一方式，且是万有之原因者——那，就是天主。\par

于是我这样回答他；他愉快地听着，并继续询问我。\par

\Person{长者}：知识岂非用于不同事物的共通术语？因为在各种技艺中，凡通晓其中某一门者，皆被称为统帅术、治国术或医道方面的巧匠。但涉及天主之事与人间之事，则并非如此。是否存在一种知识，能使人兼具理解天主之事与人间之事，继而彻底认识神性与其中所含的义理？\par

\Person{犹斯定}：确是如此。\par

\Person{长者}：那么，我们认识人和天主，是否如同认识音乐、算术、天文或任何其他类似学科一般？\par

\Person{犹斯定}：绝非如此。\par

\Person{长者}：那么，你先前回答得不正确；因为有些\OriginalTerm{知识}{branch of knowledge}是通过学习或某种运用而获得的，另一些则是凭眼见认知。倘若有人告诉你，印度有一种动物，本性独一无二，形貌奇特、变化多端，那么你在亲眼见到之前，便无从知晓；除非听见过它的人描述，否则你也不能胜任任何说明。\par

\Person{犹斯定}：确实不能。\par

\Person{长者}：那么，哲学家们既未认识天主（他们从未见过祂，也未曾听过祂），怎能正确判断天主或讲论任何真理呢？\par

\Person{犹斯定}：但是，长者，天主不能像其他活物那样单凭肉眼看见，只能用心智把握，正如柏拉图所说，我也信服他。\par

% structure: p0044 source=h2 target=subsection reason=relative-heading-rank; base=h2

\subsection{灵魂凭自身不能看见天主}

\begin{quotation}
\Emph{}
\end{quotation}

\Person{长者}：那么，我们的心智是否具有如此巨大的能力？人难道不能凭感官更先察觉？人的心智若不蒙圣神教导，何时能看见天主？\par

\Person{犹斯定}：柏拉图的确说过：心智之眼具有这样的本性，并被赋予这个目的——当心智本身纯洁时，我们便能看见那万有的本体，祂是心智所认知的一切的起因，没有颜色，没有形状，没有大小——凡肉眼所见的任何事物都没有；祂超越一切本质，不可言说，不可解释，唯独是尊贵而良善的，突然进入处于良好状态的心灵，因它们与祂有亲缘关系并渴望见祂。\par

\Person{长者}：那么，我们与天主之间有什么亲缘关系？灵魂也是神圣而不死不灭的，且是那至高心智的一部分？如同那心智看见天主一样，我们也能在心智中领会天主，并因此成为幸福的吗？\par

\Person{犹斯定}：确是如此。\par

\Person{长者}：所有活物的灵魂都能理解祂吗？还是人的灵魂是一种，马和驴的灵魂是另一种？\par

\Person{犹斯定}：不。一切灵魂都是相似的。\par

\Person{长者}：那么，马和驴也能看见天主吗？或者说，它们曾在某个时候看见过天主？\par

\Person{犹斯定}：不，因为连大多数人也不能看见，除非那些生活正义、因义德和一切其他德行而得到净化的人。\par

\Person{长者}：那么，人看见天主，并非由于亲缘关系，也非因拥有心智，而是因为节制与正直？\par

\Person{犹斯定}：是的，也是因为他有那个用以领悟天主的官能。\par

\Person{长者}：那么，山羊或绵羊会伤害任何人吗？\par

\Person{犹斯定}：在任何方面都不伤害人。\par

\Person{长者}：那么，这些动物将看见\OriginalTerm{天主}{God}，按你的说法。\par

\Person{犹斯定}：不；因为它们的身体具有那样的本性，对它们构成了障碍。\par

\Person{长者}：如果这些动物能讲话，请放心，它们会更加嘲笑我们的身体。但现在暂且不谈这个，并按你所说的让步。然而，请告诉我：灵魂是在身体内时看见\OriginalTerm{天主}{God}，还是离开身体之后？\par

\Person{犹斯定}：只要灵魂处于人的形态，便有可能借助心智达到这一点；但尤其当它脱离身体、独处一隅时，便获得了它素来不断且全然爱慕的那一位。\par

\Person{长者}：那么，当灵魂再次回到人里面时，它还记得这事（即看见天主）吗？\par

\Person{犹斯定}：我看并非如此。\par

\Person{长者}：那么，那些见过\OriginalTerm{天主}{God}的人有什么益处？除非他记得见过的事实，否则见过的人与没见过的人有什么差别？\par

\Person{犹斯定}：我说不上来。\par

\Person{长者}：那些被认为不配看见这一景象的人又会受什么苦？\par

\Person{犹斯定}：他们被囚禁在某些野兽的身体内，这便是他们的惩罚。\par

\Person{长者}：那么，他们知道自己是因此缘故而处于那样的形态，并知道自己犯了罪吗？\par

\Person{犹斯定}：我不这样认为。\par

\Person{长者}：那么，他们从惩罚中得不到任何益处。而且我还要说，除非他们意识到惩罚，否则他们并未受罚。\par

\Person{犹斯定}：的确不是。\par

\Person{长者}：因此，灵魂既不能看见天主，也不能转生到其他身体内；因为那样他们就会知道自己在受罚，并且日后连最微小的罪也不敢犯。但是，我认为灵魂能感知天主的存在，并知道正义与虔敬是值得尊崇的，这一点我完全同意你。\par

\Person{犹斯定}：您说得对。\par

% structure: p0074 source=h2 target=subsection reason=relative-heading-rank; base=h2

\subsection{灵魂按其本性并非不死不灭}

\begin{quotation}
\Emph{}
\end{quotation}

\Person{长者}：这些哲学家对此一无所知；因为他们不能说明灵魂是什么。\par

\Person{犹斯定}：似乎如此。\par

\Person{长者}：也不应称之为不死不灭的；因为如果它是不死不灭的，那么它显然就是无始的。\par

\Person{犹斯定}：按照一些被称为柏拉图主义者的人的看法，灵魂既是无始的，也是不死不灭的。\par

\Person{长者}：你是否也说世界是无始的？\par

\Person{犹斯定}：有些人这样说，但我不同意他们。\par

\Emph{老人：}你说得对；有什么理由认为一个如此坚实的、具有抵抗力的、复合的、可变的、腐朽的、每日更新的身体，不是从某种原因产生的呢？但如果世界是受生的，灵魂也必然是受生的；并且也许在某个时候它们并不存在，因为它们是为人和其他生物创造的，如果你愿意说它们是完全分开受生的，而不是与它们各自的身体一起受生的。\par

\Emph{犹斯定：}这似乎是正确的。\par

\Emph{老人：}那么，它们不是不死的吗？\par

\Emph{犹斯定：}不；因为世界在我们看来是受生的。\par

\Emph{老人：}但我并不是说，所有的灵魂都会死；因为那对恶人来说确实是一件好事。那么怎样？虔诚者的灵魂留在一个更好的地方，而不义和恶者的灵魂则在更坏的地方，等候审判的时刻。因此，那些在天主眼中显得配得的灵魂永远不会死；但其他的灵魂则要受惩罚，只要天主愿意它们存在并受惩罚。\par

\Emph{犹斯定：}那么你所说的，是否类似于\Person{柏拉图}在\WorkTitle{蒂迈欧篇}中关于世界的暗示？他说世界确实是可朽坏的，因为它被创造，但不会因天主的旨意而解体或遭遇死亡。你是否认为同样的话也可以用于灵魂，以及一般来说所有事物？因为那些存在于天主之后，或任何时候将要存在的事物，具有可朽坏的性质，可能被抹去并停止存在；只有天主是非受生和不朽坏的，因此他是天主，而其他一切在他之后的都是受造的和可朽坏的。由于这个原因，灵魂既死也受惩罚：因为如果它们是非受生的，它们就不会犯罪，也不会充满愚昧，也不会胆怯，又凶残；它们也不会自愿变成猪、蛇和狗；如果它们是非受生的，那么强迫它们也的确是不公正的。因为非受生的事物与另一个非受生的事物相似、相等、相同；在能力或尊荣上，一个不应优于另一个，因此非受生的事物并不多：因为如果它们之间有什么差异，即使你搜索，也无法发现差异的原因；但当你让思想永远漫游到无限时，你最终会疲惫地落脚于一个非受生者，并说这是万物的原因。这些事难道没有逃脱\Person{柏拉图}与\Person{毕达哥拉斯}——那些对我们来说如同哲学之墙和堡垒的智者——的眼目吗？\par

% structure: p0088 source=h2 target=subsection reason=relative-heading-rank; base=h2

\subsection{第六章 这些事是\Person{柏拉图}和其他哲学家所不知道的}

\Emph{老人：}对我来说，无论\Person{柏拉图}、\Person{毕达哥拉斯}，或者总之任何其他人持有这样的观点，都无关紧要。因为真理就是如此；你会从这一点看出。灵魂肯定是生命或拥有生命。那么，如果它是生命，它会使别的东西，而不是它自己活着，正如运动会推动别的东西而非自己。没有人会否认灵魂是活着的。但如果它活着，它活着不是因为它是生命，而是因为它分有生命；但分有任何东西的东西，与它所分有的东西不同。现在灵魂分有生命，因为天主愿意它活着。因此，当天主不愿意它活着时，它甚至不会分有生命。因为活着不是它的属性，如同是天主的属性；正如人不会永远活着，灵魂也不会永远与身体结合，因为每当这个和谐必须被打破时，灵魂离开身体，人就不再存在；同样，每当灵魂必须停止存在时，生命之灵就被从它那里拿走，就不再是灵魂，而是回到它被取来的地方。\par

% structure: p0091 source=h2 target=subsection reason=relative-heading-rank; base=h2

\subsection{第七章 唯独应从先知寻求真理的知识}

\Emph{犹斯定：}那么，是否有人应该使用一位教师？或者，如果连他们之中也没有真理，人可以从何处得到帮助呢？\par

\Emph{老人：}在很久以前，有一些人，比所有被称为哲学家的人都更古老，他们是义人，蒙天主所爱，他们藉着圣神说话，预言将要发生并且现在正在发生的事件。他们被称为先知。只有这些人看到了真理并向人宣告了真理，不畏惧也不惧怕任何人，不受荣誉欲的影响，只讲他们所看见和听到的事，充满圣神。他们的著作仍然存在，阅读它们的人，如果相信它们，在认识事物的始末以及哲学家应当知道的那些事情上，会得到很多帮助。因为他们并不在论文中使用论证，因为他们是不言自明的真理的见证，是值得相信的；已经发生和正在发生的事件迫使你同意他们所说的话，尽管他们因所行的奇迹本已值得信任，因为他们既荣耀了万物的创造者——天主和父，也宣布了祂所派遣的儿子基督。而假先知们，充满谎言之不洁神灵，不仅没有这样做，而且胆敢行一些奇事以迷惑人，并荣耀谬误的神灵和恶魔。但你要祈求，最重要的是，光明的门为你打开；因为这些事情不是所有人都能看见或理解的，只有天主和祂的基督赐予智慧的人才能理解。\par

% structure: p0095 source=h2 target=subsection reason=relative-heading-rank; base=h2

\subsection{第八章 犹斯定因交谈而被激发对基督的爱}

\Person{犹斯定}: \Quote{他说了这些和许多其他目前没时间提及的事情后，就离开了，嘱咐我注意这些事；此后我再没有见过他。但立刻有一团火焰在我灵魂中燃起；对先知和对那些基督之友的爱占据了我；当我反复思考他的话时，我发现唯有这一哲学是安全且有益的。因此，为了这个缘故，我成了一个哲学家。此外，我希望所有像我这样下定决心的人，不要远离救主的话语。因为这些话本身具有可怕的力量，足以使那些偏离正道的人心生敬畏；同时，对于那些勤勉实践它们的人，则赐予最甜蜜的安息。因此，如果你对自己有所关心，如果你热切寻求救恩，如果你相信天主，那么——既然你并非对此漠不关心——你就能认识天主的基督，并在入门之后，过上幸福的生活。}\par

我说完这话，我亲爱的朋友们——那些与\Person{特黎丰}在一起的人——笑了起来；但\Person{特黎丰}只是微笑着说：\par

\Person{特黎丰}: \Quote{我赞成你的其他言论，并钦佩你研究神圣事物的热忱；但对你来说，更好的做法是仍然坚守柏拉图或其他人的哲学，培养忍耐、自制和适度，而不是被虚假言辞所欺骗，追随那些无名之辈的意见。因为如果你停留在那种哲学中，并且无可指责地生活，更好的命运仍会留给你；但当你离弃了天主，信赖了人，还有什么安全等待你呢？因此，如果你愿意听我的（因为我已经把你当作朋友），首先受割礼，然后遵守关于安息日、节日和天主月朔所制定的各项规条；总之，遵守律法上所写的一切：这样你或许会从天主那里获得怜悯。但基督——如果他确实已经出生，且存在于某处——是无人知晓的，他甚至不认识自己，并且毫无权能，直到\Person{厄里亚}来为他傅油，并使他显明给众人。而你们，接受了无根据的传闻，为自己造出了一个基督，并且为了他而轻率地丧亡。}\par

% structure: p0099 source=h2 target=subsection reason=relative-heading-rank; base=h2

\subsection{第九章 基督徒没有相信无根据的故事}

\Person{犹斯定}: \Quote{我原谅并宽恕你，我的朋友，因为你不晓得你说的是什么，而是被那些不明白圣经的教师说服了；你像占卜者一样，信口开河。但如果你愿意听一听关于他的讲述，即我们如何没有被欺骗，并且不会停止承认他——尽管人的辱骂堆积在我们身上，尽管最可怕的暴君强迫我们否认他——我就站在这里向你证明，我们并没有相信空洞的神话，或毫无根据的言辞，而是相信充满天主圣神、满有权能、洋溢着恩宠的话语。}\par

于是那些与他在一起的人又笑起来，并以不雅的方式喧嚷。然后我站起来，正要离开；但他抓住我的衣服，说我不该那样做，直到我履行了我所承诺的。\par

\Person{犹斯定}: \Quote{那么，不要让您的同伴这样喧闹，或如此不体面。但如果他们愿意，就让他们安静地听；或者，若有更好的事情妨碍他们，就让他们离开；而你我，退到某个地方，在那里休息，就可以结束讨论。}\par

\Person{特黎丰}认为我们这样做很好；于是，一致同意后，我们退到了\Place{克西司托斯}的中央空间。他的两个朋友，在嘲笑和戏弄了我们的热忱之后，离开了。我们到达那个两边都有石凳的地方后，那些与\Person{特黎丰}在一起的人在一侧坐下，彼此交谈，其中一人插话谈到了在\Place{犹太地}进行的战争。\par

\SourceRecord{Translated by Marcus Dods and George Reith.；Ante-Nicene Fathers；Vol. 1.；Edited by Alexander Roberts, James Donaldson, and A. Cleveland Coxe.；Buffalo, NY: Christian Literature Publishing Co.,；1885.；<http://www.newadvent.org/fathers/01281.htm>.}

% structure: p0001 source=h1 target=section reason=page-title

\section{与特黎丰对话录（第10至30章）}

% structure: p0002 source=h2 target=subsection reason=relative-heading-rank; base=h2

\subsection{第十章 特黎丰唯独为此责备基督徒——不遵守法律}

他们停止后，我又向他们发言。\par

\Person{尤斯定}：朋友们，除了我们不按照法律生活、不像你们祖先那样在肉身受割损、不像你们那样守安息日之外，你们还责备我们其他什么事吗？我们的生活和习俗也被你们诽谤吗？我问这个问题：你们是否也相信关于我们的谣言，说我们吃人，并且在宴席后熄灯后进行淫乱同居？或者你们仅仅因为我们在这些教义上持守，相信一种你们认为不真实的见解而谴责我们？\par

\Person{特黎丰}：这正是我们感到惊讶的，但群众所谈论的那些事并不值得相信，因为它们与人性极其相悖。此外，我知道你们在所谓的福音中的规条是如此奇妙而伟大，以至于我怀疑没有人能遵守它们，因为我仔细读过。但让我们最困惑的是：你们自称虔诚，自以为比别人更好，却丝毫没有与他们分别，也没有改变你们的生活模式与外邦人不同，你们不守任何节日或安息日，也没有割损的礼仪；而且，你们将希望寄托在一个被钉十字架的人身上，却指望从天主那里得到好处，同时又不遵守祂的诫命。你们没有读过吗？凡第八日不受割损的人，应从他的百姓中被剪除。这同样适用于外方人和奴隶。而你们，轻率地蔑视这盟约，拒绝由此而来的义务，并试图说服自己你们认识天主，然而你们并未实行任何敬畏天主之人的行为。因此，如果你们能在这些点上为自己辩护，并显明你们如何希望在丝毫不遵守法律的情况下获得任何好处，我们将非常乐意听你们说，并且我们还会进行其他类似的探究。\par

% structure: p0006 source=h2 target=subsection reason=relative-heading-rank; base=h2

\subsection{第十一章 法律被废除；天主应许并赐予新约}

\Person{尤斯定}：特黎丰，除了那位创造并安排了这整个宇宙的天主之外，没有别的神，从永远也不曾有过。我们不认为我们有一位神，你们有另一位，而是那位以大能的手和伸展的臂膀引领你们祖先出离埃及的天主是唯一的天主。我们也不信任其他任何神（因为没有别的神），而是信任你们也信赖的那位，即亚巴郎、依撒格和雅各伯的天主。但我们并非通过梅瑟或通过法律来信赖；因为那样的话，我们就会与你们一样。但现在——（因为我曾读到，将有一部终极的法律和盟约，是万有之中最首要的，凡是寻求天主产业的人，如今都必须遵守它。因为在曷勒布颁布的法律如今已陈旧，只属于你们自己；而\Emph{这部}却是为万民普遍设立的。如今，法律对立于法律，废除了之前的法律；后来的盟约同样终止了先前的盟约；而一部永恒而终极的法律——即基督——已赐予了我们，这盟约是可信的，此后将没有法律、没有诫命、没有条例。你们没有读过依撒意亚所说的这段话吗：\InnerQuote{你们要听从我，听从我，我的百姓；列王啊，要侧耳听我：因为法律要从我发出，我的公义要作万民的光。我的正义迅速临近，我的救恩将要发出，列国要依靠我的膀臂。}以及耶肋米亚论及这同一新盟约时这样说：\InnerQuote{看，日子将到，上主说，我要与以色列家和犹大家订立新约，不像我与他们祖先所立的约，在那日我牵着他们的手，领他们出离埃及地\BibleRef{耶 31:31}}。）因此，既然天主预告了将要订立的新盟约，且这盟约要做万民的光，我们看见并确信，人们正通过那位被钉十字架者——耶稣基督的名，离开他们的偶像和其他不义，来接近天主，并且持守他们的宣认直到死亡，保持虔诚。此外，通过工作及伴随的神迹，所有人都能理解祂就是新法律、新盟约，以及从万民中期待天主恩惠之人的盼望。因为真正的精神的以色列，以及犹大、雅各伯、依撒格和亚巴郎（亚巴郎在未受割损时因信德蒙天主赞许和祝福，被称为万民之父）的后裔，正是我们这些通过这位被钉十字架的基督被引领到天主面前的人，这将在我们继续讨论时得到证明。\par

% structure: p0008 source=h2 target=subsection reason=relative-heading-rank; base=h2

\subsection{第十二章 犹太人违反永恒的法律，并曲解梅瑟的法律}

我也引用了另一段依撒意亚的话，他呼喊说：\Quote{请听我的话，你的灵魂必要生活；我要与你订立永久的盟约，即达味可靠的恩慈。看，我已立祂为万民的证人：不认识你的民族要呼求你，不认识你的百姓要投奔你，因为你的天主、以色列的圣者已荣耀了你。}\par

\Person{Justin}：你轻视了这同一法律，藐视了祂的新而圣洁的盟约；如今你既不接受它，也不悔改你的恶行。因为\Quote{你的耳朵塞住了，你的眼睛瞎了，你的心变硬了}，\Person{耶肋米亚}曾呼喊；然而即便如此你仍不听从。立法者就在眼前，你却不看见祂；福音传给了穷人，盲人看见了，你却不明白。你现在需要第二次割损，虽然你大大夸耀肉身的割损。新法律要求你遵守永远的安息日，而你因为闲散一天，就自以为虔诚，却不分辨为何这样命令你；若你吃无酵饼，你就说天主的旨意已完成了。我们的主天主并不喜悦这样的遵守：若你们中间有发假誓或偷窃的，让他停止这样做；若有奸淫的，让他悔改；这样他就遵守了天主甜蜜而真正的安息日。若有人手不洁净，让他洗涤而成为洁净。\par

% structure: p0011 source=h2 target=subsection reason=relative-heading-rank; base=h2

\subsection{依撒意亚教导：罪过因基督的血而获得赦免}

\Person{Justin}：因为依撒意亚并没有差你到浴水里，在那里洗去谋杀和其他罪过（这连全部海水也不足以洗净）；而是，正如所料，那是古时拯救的浴水，跟随那些悔改的人，他们不再被山羊和绵羊的血，或小母牛的灰，或细面的献仪所洁净，而是藉着因基督的血而有的信德，并藉着祂的死——祂正是为此而死——正如依撒意亚他自己所说的，当他这样说时：\par

\BibleQuote{主将在万民眼前显露祂的圣臂，万民和地极都要看见天主的救恩。离开，离开，离开，从那里出去，不要触摸不洁之物；从她中间出来，携带主器皿的人哪，要洁净自己，因为你们不是仓促而行。因为主必在你们前面行；主，以色列的天主，必聚集你们。看哪，我的仆人必行事明智；祂必被高举，大大受荣耀。许多人因你惊奇，同样，你的容貌和你的荣耀必被人毁损，过于常人。照样，许多国民要因祂惊奇，君王要闭口无言；因为未曾传与他们的，他们必看见；未曾听闻的，他们必明白。主啊，谁相信了我们的报告？主的臂膀向谁显露了？我们在祂面前宣布祂如一个孩童，如干旱地中的根。祂无威仪，也无美貌；我们看见祂时，祂无美貌，也无佳形；但祂的容貌被污损，比世人更甚。祂是个受苦的人，熟悉病痛，因为祂的脸被转开；祂被藐视，我们也不尊重祂。祂担当我们的罪，为我们忧愁；我们却以为祂受责罚，被天主击打苦待了。但祂为我们的过犯受害，为我们的罪孽压伤；因祂受的刑罚，我们得平安；因祂受的鞭伤，我们得医治。我们都如羊走迷，各人偏行己路；主把我们众人的罪孽都归在祂身上。祂因受欺压，不开口。祂像羊羔被牵到宰杀之地，又像羊在剪毛的人面前无声，祂也是这样不开口。祂受辱时，审判被夺去。谁能述说祂的世代？因为祂的生命从地上被夺去。因我百姓的过犯，祂被交付至死。我使恶人作祂的坟墓，使富人作祂的死地，因为祂没有行过强暴，口中也没有诡诈。主却愿将祂从痛苦中洁净。若祂被交为赎罪祭，你的灵魂必看见长寿的后裔。主愿将祂的灵魂从忧苦中带走，使祂见光，并以智慧塑造祂，使那公义者——祂为多人服务——成义。祂必担当我们的罪；因此祂必承受许多，并分得强者的掠物，因为祂的灵魂被交付至死；祂被列在罪犯之中，祂担当多人的罪，又为他们的过犯被交。你这不生育、未生产的，歌唱吧；你这未曾经历产痛的，发声欢呼吧；因为没有丈夫的比有丈夫的儿女更多。因为主说：扩张你帐幕之地，展开你居所的幔子，不要限制；要放长你的绳子，坚固你的橛子；向左向右开展；你的后裔必得外邦人为产业，并使荒凉的城邑有人居住。不要惧怕，因为你必不蒙羞；不要抱愧，因为你必不受辱；你必忘记永远的羞耻，不再记念你寡居的羞辱，因为主为自己立了名，那救赎你的必在全地被称为以色列的天主。}\par

% structure: p0014 source=h2 target=subsection reason=relative-heading-rank; base=h2

\subsection{正义不在于犹太仪式，而在于藉基督在圣洗中赐予的内心悔改}

\Emph{儒斯定：}因此，藉这为天主子民的过犯而设立的悔改与认识天主的洗濯（正如依撒意亚所呼喊的），我们相信并见证，那他所宣布的圣洗本身足以洗净悔改的人；这就是生命之水。但你们为自己所挖的水池是破裂的，对你们毫无益处。那仅洗净肉体和身体的圣洗有何用呢？给灵魂行圣洗，从愤怒、贪婪、嫉妒和仇恨中洗净；看！身体就洁净了。因为无酵饼的象征意义就在于此：你们不要做出邪恶之酵的旧行径。但你们以肉性的意义理解一切，以为做这些事就是虔诚，而你们的灵魂却充满欺骗，总之，充满各种邪恶。因此，在吃无酵饼的七天之后，天主命令他们混合新酵，即履行其他善功，而不是效法旧的和邪恶的行为。因为这位新立法者要求你们这样，我将再次引用我已引述的话，以及其他一些被省略的话。这些是依撒意亚以下述方式叙述的：\par

\Quote{你们要听从我，你们的灵魂就必存活；我要与你们立定永久的盟约，即达味的可靠恩慈。看，我已立他为万民的见证，为列国的领袖和司令。不认识你的民族要呼求你；不认识你的百姓要逃往你那里，因你的天主，以色列的圣者，因为他荣耀了你。你们要寻求天主；当你们找到他时，要呼求他，趁他临近的时候。让恶人离弃自己的道路，不义的人离弃自己的思想；让他回归主，他必得怜悯，因为他必会大量赦免你们的罪恶。因为我的思想不是你们的思想，我的道路也不是你们的道路；正如天离地那样远，我的道路离你们的道路，你们的思想离我的思想也同样遥远。因为正如雪和雨从天降下，不再返回，直到浇灌大地，使地发芽结实，给播种者种子，给食者粮食；同样，从我口中发出的话也必如此：它必不返回，直到完成了我所喜悦的一切，并使我的命令亨通。因为你们必欢欢喜喜而出，平平安安被引导。大山小山必在你们面前跳跃，田野的一切树木也都鼓掌；松树要长出代替荆棘，番石榴要长出代替蒺藜。主要为留名，为永久的记号，永不废止！}\par

这些和类似的先知们所写的话，哦，特里丰，有些是指基督的第一次来临，其中他被传讲为无荣耀、隐晦、有死的外貌；但另一些是指他的第二次来临，那时他要从荣耀中并在云上显现；你们的民族要看见并认识他们所刺透的那一位，正如十二先知之一的欧瑟亚和达尼尔所预言的。\par

% structure: p0018 source=h2 target=subsection reason=relative-heading-rank; base=h2

\subsection{第十五章 真正的斋戒是什么}

\Emph{儒斯定：}所以，要学习持守天主的真正斋戒，正如依撒意亚所说的，使你们能悦乐天主。依撒意亚曾这样呼喊：\par

\Quote{你要大声呼喊，不要吝惜：扬起你的声音像吹号，向我的百姓宣告他们的过犯，向雅各伯家宣告他们的罪恶。他们天天寻求我，乐意明白我的道，好像行义的国民，不离开天主的典章。他们向我求问公义的判断，喜悦亲近天主，说：我们为什么斋戒，你看不见？我们刻苦己心，你却不理会？看哪，你们斋戒的日子仍找自己的喜好，欺压一切属下的人。看哪，你们为争竞和辩论而斋戒，用拳头击打卑微人。你们今日这样斋戒，岂能使你们的声音听闻于上呢？我所拣选的斋戒，不是这样吧？是人刻苦己心的日子吗？难道你像环子般弯下头颈，披上麻布和灰烬，你就可以称这是斋戒，是主所悦纳的日子吗？这不是我所拣选的斋戒，主说：要解开一切不义的绳索，废除强暴的盟约，使受压迫的得到自由，避免一切不义的契约。把你的饼分给饥饿的人，将无家可归的穷人领到你屋下；看见赤身的，给他衣服穿；不要躲避你的骨肉。那时，你的光明要像晨光一样破晓，你的衣服要迅速升起；你的公义要在你前面行，天主的荣耀要环绕你。那时你呼求，主必应允；你说话，他必说：我在这里。如果你从你中间除掉轭，伸手指的，和说怨言的话；并诚心把你的饼分给饥饿的人，满足困苦的灵魂；那么，你的光明必要在黑暗中升起，你的幽暗要如日中天；你的天主必时常与你同在，你必因你灵魂的渴望得以满足，你的骨头将变得肥胖，像浇灌的园子，像水泉，像水流不竭之地。\BibleRef{依 58:1}}\par

\Quote{所以，你们要割损你们心里的包皮，}正如在这些段落中天主的话所要求的。\par

% structure: p0022 source=h2 target=subsection reason=relative-heading-rank; base=h2

\subsection{第十六章 割损作为记号，使犹太人因其对基督和基督徒的恶行而被驱逐}

\Person{儒斯定}: 天主曾亲自藉梅瑟发言，这样说：\Quote{你们要割去心里的硬皮，不要再硬著颈项。因为上主你们的天主是万主之主，伟大的、大能的、可畏的天主，不看情面，也不受贿赂。}又在肋未纪中说：\BibleQuote{因为他们得罪了我，藐视了我，又因他们行为与我反对，我也与他们反对，将他们剪除在他们仇敌之地。那时他们未受割礼的心必归顺。\BibleRef{肋 26:40}} 因为按肉体受割礼是从亚巴郎而来，是作为一个记号，使你们与其他民族分离，也与我们分离；并且你们独自承受现在正当地受的苦；你们的土地荒凉，你们的城邑被焚烧；外人在你们眼前吃你们的果实，你们中没有一人能上耶路撒冷。你们在其他人中除肉体的割礼外没有别的记号被认出。我想你们中间没有人敢说天主从未预见或将来的事，或没有预定各人应得的报应。因此，这些事公平而正义地临到你们，因为你们杀了那义者和祂之前的先知；现在你们又拒绝那些仰望祂以及那派遣祂者——全能的造物主天主——的人，并在你们的会堂里咒骂那些相信基督的人。因为你们没有能力加害我们，由于现在掌权者的缘故。但每当你们能够时，你们就这样做了。因此天主藉依撒意亚向你们呼吁说：\BibleQuote{看，义人怎样灭亡，无人放在心上。因为义人被从邪恶面前取去。他的坟茔必得平安，他从中间被取去。你们这些无法无纪的儿女，奸夫的后裔，妓女的子孙，前来这里。你们向谁戏弄？向谁张口？向谁伸舌头？\BibleRef{依 57:1}}\par

% structure: p0024 source=h2 target=subsection reason=relative-heading-rank; base=h2

\subsection{第十七章。犹太人派遣人到全地散布对基督徒的诽谤}

\Person{儒斯定}: 因为其他民族并没有像你们这样加害我们和基督；你们实在是那恶毒偏见反对义者和我们这些持守祂的人的始作俑者。当你们把祂——那唯一无玷而正义的人——钉在十字架上以后，因祂的创伤，那些藉祂亲近圣父的人得到治愈；你们明知祂已从死者中复活并升天，正如先知们所预言的，你们非但不为自己所作的恶行悔改，反而在那个时候从耶路撒冷挑选并派遣人到全地，去宣扬基督徒的无神异端已经兴起，并散布那些不认识我们的人所说的一切反对我们的话。因此，你们不仅成为自己不义的缘由，实际上也成为所有其他人不义的缘由。依撒意亚公正地呼喊：\BibleQuote{因了你们，我的名在外邦人中受亵渎。\BibleRef{依 52:5}} 又说：\BibleQuote{祸哉，他们的灵魂！因为他们对自己图谋了恶计，说：我们捆绑义人，因为他使我们厌烦。因此他们必吃自己行为的果实。祸哉，恶人！他必按他手所作的受恶报。} 此外，又说：\BibleQuote{祸哉，那些以长绳牵引自己的罪孽，以牛轭的绳索牵引自己的过犯的人！他们说：愿他的速度临近，愿以色列圣者的计划来临，叫我们知道。祸哉，那些称恶为善、称善为恶，以暗为光、以光为暗，以苦为甜、以甜为苦的人！} 因此，你们十分热心地在全地散布关于那由天主派遣的唯一无玷而正义之光的苦毒、黑暗和不义之事。\par

\Person{儒斯定}: 因为当祂在你们中间呼喊时，祂使你们厌烦：\BibleQuote{经上记载：我的殿宇是祈祷的殿宇，但你们竟把它弄成了贼窝！\BibleRef{玛 21:13}} 祂也在圣殿中推翻了兑换银钱者的桌子，并且大声说：\BibleQuote{祸哉，你们经师和法利塞人，假善人！因为你们捐献薄荷、芸香，却不遵守对天主的爱和正义。你们是粉刷的坟墓！外面看来好看，里面却装满死人的骨头。} 又对经师们说：\BibleQuote{祸哉，你们经师！因为你们拿着钥匙，自己不进，也阻止那些要进的人；你们是盲目的向导！}\par

% structure: p0027 source=h2 target=subsection reason=relative-heading-rank; base=h2

\subsection{第十八章。基督徒会遵守法律，若他们不知道为何设立法律}

\Person{儒斯定}: \Person{特黎风}啊，既然你亲自承认读过我们救主所教导的道理，我想我在先知话语之外加上祂的一些简短言语并没有做愚昧的事。所以，你们当洗濯，现今成为洁净，从你们灵魂中除去不义，正如天主命令你们在这水盆中洗濯，以真割礼受割。因为我们也愿意遵守肉体的割礼、安息日，以及总而言之所有庆节，如果我们不知道它们为何被命令给你们——即因了你们的过犯和你们的心硬。因为如果我们忍受恶人和魔鬼对我们所设计的一切，甚至在不言而喻的残酷、死亡和痛苦中，我们为那些施加如此之事的人祈求怜悯，并不愿对任何人作最小的反驳，正如新立法者所命令我们的：\Person{特黎风}啊，我们怎么会不遵守那些不伤害我们的仪式呢？——我说的是肉体的割礼、安息日和庆节？\par

% structure: p0029 source=h2 target=subsection reason=relative-heading-rank; base=h2

\subsection{第十九章。割礼在亚巴郎之前无人知晓。法律因他们心硬而由梅瑟赐下}

\Person{儒斯定}: 这正是我们所困惑的，也有理由，因为你们既然忍受这样的事，却不遵守我们现在所讨论的其他一切习俗。\par

然而，这割损并非为所有人必要，而只为你们，好使你们，如我早已说过，承受你们现今所正当地承受的这些事。我们也不接受那无用的蓄水池的圣洗，因为它与这生命的圣洗毫无关系。因此天主也曾宣布，你们离弃了祂这活水的泉源，为自己凿出了破裂而不能蓄水的池子。就连你们这些按肉体受割损的，也需要我们的割损；但我们有了后者，就不需要前者。因为倘若这必要，如你们所想，天主就不会使亚当不受割损；不会在亚伯尔未受割损而献祭时眷顾他的供物；也不会因未受割损的哈诺客而喜悦，他不见了，因为天主将他接去了。未受割损的罗特，因天使们和主亲自救他出来，而得以从索多玛得救。诺厄是我们种族的始祖，却未受割损，与他的儿女一同进了方舟。至高者的司铎默基瑟德未受割损，首先按肉体领受割损的亚巴郎向他缴纳了什一之物，而他却祝福了亚巴郎；天主曾藉达味的口宣告，要按照他的品级建立永远的司铎。因此，这割损只对你们是必要的，为叫这人民不再成为人民，这民族不再成为民族；正如十二先知之一的欧瑟亚所宣布的。此外，以上提到的所有义人，虽然不守安息日，却中悦天主；在他们之后，亚巴郎和他的一切后裔直到梅瑟，在梅瑟之下，你们民族显出对天主不义和忘恩，在旷野里铸造了牛犊。因此天主屈就于那民族，也命令他们献祭，像是献于祂的名，为叫你们不事奉偶像。然而你们并没有遵守这条诫命；反而把你们的儿女祭献给魔鬼。并且你们被命令守安息日，为叫你们保留对天主的纪念。因为祂的话这样宣布说：\BibleQuote{为叫你们知道我是救赎你们的天主。}\BibleRef{则 20:12}\par

% structure: p0032 source=h2 target=subsection reason=relative-heading-rank; base=h2

\subsection{为何规定肉食的选取}

\Person{尤斯定}：此外，你们被命令禁戒某些种类的食物，好使你们在吃喝时把天主摆在眼前，因为你们惯于且极易偏离对祂的认识，正如梅瑟也证实的：\BibleQuote{百姓坐下吃喝，起来玩耍。}\BibleRef{出 32:6} 又说：\BibleQuote{雅各伯吃了，得了饱足，就肥胖起来；那蒙爱的却踢跳：他肥胖了，粗壮了，就离弃了造他的天主。}\BibleRef{申 32:15} 因为梅瑟在\WorkTitle{创世纪}中曾告诉你们，天主准许义人诺厄吃各样动物，但不可吃带血的肉，那是\Emph{死的}。\par

当他正要说话时，说\Quote{如同青菜一样，}我先打断了他：\par

你们为何不接受\Quote{如同青菜一样}这句话，按天主所赐予的意义呢？即是说，如同天主赐予青菜供人维持生计，同样祂也赐予动物作为肉食。但你们说，此后对诺厄有一个区别，因为我们不吃某些青菜。照你们这样解释，事情就不可信了。首先我不必深究这个，虽然我能说并能主张每种蔬菜都是食物，都可吃。但尽管我们区别青菜，不吃所有种类，我们禁戒某些，不是因为它们是普通或不洁的，而是因为它们是苦的、有毒的或多刺的。但我们取用并吃所有甘甜、滋养、有益的青菜，无论是海产还是陆产的。同样，天主藉梅瑟的口命令你们禁戒不洁、不当和凶猛的动物；此外，你们在旷野吃玛纳，看见天主为你们所行的一切奇事，却制造并敬拜了金牛犊。因此祂不断地、公正地呼喊：\BibleQuote{他们是愚昧的儿女，毫无信德。}\BibleRef{申 32:6}\par

% structure: p0036 source=h2 target=subsection reason=relative-heading-rank; base=h2

\subsection{安息日是为人民的罪而设立，并非为义德之工}

\Person{尤斯定}：此外，天主命令你们守安息日，并施加其他诫命作为记号，如我早已说过的，是因你们和你们祖先的不义——祂宣告，为了外邦人的缘故，免得祂的名在他们中间被亵渎，祂容许你们中的一些人存留——这些话能向你们证明：它们由厄则克耳如此叙述：\par

\BibleQuote{我是上主你们的天主；你们要遵行我的律例，谨守我的典章，不可随从埃及的习俗；要守我的安息日为圣；它们要在我与你们之间作为记号，叫你们知道我是上主你们的天主。然而你们背叛了我，你们的儿女不遵行我的律例，也不谨守我的典章去遵行；人若遵行，就必因此活着。但他们亵渎了我的安息日。我本想在旷野向他们倾泻我的愤怒，在他们身上发尽我的怒气，我却未这样做，免得我的名在外邦人眼前完全被亵渎。我在他们眼前领他们出来，在旷野向他们举起我的手，要把他们分散在外邦人中，分散在各国；因为他们没有遵行我的典章，藐视我的律例，亵渎我的安息日，他们的眼目追随他们祖先的计谋。因此我也赐给他们不好的律例，以及不能使他们存活的典章。我要在他们自己的礼物中玷污他们，为要毁灭一切开胎的，当我经过他们时。}\BibleRef{则 20:19}\par

% structure: p0039 source=h2 target=subsection reason=relative-heading-rank; base=h2

\subsection{祭献和供物也是如此}

\Person{Justin}：但为使你们知道，正是因为你们民族的罪恶和他们的偶像崇拜，而非因为那些祭献有何必要，才同样规定了这些祭献，请听祂通过十二先知之一的亚毛斯所说的话：\BibleQuote{祸哉，你们盼望上主日子的人！上主的日子为你们有什么好处？那是黑暗，而非光明，如同人逃避狮子，却遇见了熊；他进入房屋，手扶墙壁，却被蛇咬。上主的日子岂非黑暗而非光明，极其黑暗，毫无光明吗？我恼恨、厌恶你们的庆节，你们盛大的集会我也不闻；所以你们虽然向我献全燔祭和素祭，我也不悦纳，也不垂顾你们肥畜的和平祭。从我面前除去你们的歌声和琴音，我不听你们的乐器。只愿公正如水滚流，正义如不竭的溪流。以色列家啊，你们在旷野中岂向我奉献过牺牲和祭品？你们抬着摩洛的帐幕和你们为自己所造的神像——勒番星的雕像。我必要将你们放逐到大马士革之外，全能的天主说。祸哉，那些在熙雍安逸，在撒玛黎雅山自恃的人！你们这些列国之首所挑选的人夺去了万民的初果；以色列家为自己进入。你们往加耳乃去观看，从那里到大哈特，再下到迦特，看看这些王国中最卓越的，是否他们的疆界比你们的疆界更大。你们这些向往邪恶日子、临近并持守虚假安息日的人；你们躺在象牙床上，舒身在榻上，吃羊群中的羔羊和牛群中的肥犊；你们在乐器声中欢呼，以为它们稳固而非短暂；你们用碗喝酒，用最好的膏油抹身，却不为若瑟的苦难忧伤。所以如今他们将被掳去，为首的被掳之人中的首先者；作恶者的房屋将被拆毁，厄弗辣因的马嘶声也将被除去。}又通过耶肋米亚说：\BibleQuote{你们收集肉块和祭品来吃吧！因为我领你们祖先出埃及的那一天，并没有吩咐他们关于祭品和奠祭的事。}又通过达味在第四十九篇圣咏中说：\BibleQuote{万神之神上主发言，从日出之地到日落之处呼唤大地。从熙雍，美丽绝伦的光辉显现。我们的天主要来临，绝不缄默；烈火在祂面前燃烧，风暴在祂四周旋转。祂呼唤上天和大地，为自己人民审判。你们召集我的圣民来，就是那些用祭品与我立约的人。诸天要传扬祂的公义，因为天主是审判者。请听，我的子民，我要发言；以色列啊，我要向你作证：我是天主，是你的天主。我不因你的祭品责备你，你的全燔祭常在我面前。我不从你的家中取公牛，也不从你的圈中取公山羊；因为林中的百兽是我的，山上的牲畜是我的。我知道天空的飞鸟，田野的美物也是我的。我若是饥饿，我不需要告诉你，因为世界和其中一切都是我的。我岂吃公牛的肉，喝公山羊的血呢？你要以赞美为祭献于天主，向至高者还你的愿，并在患难之日呼求我，我必拯救你，你也要光荣我。}但天主对恶人说：\BibleQuote{你凭什么宣读我的律例，将我的盟约放在你的口中？你恼恨管教，将我的话抛在脑后。你见了盗贼，便与他同伙；与行奸淫的人有分。你的口出恶言，你的舌编造诡计。你坐着毁谤你的兄弟，诬蔑你母亲的儿子。你做了这些，我默然不语；你竟以为我像你一样。但我要责备你，将你的罪行摆在你眼前。你们这些忘记天主的人，要思想这事，免得我将你们撕碎，无人搭救。赞美祭光荣我，我要将我的救恩显示给他。}因此，祂既不从你们那里取祭物，起初也没有命令献祭，因为它们对祂并非必要，而是因为你们的罪过。诚然，那称为耶路撒冷圣殿的，祂承认是祂的房屋或庭院，并非祂需要它，而是为使你们因这观念而奉献自己给祂，不再崇拜偶像。依撒意亚说：\BibleQuote{你们为我建造了什么房屋？上主说：天是我的宝座，地是我的脚凳。}\BibleRef{依 66:1}\par

% structure: p0041 source=h2 target=subsection reason=relative-heading-rank; base=h2

\subsection{第二十三章 犹太人对法律的看法对天主造成伤害}

\Person{Justin}：但若我们不承认这一点，我们就会陷于愚昧的意见，仿佛那在哈诺客时代以及所有其他未按肉体割损、未守安息日及其它礼规的人（因为梅瑟规定了这些）的时代是同一个天主；或者天主不愿人类各族持续行同样的正义行为：承认这一点似乎荒谬可笑。因此我们必须承认，永远不变的那一位，因着有罪的人而命令了这些及其类似的制度，我们必须宣称祂是仁慈、预知、一无所缺、公义且善良的。但如果并非如此，先生，请告诉我，你对我们正在探讨的这些事有什么想法？\par

无人回应：\par

为此，\Person{特黎风}，我要向你和那些愿意成为皈依者的人，宣讲我从那人那里听到的神圣信息。你们难道不见元素并非怠惰，也不守安息日么？保持你们生来的样子吧。因为若在亚巴郎之前不需要割损，或在梅瑟之前不需要守安息日、庆节和祭献；那么现在，按照天主的旨意，天主子耶稣基督由亚巴郎的后裔中的童贞女无罪地诞生之后，就更不需要它们了。因为当亚巴郎本人尚在未受割损时，他已因对天主的信德而成义并蒙受祝福，正如圣经所载。此外，圣经和事实本身迫使我们必须承认，他接受割损是作为标记，而非为了成义。因此，关于这民族恰当地记载道：凡在第八天未受割损的灵魂，应从其家族中剪除。而且，女性无法接受肉体上的割损，这就证明此割损是作为标记赐下的，而非作为成义的行为。因为天主同样赐给女性能力，遵守一切正义和德行之事；但我们看到，男性的身体形态被造得与女性的不同；然而我们知道，没有人单因此故而为义或不义，而是由虔敬和正义来断定。\par

% structure: p0045 source=h2 target=subsection reason=relative-heading-rank; base=h2

\subsection{第二十四章　基督徒的割损远为卓越}

\Person{儒斯定}：诸君，我们现在可以说明第八天所蕴含的某种奥秘意义，是第七天所没有的，且是天主通过这些礼仪所昭示的。但为免我此刻似乎离题，请听我的话：那割损的血已经过时，我们信赖救恩的血；如今有另一盟约，另一法律已从熙雍发出。耶稣基督——如以上所宣告的——用石刀为所有愿意的人施行割损，\BibleRef{苏 5:2} 使他们成为正义的国民，持守信德的民族，持守真理，维护和平。那么，凡敬畏天主、愿见耶路撒冷福祉的人，请跟我来。来，让我们走向上主的光明；因为祂已解放了祂的子民，雅各伯的家。来，万国；让我们在耶路撒冷聚集，不再因她子民的罪恶受战争困扰。\BibleQuote{因为我没有寻找我的，我让他们看见；我没有求问我的，他们找到我；}祂藉依撒意亚喊道：\BibleQuote{我说：看哪，我在这里，对那些未奉我名的民族。我整天向悖逆和抗辩的子民伸开双手，他们行走在不善的道路上，随从自己的罪恶。这是一个在我面前惹我发怒的民族。}\BibleRef{依 65:1}\par

% structure: p0047 source=h2 target=subsection reason=relative-heading-rank; base=h2

\subsection{第二十五章　犹太人徒然夸耀自己是亚巴郎的子孙}

\Person{儒斯定}：那些自称为义，并自称是亚巴郎子孙的，将甚愿与你们一同继承产业；正如圣神藉依撒意亚的口，以他们的口吻喊说：\par

\BibleQuote{从天上回来，从祢圣洁荣耀的居所垂视。祢的热忱和力量在哪里？祢丰厚的仁慈在哪里？因为祢扶持了我们，上主。因为祢是我们的父，亚巴郎不认识我们，以色列也不承认我们。但是祢，上主，我们的父，拯救我们；从起初，祢的名就在我们身上。上主，祢为什么使我们偏离祢的道路？使我们的心刚硬，以致不敬畏祢？为祢仆人的缘故，为祢产业的支派，求祢回来，使我们稍微继承祢的圣山。我们如同起初一样，当祢没有管辖我们，祢的名没有呼求在我们身上时。倘若祢开启诸天，群山必在祢面前颤栗；它们要熔化，如同蜡在火前熔化；火要吞灭敌人，祢的名要在敌人中显现；列国要在祢面前混乱。当祢行奇事时，群山必在祢面前颤栗。从起初我们未曾听闻，我们的眼也未见过除祢以外的神；而祢的作为，就是祢要向悔改者显示的仁慈。祂必迎接行义的人，他们要记念祢的道路。看哪，祢发怒，我们犯罪。因此我们迷失了，全都成了不洁的，我们所有的义德如同污秽的破布；我们因罪孽而如叶子枯干；因此风必把我们吹去。没有人呼求祢的名，或记得抓住祢；因为祢转脸不顾我们，因我们的罪恶而放弃了我们。现在求祢回来，上主，因为我们都是祢的子民。祢的圣城已变为荒凉。熙雍成了旷野，耶路撒冷成了咒诅；我们圣洁的殿宇，以及我们祖先所祝福的荣耀，已被火焚烧；所有荣耀的民族也与之同亡。除了这些灾难之外，上主，祢克制自己，缄默不语，大大地谦抑了我们。}\par

\Person{特黎风}：你说什么？我们中间没有一个人能在天主的圣山上继承产业吗？\par

% structure: p0051 source=h2 target=subsection reason=relative-heading-rank; base=h2

\subsection{第二十六章　犹太人除非藉着基督，否则无法得救}

\Person{儒斯定}：我并不这样说；但是那些迫害并且仍在迫害基督的人，如果他们不悔改，将不能在圣山上继承任何产业。而那些相信祂并悔改所犯罪过的外邦人，他们要同祖先、先知和出自雅各伯的义人一同继承产业，即使他们既不守安息日，也不受割损，也不守庆节。他们确实要承受天主圣洁的产业。因为天主藉依撒意亚这样说：\par

\BibleQuote{我，上主天主，因正义召叫了祢，并握着祢的手，坚强祢；我立了祢作为人民的盟约，外邦人的光明，为开启盲人的眼目，把被囚的从锁链中领出，把坐在黑暗中的从监牢中领出。}\BibleRef{依 42:6}\par

又说：\par

\BibleQuote{为万民竖立旗帜；因为，看，主已使声音传到地极。向熙雍女子说：看，你的救主来了；他带着他的赏报，他的工作在他面前；他要称它为圣洁的民族，为主所救赎。你将被称作被寻找的城，不被遗弃。那从厄东而来，身着红衣从波责辣来的是谁？那衣着华美，以大力上行的是谁？我讲论正义和救恩的审判。为何你的衣裳是红的，你的衣服像从酒榨中踩出来的？你满是被踩的葡萄。我独自踩了酒榨，万民中没有人同我在一起；我在烈怒中践踏他们，将他们压碎在地上，将他们的血洒在地上。因为复仇的日子临到他们，救赎之年已到。我观看，无人帮助；我诧异，无人扶持；我的手臂为我施行拯救，我的烈怒临到他们；我在烈怒中践踏他们，将他们的血洒在地上。}\par

% structure: p0056 source=h2 target=subsection reason=relative-heading-rank; base=h2

\subsection{第二十七章 为何天主借着先知教导与梅瑟相同的事}

\Person{特黎丰}: 为什么你从先知著作中选择和引用你想要的，却不提及那些明确命令守安息日的章节？因为依撒意亚这样说：\par

\BibleQuote{你若在安息日转离你的脚步，不在我的圣日做你喜欢的事，并称安息日为你天主的圣乐；你若不在工作中抬脚，不从你口中说出言语；那么你必信赖主，他必使你上升，享受土地的美物，并以你祖先雅各的产业喂养你：因为主亲口说了。\BibleRef{依 58:13}}\par

\Person{儒斯定}: 我略过了它们，朋友们，不是因为这样的预言与我相悖，而是因为你们已经明白并确实明白，虽然天主通过众先知命令你们做同样的事，他也通过梅瑟命令了，但因你们心硬和对他的忘恩，他不断宣告它们，以便即使这样，如果你们悔改，你们能悦乐他，既不献你们的儿女给魔鬼，也不与盗贼同伙，不爱礼物，不追求报复，不忽略为孤儿伸冤，不疏忽对寡妇的公义，也不使你们的手沾满血。\BibleQuote{因为熙雍的女子高傲地行走，用眼睛抛媚眼，拖着衣裙行走。\BibleRef{依 3:16}} 他宣告说：\BibleQuote{他们都偏离了正路，都成了无用的。没有一个明白的，连一个也没有。他们用舌头行欺骗，他们的喉咙是敞开的坟墓，他们嘴唇下有蛇的毒汁，毁灭与痛苦在他们的路上，和平的道路他们不曾知道。} 所以，如同起初，这些事因你们的邪恶而吩咐你们，同样，因你们固守邪恶，或者更确切地说，你们更倾向于邪恶，借着同样的诫命，他呼唤你们记起或认识它。但你们是心硬、无知的百姓，又瞎又瘸，没有信德的儿女，正如他自己所说，只用嘴唇尊崇他，心却远离他，教导你们自己的道理而不是他的。因为，告诉我，天主希望司铎们在安息日献祭时犯罪吗？或者那些受割礼并在安息日行割礼的人犯罪吗？因为他命令在第八天——即使恰逢安息日——出生的婴孩总要受割礼。难道不能在前一天或后一天给婴儿施行手术吗，如果他知道在安息日行这事是犯罪？或者他为什么不教导那些被称为义人、悦乐他、在梅瑟和亚巴郎之前生活、未受割礼、不守安息日的人遵守这些制度呢？\par

% structure: p0060 source=h2 target=subsection reason=relative-heading-rank; base=h2

\subsection{第二十八章 真正的义德由基督获得}

\Person{特黎丰}: 我们刚才听见你提出这个考虑，我们也留意了：因为说实话，这值得注意；而那个最令人满意的答案——即，他以为好就这样——并不使我满意。因为那些无法回答问题的人总是采取这种托词。\par

\Person{尤斯定}：既然我从圣经和事实本身带来证据和劝诫，你们切莫迟延或犹豫，相信我，尽管我是一个未受割损的人；留给你们成为归化者的时间不多了。如果基督的来临抢在你们前面，你们将徒然后悔，徒然哭泣；因为祂不会垂听你们。\Person{耶肋米亚}曾向百姓呼喊说：\BibleQuote{犁开你们的荒地，不要撒种在荆棘中。你们要为主行割损，割损你们心头的包皮。}\BibleRef{耶 4:3}因此，不要撒种在荆棘中，不要撒种在未耕之地，那里结不出果实。认识基督，看哪，那良好、肥沃的荒地就在你们心里。\Quote{因为，看，日子将至——上主说——我必惩罚一切肉身受割损的人：埃及、犹大、厄东、摩阿布的子孙。因为所有外邦人都是未受割损的，而以色列全家心中也未受割损。}你们看出天主不是说这作为标记的割损吗？因为它对埃及人、摩阿布的子孙、厄东的子孙都毫无用处。但即使一个人是斯提特人或波斯人，如果他认识天主和祂的基督，遵守永远的正义法令，他就行了那良好而有益的割损，是天主之友，天主因他的礼物和祭献而喜悦。但我将把天主的话摆在你们面前，朋友们，祂曾通过十二位先知之一的玛拉基亚对百姓说：\BibleQuote{我不喜悦你们——上主说——也不从你们手中接纳祭献。因为从日出到日落，我的名在外邦人中受尊崇；处处有人向我的名献祭，甚至纯洁的祭品；因为我的名在外邦人中受尊崇——上主说——你们却亵渎了它。}\BibleRef{拉 1:10}天主又通过达味说：\Quote{我不认识的百姓事奉了我；他们一听见就服从了我。}\par

% structure: p0063 source=h2 target=subsection reason=relative-heading-rank; base=h2

\subsection{第二十九章：基督对守法律者无用}

\Person{尤斯定}：让我们一起光荣天主，万民聚集，因为祂也眷顾了我们。让我们通过光荣之王、万军之主光荣祂。因为祂也恩待了外邦人；我们的祭献比你们的更蒙祂悦纳。那么，我这被天主作证的人，还需要割损吗？我这已受圣神洗礼的人，还需要那另一种洗礼吗？我想，我说这些，即使是智力有限的人也能说服。因为这些话不是由我预备的，也不是用人的技巧修饰的；而是达味歌唱的，依撒意亚宣讲的，匝加利亚宣告的，梅瑟写下的。特黎风，你认识这些吗？它们包含在你们的圣经里，更确切地说，包含在我们的圣经里。因为我们相信它们；但你们虽然读了，却捕捉不到其中的精神。不要因天主创造我们时赋予的身体上的未受割损而生气或指责我们；也不要认为我们在安息日喝热水奇怪，因为天主在这一天治理宇宙如同在其他日子一样；司祭在安息日也如其他日子一样奉命献祭；更有许多义人从未履行这些法律规定的礼仪，却仍被天主亲自作证。\par

% structure: p0065 source=h2 target=subsection reason=relative-heading-rank; base=h2

\subsection{第三十章：基督徒拥有真正的义德}

\Person{尤斯定}：但将此归咎于你们自己的邪恶，以致那些无知的人甚至可以指控天主不总以同样的正义法令教导所有人。因为许多没有领受恩宠的人认为这类制度不合理、与天主不配，因为他们不知道你们的民族在罪恶状态和属灵疾病中蒙召进行悔改和心灵的皈依；而且梅瑟死后宣告的预言是永远的。朋友们，这在圣咏中也有所提及。而我们因这些预言变得有智慧，承认上主的法令比蜂蜜和蜂房更甘甜，这从以下事实显而易见：即使面对死亡威胁，我们也不否认祂的名。此外，所有相信祂的人都看见，我们祈求祂保护我们脱离邪魔，即邪恶和欺骗的灵；正如预言的话以比喻的方式，模拟一个相信祂的人所宣告的。因为我们确实通过耶稣基督不断恳求天主，保守我们脱离那些敌对敬拜天主的魔鬼，我们从前事奉过他们，为的是在我们经由祂归向天主后，成为无可指摘的。因为我们称祂为助佑者和救赎者，祂的名号甚至使魔鬼恐惧；如今，当魔鬼在般雀比拉多（犹太总督）治下被钉十字架的耶稣基督的名号被驱魔时，它们就被制服了。因此，众人清楚可见，祂的父赐予祂如此大的权能，藉此魔鬼屈服于祂的名号和祂受苦的救世计划。\par

\SourceRecord{Translated by Marcus Dods and George Reith.；Ante-Nicene Fathers；Vol. 1.；Edited by Alexander Roberts, James Donaldson, and A. Cleveland Coxe.；Buffalo, NY: Christian Literature Publishing Co.,；1885.；<http://www.newadvent.org/fathers/01282.htm>.}

% structure: p0001 source=h1 target=section reason=page-title

\section{\WorkTitle{与特里弗对话录}（第31—47章）}

% structure: p0002 source=h2 target=subsection reason=relative-heading-rank; base=h2

\subsection{第31章 如果基督现今的权能如此伟大，那么第二次来临时的权能该何等伟大呢！}

\Emph{\Person{犹斯定}:} 但是，如果如此大的权能已经伴随并依旧伴随着祂受苦的救世工程，那么伴随祂光荣来临的权能又该何等的伟大呢！因为祂要像人子一样驾云降临，正如达尼尔所预言的，祂的天使们也要同祂一起来。这些话是：\BibleQuote{我观看，直到宝座设立了，有一位万古常存者坐着，祂的衣服洁白如雪，头发如纯净的羊毛。祂的宝座是火焰，祂的车轮是燃烧的火。一条火河从祂面前流出。事奉祂的有千千万万，侍立在祂面前的有万万。审判者已坐庭，案卷都展开了。我因那角所说狂傲的话而观看，直到那兽被击杀，身体被毁，扔进火里焚烧。其余野兽的权柄都被夺去，它们的寿命被限定直到一段时期。我在夜间的异象中观看，见有一位像人子的，驾着天云而来，来到万古常存者面前，被领到祂跟前。祂得了权柄、尊荣和国度，各民族、各邦国、各语言的人都事奉祂。祂的统治是永远的统治，永不会消逝；祂的王国不会被毁灭。我的精神在我体内发颤，我头脑中的异象使我困扰。我走近一位侍立者，问他这一切究竟是什么意思。他回答我，向我解释这事：这些巨大的野兽是四个王国，它们将从地上消灭，不再永久执掌权柄，直到永远。然后我渴望确切知道那第四只兽，它毁灭了所有其余的，极其可怕，它有铁牙铜爪，吞吃、压碎，用脚践踏残余；也关于它头上的十只角，以及那只新长出的角，它使前三只角掉落了。那角有眼睛和一张说大话的嘴；它的面貌超过其他。我观看那角与圣徒作战，战胜了他们，直到万古常存者来到；祂为至高者的圣徒伸冤。时候到了，至高者的圣徒得了王国。关于第四只兽，我被告知：在地上将有第四个王国，它要胜过所有这些王国，要吞吃全地，要毁灭它，使它彻底荒废。那十只角是十个王；他们之后将出现另一个王；他要在邪恶的事上超过先前的，他要制服三个王，他要说攻击至高者的话，要推翻至高者其余的圣徒，想要改变节期和法律。它将被交在他手中一段时间，二段时间和半段时间。审判已设立，他们要夺去他的权柄，要消灭它，直到永远。国度、权柄和天下诸国的大权都赐给了至高者的圣民，他们要在永远的国度中执政；所有的权柄都要服从祂，听命于祂。事情到此为止。我达尼尔极其惊惶，我的言语在我里面改变了；但我将这事存在心里。}\BibleRef{达 7:9}\par

% structure: p0004 source=h2 target=subsection reason=relative-heading-rank; base=h2

\subsection{第32章 特里弗反对说，达尼尔描述基督是光荣的；犹斯定区分两次来临}

\Emph{\Person{特里弗}:} 先生，这些以及类似的圣经教导迫使我们等待那位作为人子从万古常存者那里接受永恒王国的那一位。但你们所谓的这位基督却是可耻而无荣耀的，以致天主法律中的最后诅咒落到了祂身上，因为祂被钉死在十字架上。\par

\begin{quotation}
\Person{犹斯定}：各位先生，倘若圣经——即我已引用的那些——并未说祂的形貌是无光彩的，祂的世代未被宣告，祂的死亡使富人丧亡，因祂的创伤我们得痊愈，祂被牵去如同羊羔；倘若我未曾解释祂会有两次来临——一次是你们刺透祂，第二次是你们要认出你们所刺透的那位，各支派要哀悼，妇女单独，男子单独——那么我所说的话必定是含糊不明的。但现在，我借着你们尊为神圣和先知性的圣经内容，试图证明我所说的一切，希望你们中有人能被发现是属于那因万军的上主的恩宠而存留的遗民，以获得永恒的救恩。为使你们更明了所探讨的问题，我愿再提及蒙福的达味所说的其他话语，由此你们将知道：主被预言的圣神称为基督；万有的圣父主将祂从地上举起，使祂坐在自己的右边，直到使祂的仇敌作祂的脚凳。这实在是从我们的主耶稣基督死后复活升天之时开始发生的，现今的时期正走向终结；而达尼尔所预言的那位将统治一个时期、两个时期和半个时期的，甚至已在门口，即将对至高者说亵渎和狂妄的话。但你们因不知他将统治多久，而持另一种见解。因为你们将\Quote{一个时期}解释为一百年。若如此，那罪恶之人至少必须统治三百五十年，以便我们计算圣达尼尔所说的\Quote{两个时期}仅为两个时期。我以上所说全是离题之言，为使你们最终能被天主所宣告的话说服，即你们是愚昧的儿子；以及\BibleQuote{因此，看，我要继续夺去这百姓，并将他们夺去；我要使智者的智慧消失，隐藏精明人的聪明。}\BibleRef{依 29:14}请你们停止自欺并欺骗听你们的人，而向我们学习，我们是由基督的恩宠获得智慧的人。达味所说的话如下：
\end{quotation}

\BibleQuote{上主对我主说：祢坐在我右边，等我使祢的仇敌作祢的脚凳。上主必从熙雍伸出祢的权杖：祢要在祢仇敌中掌权。在祢的圣者光辉中，祢的百姓在祢能力之日，都甘愿献身。从曙光之前，我从母腹中生了祢。上主发了誓，决不反悔：祢是照默基瑟德的品位永为司祭。上主在祢右边：祂在发怒之日击碎了君王。祂要在外邦中审判，使尸首遍地。祂要喝路旁的溪水，因此必将抬起头来。}\par

% structure: p0008 source=h2 target=subsection reason=relative-heading-rank; base=h2

\subsection{第三十三章 第一一〇篇不是关于希则克雅。证明基督先卑微，后荣耀}

\Person{犹斯定}：我知道你们竟敢将这篇圣咏解释为指希则克雅王；但你们错了，我将立即从这些话本身向你们证明。经上说：\BibleQuote{上主发了誓，决不反悔}，以及\BibleQuote{祢是照默基瑟德的品位永为司祭}，连同前后的句子。连你们也不敢反对说希则克雅是司祭，或是天主的永司祭；但这些话是指我们的耶稣说的，这些表达已表明。但你们的耳朵闭塞，心也迟钝。因为藉着这句\BibleQuote{上主发了誓，决不反悔：祢是照默基瑟德的品位永为司祭}，天主以誓言（因你们不信）表明祂是大司祭，照默基瑟德的品位。即：正如梅瑟所描述的默基瑟德是至高者的司祭，他是未受割礼者的司祭，且祝福了给他献上什一之物的受割礼的亚巴郎；同样，天主表明祂的永司祭，也被圣神称为主，将是未受割礼者的司祭。那些在割礼中接近祂的人，即相信祂并寻求祂祝福的人，祂也必接纳并祝福。而且祂将先像人一样卑微，然后被高举，圣咏末尾的话显明了：\BibleQuote{祂要喝路旁的溪水}，然后\BibleQuote{因此必将抬起头来}。\par

% structure: p0010 source=h2 target=subsection reason=relative-heading-rank; base=h2

\subsection{第三十四章 第七二篇不适用于撒罗满，基督徒对撒罗满的过犯感到战栗}

\Person{犹斯定}：为使你们相信你们对圣经一无所知，我还要提醒你们另一篇圣咏，由圣神启示达味所作的，你们说是指你们的王撒罗满，但其实也是指我们的基督。但你们被含糊的言辞所欺骗。因为当它说\BibleQuote{上主的法律是完善的}时，你们不理解为梅瑟之后将颁布的法律，而是梅瑟所赐的法律，尽管天主宣告要订立新的法律和新的盟约。当它说\BibleQuote{天主，求祢将审判权赐给君王}时，既然撒罗满是王，你们就说这篇圣咏指他，但圣咏的话明确宣告是指永生的君王，即基督。因为基督是君王、司祭、天主、主、天使、人、元帅、磐石、所生的儿子、先受苦难、然后回到天上、再带着荣耀来临，并宣讲祂永生的王国：我从所有圣经中这样证明。为使你们明白我所说的话，我引述这篇圣咏的话如下：\par

\Emph{天主，请将祢的审判权授与君王，将祢的正义授与君王的儿子；使祂以正义审判祢的人民，以判断扶助祢的穷人。大山将给人民带来和平，小山带来正义。祂要审判百姓中的穷人，拯救贫苦者的子孙，并要压伏欺压者。祂要与日月共存，直到万代。祂要像雨露降在禾捆上，像甘霖滋润大地。在祂的日子，正义要兴盛，并充满和平，直到月亮被除去。祂的权柄要从这海到那海，从大河直到地极。埃塞俄比亚人要在祂面前俯伏，祂的仇敌要舔尘土。塔尔史士和群岛的君王要献上礼物，阿剌伯和舍巴的君王要进贡；众君王都要崇拜祂，万民都要事奉祂：因为祂从有权势者手中救出了穷人，和那没有人帮助的贫苦者。祂要怜恤穷人和贫苦者，并拯救贫苦者的灵魂：祂要从重利和不义中救赎他们的灵魂，祂的名在他们面前是尊贵的。祂要活着，人要把阿剌伯的金子献给祂，人们要不断为祂祈祷，终日赞美祂。在地面上要有根基，在高山之巅被高举：祂的果实要在黎巴嫩，城里的人要像地上的草一样茂盛。祂的名要永远受赞美。祂的名要在太阳前存留，地上的万族都要因祂蒙福，万民都要称祂有福。愿上主以色列的天主受赞美，唯独祂行奇事；愿祂荣耀的名永远受赞美，直到永远；愿全地充满祂的荣耀。阿们，阿们。}\par

在这篇我所引用的诗篇的结尾，写着：\Quote{耶西的儿子达味的圣咏到此为止。}此外，我知道\Person{撒罗满}是一位著名的伟大君王，他建造了那座称为\Place{耶路撒冷}的圣殿；但显然，诗篇中所提及的那些事没有一件发生在他身上。因为列王并没有都崇拜他；他也没有统治到地极；他的仇敌也没有在他面前俯伏舔土。不仅如此，我还敢重复记载在\WorkTitle{列王纪}中的他所犯下的事：他如何因一个妇人的影响而崇拜了\Place{西顿}的偶像；而那些通过被钉在十字架上的\Person{耶稣}认识\OriginalTerm{天主}{God}、万有的创造者的外邦人，却不敢这样做，反而忍受各种酷刑和报复，甚至死亡，也不肯崇拜偶像，或吃祭偶像的肉。\par

% structure: p0014 source=h2 target=subsection reason=relative-heading-rank; base=h2

\subsection{第三十五章 异端者坚定了公教徒的信德}

\Emph{\Person{特黎丰}：}我相信，然而，许多自称承认耶稣、被称为基督徒的人，吃祭偶像的肉，并声称他们因此丝毫未受损害。\par

\Emph{\Person{尤斯定}：}事实上，有这样一些人，他们承认自己是基督徒，并承认被钉十字架的耶稣是主和基督，却不教导祂的教义，而是教导错谬之神的教义，这使我们这些真正纯正教义的门徒更加坚定在祂所宣布的希望中。因为祂曾预言要在祂的名下发生的事，我们亲眼目睹正在实现。祂说：\BibleQuote{许多人要在我名里来，外面披着羊皮，里面却是凶残的豺狼。}\BibleRef{玛 7:15}又说：\BibleQuote{将会有分裂和异端。}\BibleRef{格前 11:19}又说：\BibleQuote{你们要提防假先知，他们来到你们跟前，外披羊毛，内里却是凶残的豺狼。}\BibleRef{玛 7:15}又说：\BibleQuote{将有许多假基督和假宗徒兴起，欺骗许多信徒。}\BibleRef{玛 24:11}因此，我的朋友们，现今和过去都有许多人，以耶稣的名出现，教导并做出亵渎和不敬的事；我们以这些教义和意见的创始人来称呼他们。（有些人以一种方式，另一些人以另一种方式，教导亵渎万物的创造者，以及祂预言要来的基督，以及\Person{亚巴郎}、\Person{依撒格}和\Person{雅各伯}的天主；我们与他们毫无共同之处，因为我们知道他们是无神论者、不敬、不义、有罪，只在名义上承认耶稣，却不崇拜祂。然而，他们自称基督徒，正如某些外邦人将天主的名字刻在自己手工作品上，参与邪恶不敬的仪式。）有些人被称为\OriginalTerm{马西昂派}{Marcians}，有些\OriginalTerm{瓦伦廷派}{Valentinians}，有些\OriginalTerm{巴西里德派}{Basilidians}，有些\OriginalTerm{撒督尼派}{Saturnilians}，还有其他名称；每个人按照其意见创始人的名字命名，正如我先前所说，每个自认为哲学家的人，都认为自己必须从特定教义之父的名字中取用他所遵循的哲学名称。因此，由于这些事件，我们知道耶稣预知祂之后会发生的事，也由于许多其他事件，祂预言了那些相信并承认祂是基督的人将要遭遇的事。因为我们所遭受的一切，甚至被朋友杀害，祂都预言会发生；所以祂的言语和行为都无可指摘。因此，我们为你们和所有仇恨我们的人祈祷，希望你们与我们一同悔改，不再亵渎那位凭祂的作为、如今仍通过祂的名所行的奇迹、祂所教导的话、以及关于祂所宣布的预言，而证明自己是无可指摘、全然纯全的基督耶稣；而是相信祂，在祂荣耀的第二次来临时获得救恩，而不致被祂丢进火里受罚。\par

% structure: p0017 source=h2 target=subsection reason=relative-heading-rank; base=h2

\subsection{第三十六章 他证明基督被称为万军之主}

\Emph{\Person{特黎丰}：}就按你说的那样吧——即预言基督要受苦，并被称作石头；在祂第一次显现（预告祂要受苦）之后，祂将在荣耀中来临，并最终成为审判者，永恒的君王和司祭。现在证明此人是否就是这些预言所指向的那位。\par

\Person{犹斯定}：正如你所愿，特黎丰，我将在适当的地方提出你所寻求的这些证据；但现在请允许我先复述那些预言，我这样做是为了证明基督被圣神以比喻称为天主、万军的上主和雅各伯。而你们的解释者，正如天主所说，是愚昧的，因为他们说那是指撒罗满而不是基督，当时他将约柜抬进他所建造的圣殿。达味的圣咏如下：\par

\BibleQuote{大地和其中的万物属于上主，世界和其间的居民属于上主。祂将地建立在海上，安置在江河之上。谁能登上上主的山？谁能站在祂的圣所？那手洁心清，不慕虚幻，不发假誓的人。他必蒙上主降福，由拯救他的天主获得正义。这就是寻求上主、求见雅各伯天主面容的族类。城门，请提高你们的门楣；古老的门户，请高抬起来；光荣的君王要进来。谁是这位光荣的君王？是强劲有力的上主，是作战有能的上主。城门，请提高你们的门楣；古老的门户，请高抬起来；光荣的君王要进来。谁是这位光荣的君王？万军的上主，祂就是光荣的君王。}\par

因此，可见撒罗满并非万军的上主；但当我们基督从死者中复活并升天时，天上的统治者们，奉天主的命令，被吩咐打开天堂的门，好让那位光荣的君王进入，并在升天后坐在圣父的右边，直到祂使仇敌成为祂的脚凳，正如另一篇圣咏所显示的。因为当天上的统治者们看见祂容貌丑陋、受人轻慢、毫无光荣时，他们没有认出祂，便问：\BibleQuote{谁是这位光荣的君王？}圣神或代表圣父，或代表自己，回答他们说：\BibleQuote{万军的上主，祂就是这位光荣的君王。}因为每个人都会承认，耶路撒冷圣殿的守门者，没有一个敢指着撒罗满——尽管他是如此光荣的君王——或指着约柜说：\BibleQuote{谁是这位光荣的君王？}\par

% structure: p0022 source=h2 target=subsection reason=relative-heading-rank; base=h2

\subsection{第三十七章  同样由其他圣咏证明}

\Person{犹斯定}：此外，在第四十六篇的\OriginalTerm{细拉}{diapsalm}中，也同样论及基督说：\BibleQuote{天主上升，有欢呼之声；上主上升，有号角之声。你们应向天主歌颂，歌颂；向我们的君王歌颂，歌颂；因为天主是全地的君王，你们要歌颂圣咏。天主作王统治万邦，天主坐在祂的圣座上。万邦的领袖聚集在一起，作为亚巴郎天主的子民；因为地上的大能者属于天主，祂极受尊崇。}在第九十八篇中，圣神责备你们，并预言那位你们不愿意祂为王的，乃是君王与主，掌管撒慕尔、亚郎、梅瑟，以及所有其他人。这篇圣咏的话如下：\par

\BibleQuote{上主为王，万邦应战栗；祂坐在革鲁宾之上，大地应震动。上主在熙雍伟大，崇高超越万民。愿他们赞美你伟大的名号，因为它是可敬可畏的；王的权柄喜爱公正。你奠定了正直；你在雅各伯中施行了公正和仁义。请尊崇上主我们的天主，向祂的脚凳下拜，因为祂是神圣的。梅瑟和亚郎是祂的司祭，撒慕尔是呼求祂名的人。他们（圣经说）呼求上主，祂就应允了他们。祂在云柱中向他们说话；因为他们遵守了祂的规诫，以及祂所赐给他们的诫命。上主我们的天主，你应允了他们；天主，你宽恕了他们，却惩罚了他们的一切恶行。请尊崇上主我们的天主，向祂的圣山下拜，因为上主我们的天主是神圣的。}\par

% structure: p0025 source=h2 target=subsection reason=relative-heading-rank; base=h2

\subsection{第三十八章  犹太人因基督被称为朝拜的对象而感到烦恼。犹斯定却由第四十五篇予以证实}

\Person{特黎丰}：先生，如果我们听从我们的老师，那对我们有好处，他们曾立下法律，不准我们与你们有任何来往，也不准我们与你们讨论这些问题。因为你们说了许多亵渎的话，竟然试图说服我们，说这位被钉十字架的人曾与梅瑟和亚郎同在，并在云柱中向他们说话；然后祂成了人，被钉十字架，升天，还要再来到地上，并且应当受朝拜。\par

\Person{犹斯定}：我知道，正如天主的话所说，天主的伟大智慧，万物的创造者、全能者，对你们是隐藏的。因此，我出于同情，极力使你们理解这些对你们来说似乎矛盾的事；如果不行，至少使我在审判之日得以无罪。因为你们会听到其他更显得矛盾的话；但不要惊慌，反而更要热心地听和探究，藐视你们老师的传统，因为他们被圣神证明不能领悟天主所教导的真理，而宁愿传授自己的教义。同样，在第四十四（四十五）篇中，这些词也论及基督：\par

\Quote{我的心脏涌出美辞；我向君王呈献我的作品。我的舌头像熟练的书记的笔。祢比世人更美：恩宠倾注在祢唇上，因此天主永远祝福了祢。大能者啊，佩剑在腰间。愿祢在威严与荣美中前进，为了真理、温良和正义而胜利为王；愿祢的右手施行可惊的事。祢的箭矢锋利，大能者啊；万民必仆倒在祢脚下；在君王仇敌的心中［箭矢钉入］。天主，祢的宝座永永远远；祢王国的权杖是正直的权杖。祢喜爱正义，憎恨不义；因此天主，祢的天主，用喜乐油膏了祢，胜过祢的同伴。［祂膏了祢］用没药、香料和肉桂，从祢的衣裳；从象牙宫中，他们使祢喜乐。君王的女儿们在祢尊荣中。王后穿戴金绣衣裳，站在祢右边。女子啊，请听，请看，侧耳细听，忘记你的百姓和你的父家；君王必恋慕你的美貌，因为祂是你的主，他们也要敬拜祂。提洛的女儿必带着礼物前来。民中的富足人要恳求祢的面。君王女儿的一切荣耀在于内里，身穿锦绣衣裳。随从她的童女要带到君王面前；她的同伴要引到祢面前；她们必欢喜快乐被引导，进入君王的圣殿。代替你的祖先，你的子孙出生了；祢要立他们作全地的首领。我要在万代中记念祢的名；因此万民要永远称谢祢，直到永远。}\par

% structure: p0029 source=h2 target=subsection reason=relative-heading-rank; base=h2

\subsection{第三十九章　犹太人憎恨相信此事的基督徒。两者之间有极大的区别！}

\Person{犹斯定}：如今，你们憎恨我们这些持这些见解并控告你们心地刚硬的人，并不奇怪。因为厄里亚与天主论及你们时，这样说：\Quote{主，他们杀了祢的先知，拆毁了祢的祭台；只剩下我一个人，他们还要寻索我的命。}祂回答他说：\Quote{我仍为自己留下了七千人，未曾向巴耳屈膝。}因此，正如天主因那七千人没有发怒，同样，祂如今也尚未施行审判，也不施行，因为祂知道每天有些［你们中的］人因基督的名成为门徒，离开错误的道路；他们也领受恩赐，各人按自己的价值，被这基督的名光照。因为有人领受明悟之神，有人领受谋略之神，有人领受能力之神，有人领受治愈之神，有人领受预知之明，有人领受教导之神，有人领受敬畏天主之神。\par

\Person{特黎丰}：我希望你知道，你说出这些见解，是神志不清了。\par

\Person{犹斯定}：朋友，请听，因为我并非疯狂或神志不清；但预言说，基督升天以后，祂将使我们脱离错误，赐给我们恩赐。经文这样说：\Quote{祂升上高天，掳掠了俘虏，将恩赐赏给人。}因此，我们这些从升上高天的基督领受了恩赐的人，从预言的话证明你们\Quote{自以为有智慧，自视为明达的人}\BibleRef{依 5:21}是愚昧的，只用嘴唇尊敬天主和祂的基督。但我们这些受教于全部真理的人，以行为、知识和心，甚至至死，尊敬他们两位。但你们却犹豫不承认祂是基督，如同圣经和以祂的名所见证的事件和行为所证明的，或许是因为这个原因：恐怕被统治者迫害，他们受那邪恶诡诈的蛇的灵的影响，必不断杀害和迫害那些承认基督之名的人，直到祂再来，消灭他们全体，并照各人的行为报应各人。\par

\Person{特黎丰}：那么，请给我们证明，这个你们说被钉十字架并升天的人，就是天主的基督。因为你已经通过你之前引用的圣经充分证明，圣经中宣告基督必须受难，并带着荣耀再来，并领受万国的永恒王国，万国都要臣服于祂；现在请向我们证明这人就是祂。\par

\Person{犹斯定}：先生们，对有耳朵的人来说，甚至从你们所承认的事实中，这事已经证明了；但为使你们不以为我困惑无能，无法像我所应许的那样给出你们所问的证明，我将在适当的地方去做。现在，我继续讨论我刚才正在论述的主题。\par

% structure: p0035 source=h2 target=subsection reason=relative-heading-rank; base=h2

\subsection{第四十章　他回到梅瑟的律法，并证明它们是关乎基督之事的预像}

\Person{犹斯定}：那么，天主命令作为逾越节祭献的羔羊的奥秘，是基督的预像；凡信祂的人，按其对祂的信德，用祂的血涂抹自己的房屋，即他们自己。因为天主所创造的受造物——即亚当——是出自天主的圣神之房屋，这是你们都能理解的。我这样证明这条命令是暂时的：天主不许逾越节的羔羊在祂的名被称呼以外的任何地方祭献，因为祂知道，在基督受难之后，日子将到，连耶路撒冷那地方也要交给你们的敌人，总之，所有祭献都要停止；而那只被命令完全烤熟的羔羊，是基督将要受的十字架苦难的象征。因为烤熟的羔羊，是烤成十字架形状而预备的：一根叉从下部直贯到头部，另一根横贯背部，羔羊的腿就系在上面。那在斋戒期间被命令献上的两只山羊，一只作为替罪羊被送走，另一只被祭献，同样宣告了基督的两次来临：第一次，你们百姓的长老和司铎们按手在祂身上，将祂处死，把祂当作替罪羊送走；第二次来临，因为在耶路撒冷同一地方，你们将认出你们所侮辱的那一位，祂是甘愿悔改的一切罪人的祭献，并且遵守依撒意亚所说的斋戒，解除暴力的契约，并遵守他所列举的其他诫命，这些我都引用过，相信耶稣的人确实遵守。此外，你们知道，在斋戒期间被命令献上的两只山羊，也不允许在其他地方同样献上，而只能在耶路撒冷。\par

% structure: p0037 source=h2 target=subsection reason=relative-heading-rank; base=h2

\subsection{第四十一章：细面的祭献是圣体圣事的预像}

\Person{犹斯定}：诸位，那为癞病人得洁净而规定献上的细面祭献，是圣体圣事之饼的预像；我们的主耶稣基督规定了这圣事的举行，以纪念祂为那些灵魂从一切不义中得洁净的人所受的苦难，为使我们同时感谢天主，因为祂为了人的缘故创造了世界和其中的万物，将我们从所处的邪恶中解救出来，并借着按祂的旨意受苦的那一位，彻底摧毁了掌权者和权势者。因此，天主借着十二位先知之一玛拉基亚的口，论及那时你们所献的祭献，如我先前所说：\InnerQuote{我不喜欢你们，上主说；我也不从你们手中接纳你们的祭献：因为从日出到日落，我的名在外邦人中受尊崇，处处有人向我的名献香和洁净的祭品：因为我的名在外邦中大受尊崇，上主说：但你们亵渎它。}\BibleRef{拉 1:10}所以祂接着论及那些外邦人，即我们，我们在各处向祂献祭，即圣体圣事之饼和圣体圣事之杯，肯定我们尊崇祂的名，而你们亵渎它。至于割损的诫命，命令在第八天给孩子行割损，这是真正割损的预像，借着那在安息日后第一天从死者中复活的一位，即我们的主耶稣基督，我们从诡诈和不义中被割损。因为安息日后第一天，虽是所有日子中居首的，但按一周全部日子的数目，却被称为第八天，并且仍然是第一天。\par

% structure: p0039 source=h2 target=subsection reason=relative-heading-rank; base=h2

\subsection{第四十二章：司祭袍上的铃铛是宗徒的预像}

\Person{犹斯定}：此外，大司祭袍上系有十二个铃铛，垂到脚边，这是十二位宗徒的象征，他们依靠基督——永恒司祭——的能力；借着他们的声音，普世充满了天主和祂的基督的光荣与恩宠。为此，达味也说：\InnerQuote{他们的声音传遍普世，他们的言语达到地极。}依撒意亚也说话，好像亲口代表宗徒们，当他们向基督说，他们相信的不是自己的报告，而是那派遣他们者的能力。因此他说：\InnerQuote{主，谁相信了我们的报告？上主的臂膀向谁启示了？我们在祂面前宣讲，好像一个孩童，好像干地上的根。}\BibleRef{依 53:1}（以及那段已引用的预言中接着的话。）但当经文从多人口中说\InnerQuote{我们在祂面前宣讲}，并加上\InnerQuote{好像一个孩童}时，它预示恶人将服从祂，听从祂的命令，大家都成为\Emph{一}个孩童。这种情况你们可以在身体中看到：虽然肢体被数算为许多，却都称为\Emph{一}，并且是一个\Emph{身体}。因为，的确，一个团体和一个教会，虽然人数众多，实际上如同一体，用一个称呼呼唤和称呼。总之，诸位，我可以列举梅瑟的一切其他规定，证明它们都是预像、象征和宣告：那些将要发生在基督身上的事，那些预知要相信祂的人，以及那些基督自己要完成的事。但既然我现在所列举的在我看已足够了，我便回到讨论的顺序上。\par

% structure: p0041 source=h2 target=subsection reason=relative-heading-rank; base=h2

\subsection{第四十三章：他总结说，法律在基督内结束，基督由童贞女所生}

\Person{犹斯定}：如割损礼始于亚巴郎，安息日、祭献、供物与庆节始于梅瑟，并且已证明这些是因你们人民心硬而施行的，因此，按照圣父的旨意，它们理应在由童贞女所生、出自亚巴郎家族、犹大支派及达味的祂内、在基督天主子内获得终结；祂被宣告将来到全世界，成为永远的律法与永远的盟约，正如前述预言所示。我们藉着祂接近天主的人，领受了不是属肉体的，而是属神的割损礼，就是哈诺客及类似他之人所遵守的。我们藉着圣洗领受了它，因为我们原是罪人，出于天主的仁慈；众人皆可同样获得。但既然祂诞生的奥秘如今需要我们的注意，我将讲述此事。依撒意亚论及基督的诞生时，在已引用的话中断言：\BibleQuote{谁能述说祂的世代？因为祂的生命从地上被夺去；因我人民的过犯，祂被处死。} \BibleRef{依 53:8} 预言之神由此肯定，那位要死去的、使我们罪人因祂的创伤而得痊愈者的诞生，是无法言说的。此外，为使信祂的人得知祂进入世界的方式，预言之神藉同一位依撒意亚预言了此事将如何发生：\par

\BibleQuote{\Emph{主又对阿哈次说：\Quote{你向你的上主天主要求一个征兆，或求诸阴府，或求诸高天。}阿哈次说：\Quote{我不要求，我不试探上主。}依撒意亚说：\Quote{达味家，你们听着：你们使人厌烦还算小事，还要使我的天主厌烦吗？因此主自己要给你们一个征兆：看，童贞女将怀孕，将生一个儿子，他的名字要叫厄玛奴耳。在他知道或拒绝邪恶、选择良善之前，他要吃奶油和蜂蜜；因为孩子知道好坏之前，他藉选择良善而拒绝邪恶。因为在孩子知道如何呼唤父亲或母亲之前，他将获得大马士革的力量和撒玛黎雅的掠物，在亚述王面前。那地将被遗弃，你们将因两位国王的临在而难以忍受。但天主将使你和你的百姓及你父家，临到从厄弗辣因从犹大夺走亚述王之日起所未曾临到的日子。}}}\par

如今众人皆知，在按血统属于亚巴郎的后裔中，除了我们的这位基督外，没有一人由童贞女所生，或据说曾由童贞女所生。但你们和你们的教师竟敢断言，在依撒意亚的预言中说的不是\BibleQuote{看，童贞女将怀孕}，而是\BibleQuote{看，年轻女子将怀孕，生一个儿子}；并且你们将预言解释为似乎是指你们的君王希则克雅。我将尽力就此点简短地与你们辩驳，并表明其中指的是我们所承认为基督的那一位。\par

% structure: p0045 source=h2 target=subsection reason=relative-heading-rank; base=h2

\subsection{第四十四章　犹太人徒然期望救恩，惟藉基督才能获得}

\Person{犹斯定}：因为就你们而言，如果我努力通过提出证据来说服你们，我将在各方面显得无罪。但如果你们因基督徒的命运——死亡——而心硬或意志薄弱，不愿赞同真理，你们将成为自身灾祸的制造者。你们自欺欺人，自以为因按血统是亚巴郎的后裔，就必完全承受由天主藉基督所应许的美物。因为任何人，甚至他们中的任何人，都无可指望，只有那些在心灵上与亚巴郎的信德相似、并认识了一切奥秘的人，才能指望。因为我说，有些诫命是关于敬拜天主和实践正义的；但有些诫命和行为也同样提及基督的奥秘，是因你们人民的心硬。这一点，天主在厄则克耳中指明，论到此事说：\BibleQuote{纵然诺厄、雅各伯和达尼尔求问儿子或女儿，他们的请求也不会得允。} \BibleRef{则 14:20} 在依撒意亚中，关于同一件事祂这样说：\BibleQuote{上主天主说：他们都要出去，观看那些背叛我的人的尸首。因为他们的虫永不灭，他们的火永不熄，他们成为一切血肉之躯的惊骇。} \BibleRef{依 66:24} 因此，你们理应从你们的心灵中根除这种希望，并赶紧认识罪的赦免和承受应许美物的希望如何属于你们。但没有别的途径——只有认识这位基督，在依撒意亚所说的泉源中受洗以得罪赦，并且在其余方面过无罪的生活。\par

% structure: p0047 source=h2 target=subsection reason=relative-heading-rank; base=h2

\subsection{第四十五章　在律法以前和律法之下的义人将由基督得救}

\Person{特里丰}：如果我似乎打断这些你认为必须探究的事情，但我要提出的问题是迫切的。请先允许我。\par

\Person{犹斯定}：请问你任意想问的，就如你所想到的；我将在问答之后努力恢复并完成讲论。\par

\Person{特里丰}：那么请告诉我，那些按梅瑟所赐的律法生活的人，在死人复活时，将与雅各伯、哈诺客和诺厄同样地生活，还是不会？\par

\Emph{犹斯定：}我引用了\Person{厄则克耳}所说的话，即\Quote{即使\Person{诺厄}、\Person{达尼尔}和\Person{雅各伯}恳求儿女，这请求也不会被应允}，而是说每一个人将因自己的义德而得救。我又说，那些按梅瑟法律调整生活的人也将同样得救。因为梅瑟法律中自然善、虔敬、正义的事，且被吩咐给遵守者去做；以及因百姓的心硬而被规定要履行的事，同样被记录，也被那些在法律之下的人所实行。既然行那普遍、自然、永恒之善的人中悦天主，他们将在复活中藉着这位\Person{基督}而得救，与在他们之前的义人，即\Person{诺厄}、\Person{哈诺客}、\Person{雅各伯}，以及其他所有人，连同那些认识这位\Person{基督}——\Person{天主子}，祂在晨星和月亮之前就已存在，且屈尊降生成人，由\Person{达味}家族的这个童贞女所生——的人一样，为的是藉着这安排，从起初犯罪的蛇和与他相似的使者被消灭，死亡被藐视，并在\Person{基督}亲自第二次来临时永远离开那些信祂并度相称生活的人，且不再存在：那时，一些人被罚入永不间断的审判和火刑中，另一些则存在于不受苦、不腐朽、不忧伤和不朽之中。\par

% structure: p0052 source=h2 target=subsection reason=relative-heading-rank; base=h2

\subsection{第四十六章 \Person{特黎丰}问：如今仍遵守法律的人能否得救。\Person{犹斯定}证明这行为对义德毫无助益}

\Emph{特黎丰：}但若有人如今仍愿遵守梅瑟所赐的规条，却信这被钉的\Person{耶稣}，承认祂是天主的\Person{基督}，且全权审判者已赐给祂，祂的王国是永远的，这样的人也能得救吗？\par

\Emph{犹斯定：}让我们一同考虑这事：人如今能否遵守所有梅瑟的规定呢？\par

\Emph{特黎丰：}不能。因为我们知道，如你所说，无论在何处都不可能献逾越节羔羊，也不能献为禁日所命献的公山羊，简言之，不能呈献所有其他的祭品。\par

\Emph{犹斯定：}那么，我求你自己说说哪些是可以遵守的；因为你将被说服，即一个人即使不遵守或没有履行永久的法令，也必能得救。\par

\Emph{特黎丰：}守安息日、受割损、守月份、凡接触梅瑟禁止之物或房事后要洗濯。\par

\Emph{犹斯定：}你以为\Person{亚巴郎}、\Person{依撒格}、\Person{雅各伯}、\Person{诺厄}、\Person{约伯}以及在他们前后的一切同样义人，还有\Person{亚巴郎}的妻子\Person{撒辣}、\Person{依撒格}的妻子\Person{黎贝加}、\Person{雅各伯}的妻子\Person{辣黑耳}和\Person{肋阿}，以及直到忠信仆人\Person{梅瑟}的母亲在内的所有人，他们都没有遵守这些（规条），会得救吗？\par

\Emph{特黎丰：}难道\Person{亚巴郎}和他的后裔不都受割损了吗？\par

\Emph{犹斯定：}我知道\Person{亚巴郎}和他的后裔受割损了。割损赐给他们的理由在前面已经详细陈述；如果之前说的不能说服你，让我们再探讨这问题。但你知道，直到\Person{梅瑟}，事实上没有任何义人遵守我们所谈论的任何这些仪礼，也没有领受一条命令去遵守，除了一开始从\Person{亚巴郎}开始的割损。\par

\Emph{特黎丰：}我们知道，也承认他们得救。\par

\Emph{犹斯定：}你看，\Person{天主}藉\Person{梅瑟}将这一切规条加在你们身上，是因你们百姓的心硬，为的是藉着众多规条，使你们常在一切行动中记念\Person{天主}，永不开始行不义或不虔敬的事。因为祂命令你们在衣服四边加上紫色穗头\BibleRef{户 15:38}，好使你们不忘记\Person{天主}；又吩咐你们佩带经文匣\BibleRef{申 6:6}，将某些字符刻在极薄的羊皮纸上，我们视之为圣物；藉这些手段激发你们不断记念\Person{天主}，同时却说服你们，你们心里连对\Person{天主}崇拜的一丝记忆都没有。然而即便如此，你们仍未从偶像崇拜中回转；因为在\Person{厄里亚}的时代，当祂数点那些没有向\Person{巴耳}屈膝的人数时，祂说数目是七千；在\Person{依撒意亚}中，祂谴责你们将儿女祭献给偶像。但我们，因为我们拒绝祭献我们从前惯常祭献的对象，而承受极端的刑罚，并因死亡而喜乐——相信\Person{天主}必藉祂的\Person{基督}复活我们，使我们不朽坏、不受搅扰、不死；并且我们知道，因你们百姓的心硬而规定的那些规条，对实行义德和虔敬毫无助益。\par

% structure: p0063 source=h2 target=subsection reason=relative-heading-rank; base=h2

\subsection{第四十七章 \Person{犹斯定}与遵守法律的基督徒共融。不少天主教徒却另作选择}

\Emph{特黎丰：}但若有人知道是这样，在认出此人是\Person{基督}后，相信并服从祂，却仍愿意遵守这些（规条），他会得救吗？\par

\Emph{犹斯定：}依我看来，\Person{特黎丰}，这样的人会得救，只要他不竭力说服其他人——我指的是那些曾被\Person{基督}从错谬中挽救的外邦人——也遵守与他相同的事，并告诉他们若不如此就不得救。你自己在讨论开始时就如此做了，当时你声称我若不遵守这些规条就不能得救。\par

\Emph{特黎丰：}那你为什么说\Quote{依我看来，这样的人会得救}，除非有些人声称这样的人不会得救？\par

\Emph{犹斯定：} 是有这样的人，\Person{特黎丰}，他们不敢与这样的人有任何来往或接待他们；但我不同意他们。但若有人，因着心性软弱，愿意遵守\Person{梅瑟}所赐的那些制度，并期望从中获得某种德行——虽然我们相信这些制度是因人民心硬而设立的——同时怀有对这位基督的望德，并履行永恒而自然的正义与虔诚的行为，却又愿意与基督徒和信徒共同生活，如我先前所说，并不诱导他们像自己一样受割损，或守安息日，或遵守任何其他此类仪式，那么我认为我们应当接纳这样的人，并以亲属和弟兄的身份与他们共同交往。但是，\Person{特黎丰}，若你们中间有些人，自称相信这位基督，却强迫那些相信这位基督的外邦人完全按照\Person{梅瑟}所赐的法律生活，或者选择不与那些人有如此亲密的交往，我同样不赞同他们。然而我相信，就连那些被他们说服，在宣认天主在基督内的同时，也遵守法律安排的人，也可能会得救。此外，我坚持认为，那些宣认并知道这人就是基督，却因某种原因退回到法律安排，并否认这人是基督，且在死前没有悔改的人，必不能得救。进一步说，我坚持认为，\Person{亚巴郎}的后裔中那些按法律生活，却在死前不信这位基督的人，同样不能得救，尤其是那些在会堂里诅咒并仍在诅咒这位基督，以及一切能使他们获得救恩并逃脱烈火惩罚之事的人。因为天主的良善和慈爱，以及祂无限的富饶，使那如\Person{厄则克耳}所讲悔改自己的罪的人，被视为义人和无罪者；而那从虔诚和正义堕入不义和不虔敬的人，则被算作有罪、不义和不敬。因此，我们的\Person{主耶稣基督}也曾说：\Quote{无论我在什么事上找到你们，就在那些事上审判你们。}\par

\SourceRecord{Translated by Marcus Dods and George Reith.；Ante-Nicene Fathers；Vol. 1.；Edited by Alexander Roberts, James Donaldson, and A. Cleveland Coxe.；Buffalo, NY: Christian Literature Publishing Co.,；1885.；<http://www.newadvent.org/fathers/01283.htm>.}

% structure: p0001 source=h1 target=section reason=page-title

\section{与特黎丰对话（第四十八至五十四章）}

% structure: p0002 source=h2 target=subsection reason=relative-heading-rank; base=h2

\subsection{第四十八章 在证明基督的天主性之前，他（特黎丰）要求先确定祂是基督}

\Emph{\Person{特黎丰}：}我们已听了你对这些事的看法。请从你中断处继续论述，并把它结束。因为其中一部分在我看来似是而非，完全无法证明。因为当你说这位基督在万世以前就作为天主存在，然后祂屈尊降生成为人，却又不是由人所生，这个论断在我看来不仅似是而非，而且是愚蠢的。\par

\Emph{\Person{儒斯定}：}我知道这说法确实似是而非，尤其对你们民族的人而言，你们总不愿理解或履行天主的诫命，却乐于遵行你们师傅的教训，正如天主亲自所宣布的。\BibleRef{依 29:13} 然则，特黎丰，此人即是天主的基督这一证明并不落空，即使我无法证明祂以前就作为万物的创造主之圣子而存在，身为天主，并因童贞玛利亚而降生为人。但是既然我已经确实证明了此人是天主的基督，不论祂是谁，即使我未能证明祂先存，并按照圣父的旨意屈尊降生成为具有与我们同样情感且拥有身体的人，那么只有在这最后一点上说我错了才是合理的，而非否认祂是基督，即使祂看来是由人所生，并且仅仅证明了祂是因拣选而成为基督。因为我的朋友们，我们民族中有些人承认祂是基督，却认为祂是由人所生；我不同意他们的观点，即使现在大多数与我意见相同的人这样说，我也不会同意；因为基督亲自吩咐我们，不要相信人的教义，而要相信蒙福的众先知所宣讲、由祂本人所教导的教义。\par

% structure: p0005 source=h2 target=subsection reason=relative-heading-rank; base=h2

\subsection{第四十九章 对于反对说厄里亚尚未来临的人，他回答说他是第一次来临的前驱}

\Emph{\Person{特黎丰}：}那些肯定祂是一个人，因拣选而被傅油，然后成为基督的人，在我看来比你所持的那些意见更为合理。因为我们都期望基督将是一个由人所生的人，并且厄里亚来时会傅油祂。但如果此人似乎是基督，那么他当然必须被确认为由人所生的人；但鉴于厄里亚尚未来临，我推断此人并非祂（基督）。\par

\Emph{\Person{儒斯定}：}圣经不是在匝加利亚书中说，厄里亚将在主伟大而可畏的日子之前来临吗？\BibleRef{拉 4:5}\par

\Emph{\Person{特黎丰}：}诚然。\par

\Emph{\Person{儒斯定}：}因此，如果圣经强制你承认基督的两次来临被预言要发生——一次祂将显现为受苦、受辱、毫无威仪；另一次祂将光荣地来临，审判众人，正如在许多前述经文中已经显明的——难道我们不应该认为，天主的话已经宣告厄里亚将是这伟大而可畏之日的先驱，即祂第二次来临的先锋吗？\par

\Emph{\Person{特黎丰}：}诚然。\par

\Emph{\Person{儒斯定}：}因此，我们的主在教导中宣布此事必要发生，说厄里亚也必会来。我们知道这事要在我们的主耶稣基督从天上光荣降临时发生；祂的第一次显现由曾在厄里亚内的天主之神作为先锋，藉着若翰（一位你们民族的先知）传达；此后你们中间再没有出现其他先知。他坐在约但河旁呼喊说：\Quote{我用水洗你们为悔改；但有一位比我更强的将要来，我就是替祂提鞋也不配；祂要用圣神和火洗你们；祂手里拿着簸箕，要扬净祂的禾场，把麦子收进仓里，把糠用不灭的火烧尽。}\BibleQuote{\BibleRef{玛 3:11}} 你们王黑落德曾把这位先知关在监里；当他的生日庆祝时，这位黑落德的侄女跳舞讨他欢心，他对她说随意求什么。那少女的母亲唆使她求问在监里的若翰的头；她求了后，黑落德就派人把若翰的头放在盘子里拿来。因此，我们的基督在地上时，对那断言厄里亚必须在基督之前来的人说：\Quote{厄里亚的确要来，并要恢复一切；但我告诉你们，厄里亚已经来了，他们却不认识他，反而任意待了他。}\BibleQuote{\BibleRef{玛 17:12}} 经上又记着：\Quote{门徒这才明白耶稣所说的是指洗者若翰。}\par

\Emph{\Person{特黎丰}：}这个说法在我看来也似是而非，即那曾在厄里亚内的天主的预言之神，也在若翰内。\par

\Emph{\Person{儒斯定}：}难道你不认为同样的事发生在纳维的儿子若苏厄身上吗？他在梅瑟之后继承管理百姓之职，当时梅瑟奉命按手在若苏厄身上，天主对他说：\Quote{我要从你的身上取一些神，放在他身上。}\par

\Emph{\Person{特黎丰}：}诚然。\par

\Emph{\Person{儒斯定}：}因此，正如当梅瑟还在人间时，天主取了梅瑟身上的神放在若苏厄身上，同样，天主也能够使厄里亚的神降临在若翰身上；为的是，正如基督首次来临显现为无荣耀，同样，那在厄里亚内始终纯洁如基督之神的神首次来临也被视为无荣耀。因为主说祂要用隐藏的手与阿玛肋克争战；你也不会否认阿玛肋克败亡了。但如果说只有在基督荣耀的来临时与阿玛肋克争战，那么经上说\Quote{天主要用隐藏的手与阿玛肋克争战}将如何应验呢！你能看到天主隐藏的能力在被钉十字架的基督内，众魔鬼和世间所有掌权者与权势在祂面前战栗。\par

% structure: p0016 source=h2 target=subsection reason=relative-heading-rank; base=h2

\subsection{第五十章 从依撒意亚证明若翰是基督的前驱}

\Person{Trypho}：在我看来，你似乎曾与许多人在我们所探究的所有问题上进行了激烈的争论，因此对向你提出的所有问题都随时准备回答。那么，请先回答我：你如何能证明除了万有的创造者之外还有另一位天主？然后你才能证明祂甘愿由童贞女所生。\par

\Person{Justin}：首先请允许我引用以赛亚预言中的某些段落，这些段落关于洗者若翰——这位在我们主耶稣基督之前的先知——所履行的前驱职责。\par

\Person{Trypho}：我允许。\par

\Person{Justin}：以赛亚如此预言若翰的前驱工作：\BibleQuote{希则克雅对依撒意亚说：主所说的话是好的：愿在我的日子有和平与公义。\BibleRef{依 39:8} 并且，安慰我的百姓，你们司铎，向耶路撒冷的心说话，安慰她，因为她卑微的日子已满。她的罪孽已赦免；因她从主手中为她所有的罪受了双倍的惩罚。在旷野中有呼喊者的声音：预备主的道，修直我们天主的路径。一切山谷将被填满，一切山岳和丘陵将被削平；弯曲的将被修直，崎岖的将成为平坦的道路；主的荣耀将显现，凡有血肉的都将看见天主的救恩：因为主已说了。有声音说：呼喊！我说：我呼喊什么？凡有血肉的如同草，人的一切荣耀如同草上的花。草枯干，花凋谢；但主的话永远立定。报喜讯给熙雍的啊，请登上高山！报喜讯给耶路撒冷的啊，请用力扬声！扬声，不要害怕！向犹大各城报告：看，你们的天主！看，主带着大能而来，祂的手臂带着权能。看，祂的赏报在祂面前，祂的功业在祂前面。祂如牧人牧放自己的羊群，用臂膀集合羔羊，抚慰那怀孕的。谁用手量了水，用掌量了天，用拳量了地？谁称了山岳，把山谷放在天平上？谁认识了主的心？谁作过主的谋士，谁曾指教祂？祂与谁商议，谁教导了祂？谁给祂指明了判断？谁使祂知道智慧的道路？万民像桶中的一滴水，像天平上的一粒灰尘，被看作唾沫。黎巴嫩不足以作祭柴，其走兽不足以作燔祭；万民都归于无有，被视为虚无。\BibleRef{依 40:1}}\par

% structure: p0021 source=h2 target=subsection reason=relative-heading-rank; base=h2

\subsection{第五十一章 证实这预言已经应验}

\Person{Trypho}：先生，你重复的预言所有话语都是模棱两可的，对于证明你所想要证明的没有力量。\par

\Person{Justin}：如果先知们没有停止，以至于在你民族中，特黎弗，在这若翰之后不再有先知，那么我关于耶稣基督所说的可能被视为模棱两可。但如果若翰先来，呼唤人悔改，而基督在若翰仍坐在约旦河边时到来，终结了他的预言和洗礼，并且亲自宣讲说天国近了，祂必须从经师和法利塞人受许多苦，被钉十字架，第三天复活，将再次在耶路撒冷显现，再次与门徒们一起饮食；并且预言在祂两次来临之间的时期，如我先前所说，假司铎和假先知将奉祂的名兴起，这些事确实发生了；那么它们又怎能是模棱两可的呢？何况你还可以被事实说服？此外，祂还提到在你民族中不再有先知，并且人们认识到天主先前宣布要颁布的新约那时已经临在，即基督自己；祂用以下的话说：\BibleQuote{法律和先知直到洗者若翰；从那时起，天国受暴力，暴力的人夺取它。如果你们能接收，他就是那要来的厄里亚。有耳可听的，应当听。\BibleRef{玛 11:12}}\par

% structure: p0024 source=h2 target=subsection reason=relative-heading-rank; base=h2

\subsection{第五十二章 雅各伯预言基督的两次来临}

\Emph{\Person{尤斯定}：}圣祖\Person{雅各伯}曾预言\Person{基督}有两次来临，并且首次祂将受苦，在祂来临之后，你们民族中不再有先知或君王（我继续说），而相信受难之\Person{基督}的万民将要期待祂未来的显现。为此，\Person{圣神}以比喻并隐晦地宣说了这些真理：因为经上记载：\BibleQuote{犹大，你的兄弟们赞美了你；你的手在你仇敌的颈上；你父亲的儿子们要向你下拜。犹大是只幼狮；我儿，你从萌芽中长出。他卧下如狮，又如幼狮：谁能唤醒他？权杖不离开犹大，领导者的棍杖不离开他两腿之间，直到那为他所存留的来到；他将成为万民的期望，把他的驴驹拴在葡萄树上，把他母驴的驹子拴在葡萄树的嫩枝上。他在酒中洗他的衣服，在葡萄的血中洗他的袍子。他的眼睛因酒而发红，他的牙齿因奶而洁白。}\BibleRef{创 49:8}此外，从起初直到这位\Person{耶稣基督}出现并受苦的时候，你们民族中从未缺少先知或统治者，你们不敢厚颜无耻地断言，也无法证明。因为虽然你们断言，祂在治下受苦的\Person{黑落德}是一个阿市刻隆人，但你们承认你们民族中有一位大司祭；如此你们当时有一个人按照\Person{梅瑟}法律奉献祭品，并遵守其他的法律仪式；也有相继的先知直到\Person{若翰}，（甚至当你们的民族被掳到\Place{巴比伦}，你们的土地被战争蹂躏，圣器被夺走时）；你们中间从未缺少一位先知，他作你们民族的主人、领袖和统治者。因为在先知内的\Person{圣神}傅油立了你们的君王，并坚立了他们。但在我们的\Person{耶稣基督}在你们民族中显现并受死之后，任何地方都没有先知了：不仅如此，你们不再拥有自己的君王，你们的土地荒废，如同葡萄园中的茅舍被遗弃；而圣经借\Person{雅各伯}的口所说的话：\BibleQuote{祂将成为万民的期望}，象征性地指祂的两次来临，以及万民将相信祂；这些事实你们如今终可辨别。因为来自万民中，借着\Person{基督}的信德而虔诚正义的人，期待祂未来的显现。\par

% structure: p0026 source=h2 target=subsection reason=relative-heading-rank; base=h2

\subsection{第五十三章 \Person{雅各伯}预言\Person{基督}将骑驴，\Person{匝加利亚}证实此事}

\Emph{\Person{尤斯定}：}那句表述：\BibleQuote{把他的驴驹拴在葡萄树上，把他母驴的驹子拴在葡萄树的嫩枝上}，既预告了祂在首次来临时所行的事工，也预告了万民将寄予祂的信心。因为他们像未上轭的驴驹，颈上未曾负轭，直到这位\Person{基督}来到，派遣祂的门徒教导他们；他们负起祂言语的轭，屈服颈项忍受一切苦难，为了祂自己所应许、他们期盼的美善事物。确实，我们的主\Person{耶稣基督}，当祂打算进入\Place{耶路撒冷}时，曾要求祂的门徒给祂牵来一匹驴和一头驴驹，它们被拴在一个叫\Place{贝特法革}的村庄入口；祂骑上它，进入了\Place{耶路撒冷}。由于祂所做的正如预言所精确描述\Person{基督}将行的事，并且应验得到了承认，这便成了祂是\Person{基督}的明证。虽然这一切都发生了，并且从圣经得到证实，你们仍然心硬。不仅如此，十二位先知中的\Person{匝加利亚}也曾预言此事将发生，说了以下的话：\BibleQuote{熙雍女子，你应尽情欢乐！耶路撒冷女子，你应欢呼，看，你的君王到你这里来，他是正义的，带来救恩，谦逊、卑微，骑在驴上，骑在驴驹上。}\BibleRef{匝 9:9}如今，预言之神与圣祖\Person{雅各伯}都提到了一匹驴和它的驴驹，它们将为祂所用；而且，正如我先前所说，祂要求门徒牵来这两匹牲口；这一事实预言了你们会堂中的人连同外邦人要相信祂。因为正如未上轭的驴驹是外邦人的象征，同样上了轭的驴是你们民族的象征。因为你们拥有先知所施加于你们的法律。此外，先知\Person{匝加利亚}预言这位\Person{基督}将被击打，祂的门徒将四散：这事也发生了。因为在祂被钉十字架后，跟随祂的门徒四散了，直到祂从死者中复活，并说服他们相信关于祂的预言如此说，祂必须受苦；于是他们被说服，走遍普世，教导这些真理。因此我们在祂的信德和教义上也坚固，因为我们既有来自先知的劝服，也有来自普世看见的以那位被钉者的名敬拜\Person{天主}之人的见证。\Person{匝加利亚}也说了以下的话：\BibleQuote{刀剑，起来攻击我的牧人，攻击我百姓中的人——万军的上主说。打击牧人，羊群就四散。}\BibleRef{匝 13:7}\par

% structure: p0028 source=h2 target=subsection reason=relative-heading-rank; base=h2

\subsection{第五十四章 葡萄的血所象征的意义}

\Emph{\Person{尤斯定}：}那段由\Person{梅瑟}写下、由圣祖\Person{雅各伯}预言的话，即：\BibleQuote{他在酒中洗他的衣服，在葡萄的血中洗他的袍子}，象征祂要用自己的血洗净那些相信祂的人。因为\Person{圣神}称那些借着祂获得罪赦的人为祂的衣服；祂在他们中间常以权能临在，但在祂第二次来临时将明显地临在。圣经提及葡萄的血显然是经过设计的，因为\Person{基督}的血不是来自人的种子，而是来自\Person{天主}的德能。因为正如是\Person{天主}而非人产生了葡萄的血，同样圣经也预言\Person{基督}的血不是来自人的种子，而是来自\Person{天主}的德能。但是各位先生，我重复的这个预言证明\Person{基督}不是由人所生的人，不是按人类常例所生的。\par

\SourceRecord{Translated by Marcus Dods and George Reith.；Ante-Nicene Fathers；Vol. 1.；Edited by Alexander Roberts, James Donaldson, and A. Cleveland Coxe.；Buffalo, NY: Christian Literature Publishing Co.,；1885.；<http://www.newadvent.org/fathers/01284.htm>.}

% structure: p0001 source=h1 target=section reason=page-title

\section{与特里福的对话（第55-68章）}

% structure: p0002 source=h2 target=subsection reason=relative-heading-rank; base=h2

\subsection{第55章 特里福要求证明基督是天主，但不用比喻。尤斯定承诺这样做}

\Emph{特里福：} 我们当记住你这番阐述，如果你用其他论据加强[解决]这个难题；但现在继续讨论，向我们证明预言之神承认在万有的创造者之外还有另一位天主，注意不可提到日月，因为经上记载，天主将这些赐给外邦人作为神明崇拜；先知们时常使用这种说话方式，说：\Quote{你们的天主是万神之神，万主之主}，常常加上\Quote{伟大、强有力、可畏的[天主]}。因为这样的表达并非因为他们真的是神，而是圣经教导我们：真天主，即创造万有的那一位，是独一的主，超越那些被称为神和主的。为使圣神说服我们，祂通过圣达味说：\Quote{外邦人的神，所谓的神，是魔鬼的偶像，并非神；}祂诅咒那些敬拜它们的人。\par

\Emph{尤斯定：} 我不愿提出这些证据，特里福，我知道那些敬拜这些[偶像]和类似之物的人已被定罪，但要提出无人能反驳的[证据]。它们对你而言似乎陌生，尽管你每天读它们；因此就连这件事我们也明白，由于你们的邪恶，天主已阻止你们辨识祂圣经的智慧；但仍有一些例外，按照祂容忍的恩宠，如依撒意亚所说，祂为你们留下救恩的种子，免得你们的种族像索多玛和哈摩辣那样完全毁灭。因此，请留意我将从圣经中记录的话，这些不需要解释，只需聆听。\par

% structure: p0005 source=h2 target=subsection reason=relative-heading-rank; base=h2

\subsection{第56章 向梅瑟显现的天主与天主圣父有别}

\Emph{尤斯定：} 那么，梅瑟，天主蒙福而忠信的仆人，宣告在玛默勒橡树下向亚巴郎显现的那一位是天主，祂带着两位天使同行，受另一位永远住在超天界、对所有人不可见、不与任何人亲自交往的天主派遣，去审判索多玛；我们相信这位是万有的创造者和圣父；因为梅瑟这样说：\Quote{天主在玛默勒橡树下向他显现，正当中午他坐在帐幕门口时。他举目观看，见有三人站在他面前；他一见，就从帐幕门口跑去迎接他们，俯伏在地，说……} \BibleRef{创 18:1} \Quote{亚巴郎清晨早起，来到他从前站在主面前的地方，向索多玛和哈摩辣以及附近的地方观望，看见地上有火焰上升，如同熔炉的烟。} \BibleRef{创 19:27}\par

我引用完这些话后，问他们是否理解了。他们说理解了，但所引的段落不能证明除了万有的创造者之外另有天主或主，或圣神如此说。\par

\Emph{尤斯定：} 既然你们理解了圣经，我要试图说服你们[相信]我所说的真理：有一位另有的天主和主，服从于万有的创造者；祂也被称为天使，因为祂向人宣布万有的创造者——其上再没有别的天主——愿意向他们宣布的一切。\par

我再次引用了先前的段落。\par

\Emph{尤斯定：} 你们认为天主在玛默勒橡树下向亚巴郎显现，如圣经所断言的吗？\par

\Emph{特里福：} 当然。\par

\Emph{尤斯定：} 祂是亚巴郎所见三人中的一位吗？即预言之神描述为人的那三位？\par

\Emph{特里福：} 不；天主在看见那三人之前已向他显现。那三位圣经称为人的是天使；两位被派去毁灭索多玛，一位向撒辣宣布喜讯，说她要生一个儿子；为此他被派出，完成任务后便离开了。\par

\Emph{尤斯定：} 那么，那三人中的一位，曾在帐幕中，说：\Quote{明年此时我要回到你这里，撒辣必生一个儿子} \BibleRef{创 18:10}，当撒辣生了儿子后，祂似乎回来了，并在此处被预言的话称为天主。为使你们清楚明白我所说的，请听梅瑟特意使用的言辞，它们是：\Quote{撒辣看见埃及婢女哈加尔给亚巴郎生的儿子与她的儿子依撒格嬉戏，就对亚巴郎说：你把此婢女和她的儿子赶走，因为这婢女的儿子不能与我的儿子依撒格一同承受产业。亚巴郎因他儿子的缘故很苦恼。但天主对亚巴郎说：你不要因这孩子和你的婢女而苦恼；凡撒辣对你说的，你都应听从，因为只有从依撒格生的，才称为你的后裔。} \BibleRef{创 21:9} 那么，你们是否看见了，那位在橡树下说要回来的，因为知道必须劝告亚巴郎做撒辣所愿的事，正如经上记载祂回来了，并且是经文中说话的天主：\Quote{天主对亚巴郎说：你不要因这孩子和你的婢女而苦恼？}\par

\Emph{特里福：} 当然；但你未能由此证明，除了向亚巴郎显现、也向其他圣祖和先知显现的那一位之外，另有天主。然而你证明了我们以为在亚巴郎帐幕中的三人全是天使是错误的。\par

\Emph{尤斯定：} 如果我不能从圣经向你们证明那三人中的一位是天主，并被称为天使，因为如前所说，祂向那些万有的创造者天主愿意[传达信息]的人传达信息，那么关于那位在世上以人形向亚巴郎显现、如同与祂同来的两位天使一样，并且在创造世界之前就已存在的那位，你们接受与你们全民族相同的看法是合理的。\par

\Emph{特里福：} 当然，直到此刻我们一直这样相信。\par

\Emph{犹斯定：}现在回到圣经，我要努力说服你们：那位曾显现给亚巴辣罕、雅各伯和梅瑟，并被称为天主的，与创造万有的那位——在数目上，而非在旨意上——是不同的。因为我断言，祂从未做过任何造物主（其上没有别的神）不愿祂做或不愿祂参与的事。\par

\Emph{特黎丰：}现在就证明这是真的，好让我们也同意你。因为我们并不认为你主张祂做过或说过任何与造物主意愿相悖的事。\par

\Emph{犹斯定：}我刚引用的圣经会向你说明这一点。它这样说：\Quote{太阳已升在地上，罗特进了左哈尔；上主从天上，从自己那里，将硫磺与火降在索多玛，毁灭了这些城和整个平原。}\BibleRef{创 19:23}\par

\Emph{那留在特黎丰身边的第四个人：}因此必须说，那两位去索多玛的天使中，有一位被梅瑟在圣经中称为\Quote{主}的，与那位也是天主且显现给亚巴辣罕的，是不同的。\par

\Emph{犹斯定：}不单基于此理由就必须绝对承认，圣神将另一位称为\Quote{主}，而非那被视为造物主的；不单基于梅瑟的话，也基于达味的话。因为他记着说：\Quote{上主对我主说：你坐在我右边，等我使你的仇敌作你的脚凳。}\BibleRef{咏 110:1} 又在另一处说：\Quote{天主，你的宝座永远常存；你王权的权杖是正直的权杖。你喜爱正义，憎恶邪恶；因此天主、你的天主，用喜乐的油傅了你，胜过你的同伴。}\BibleRef{咏 45:6-7} 因此，如果你们主张，除了万有的圣父和祂的基督之外，圣神还称另一位为天主和主，那就回答我；因为我承担从圣经本身向你们证明，圣经称为\Quote{主}的，并非去索多玛的两位天使之一，而是与他们同在、被称为天主的那位，就是显现给亚巴辣罕的。\par

\Emph{特黎丰：}证明这一点；因为如你所见，天色已晚，我们并未准备好作这样危险的回答；因为我们从未听过任何人探究、查考或证明这些事；若非你将一切引证到圣经，我们绝不会容忍你的谈话；因为你非常热切地从圣经中引证，而且你认为在造物主之上没有别的天主。\par

\Emph{犹斯定：}那么你们知道，圣经说：\Quote{上主对亚巴辣罕说：\InnerQuote{撒辣为什么笑，说：我果真能生育吗？我已经老了。难道天主有难成的事吗？到了定期，我要回到你这里，按这生命的时候，撒辣必生一个儿子。}}\BibleRef{创 18:13} 过了一会儿：\Quote{那些人就从那里起身，面向索多玛观望；亚巴辣罕与他们同行，要送他们一程。上主说：\InnerQuote{我岂能瞒着亚巴辣罕，我所要作的事呢？}}\BibleRef{创 18:16} 又过了一会儿，这样说：\Quote{上主说：\InnerQuote{索多玛和哈摩辣的喊声甚大，他们的罪极其严重。我要下去，看看他们所行的，是否完全像那达到我这里的喊声一样；若不然，我也要知道。}那些人就转身离开那里，往索多玛去；但亚巴辣罕仍旧站在上主面前。亚巴辣罕近前来说：\InnerQuote{你要将义人与恶人一同毁灭吗？}}\BibleRef{创 18:20}\par

如此等等，因为我不认为需要重写同样的词句（之前都已写过），但必须给出那些我用来向特黎丰及其同伴建立证明的段落。于是我继续往下，记有这些话：\par

\Emph{犹斯定：}\Quote{上主与亚巴辣罕说完了话就走了；亚巴辣罕也回到自己的地方去了。黄昏时，两位天使来到索多玛；罗特正坐在索多玛城门口。}\BibleRef{创 18:33} 以及直到\Quote{那两个人伸出手来，将罗特拉进屋里，关上门。}\BibleRef{创 19:10} 接着直到\par

\Quote{天使拉着他的手，和他妻子的手，并他两个女儿的手，因为上主怜悯他。领他们出来以后，就说：\InnerQuote{逃命吧！不可回头看，也不可在平原任何地方站住；要逃到山上，免得你被剿灭。}罗特对他们说：\InnerQuote{我主啊，不要这样！你的仆人已经在你眼前蒙恩宠，你又向我显大正义，救我的性命；我不能逃到山上，恐怕灾祸追上我，我就死了。看哪，这座城又小又近，容易逃到；那不是一座小城吗？求你容我逃到那里，我的性命就得存活。}天使对他说：\InnerQuote{这事我也应允你，不倾覆你所说的这城。你要速速地逃到那里；因为你还没有到那里，我不能作什么。}因此那城名叫左哈尔。太阳已升起在地上；罗特进了左哈尔。上主从天上，从自己那里，将硫磺与火降在索多玛和哈摩辣，毁灭了这些城和整个平原。}\BibleRef{创 19:16}\par

\Person{Justin}: \Emph{(又停顿片刻)} 朋友们，难道你们还没有觉察到，那三者之一，既是天主又是主，并服务于那在天上的那位，乃是两位天使的主吗？因为当[天使]前往索多玛时，祂留下来，用梅瑟记录的话语与亚巴郎交谈；当祂谈完离去后，亚巴郎回到了自己的地方。当祂[来到索多玛]时，两位天使不再与罗特交谈，而是祂亲自与他交谈，正如圣经所明示的；祂就是那位从在天上的主（即万有的创造者）领受使命，以降罚于索多玛和哈摩辣的主，圣经以这样的话描述说：\BibleQuote{上主从天上从主那里将硫磺与火降于索多玛和哈摩辣。}\par

% structure: p0029 source=h2 target=subsection reason=relative-heading-rank; base=h2

\subsection{第五十七章 犹太人反对：如果祂是天主，为何说祂吃了？尤斯定的回答}

\Person{Trypho}: \Emph{（当我沉默时）} 显然，圣经迫使我们必须承认这一点；但有一点我们理所当然感到困惑——即关于[主]吃了亚巴郎所准备并摆在他面前的东西的记载；而你会承认这一点。\par

\Person{Justin}: 经上记载他们吃了；如果我们相信所记载的是那三位吃了，而不仅仅是两位——他们实际上是天使，在天上受滋养，正如我们所明见的，尽管他们不是靠与凡人相同的食物来滋养（因为关于在旷野中供养你们祖先的玛纳，圣经这样说道，他们吃了天使的食物）——[如果相信三位都吃了]，那么我要说，记载他们吃了的圣经，其意义就如同我们说火吞灭了万物一样；然而可以肯定，他们并不是用牙齿和下颌咀嚼而吃。所以，即使是这一点，我们也不应丝毫感到困惑，只要稍微熟悉比喻的表达方式，并能超越它们。\par

\Person{Trypho}: 关于吃的方式，或许可以这样解释：即经上记载他们拿了并吃了亚巴郎所准备的东西；这样，你现在可以向我们解释，这位向亚巴郎显现、并服务于创造万有的天主、由童贞女所生、成为人、具有与我们相同的情感，正如你先前所说的。\par

\Person{Justin}: 请允许我先收集一些其他关于此点的证据，好使你因数量众多而被说服[接受真理]，之后我将解释你所问的。\par

\Person{Trypho}: 你看着办吧；我会非常满意的。\par

% structure: p0035 source=h2 target=subsection reason=relative-heading-rank; base=h2

\subsection{第五十八章 从向雅各伯显现的异象中证明同一件事}

\Person{Justin}: 我打算向你们引用圣经，并非仅仅为了技艺性的言辞展示；因为我没有这种能力，唯独天主的恩宠赐予我理解祂圣经的能力，我劝勉众人白白地、慷慨地分享这恩宠，以免因缺乏它而在万物创造者天主将要通过我主耶稣基督进行的审判中遭受定罪。\par

\Person{Trypho}: 你所做的配得上对天主的敬拜；但你说你没有技艺性言辞的储备，在我看来你是在假装无知。\par

\Person{Justin}: 既然你这样认为，那就这样吧；但我确信我说的是实话。现在请听我继续引出其余的证明。\par

\Person{Trypho}: 请讲。\par

\Person{Justin}: 我的兄弟们，梅瑟又记载说，那位被称为天主并向先祖显现的，同时被称为天使和主，好使你们由此明白，祂是万物之父的仆役，正如你们已经承认的，并因更多的论据而坚定确信。因此，天主的圣言通过梅瑟论到亚巴郎的孙子雅各伯时这样说：\par

\BibleQuote{羊群配合的时候，我在梦中举目一看，见跳在羊群上的公山羊和公绵羊都是有白纹的、有斑点的、有杂色的。天主的使者在梦中对我说：雅各伯！雅各伯！我答说：主，有什么事？祂说：你举目一看，跳在羊群上的公山羊和公绵羊都是有白纹的、有斑点的、有杂色的，因为我看见了拉班对你所作的一切。我是那在贝特耳向你显现的天主，你在那里曾立过石柱，并向祂许过愿。现在起身，离开这地方，回到你的出生地去吧，我必与你同在。}\BibleRef{创 31:10}\par

又用别的话论到同一个雅各伯，这样说：\par

\BibleQuote{他当夜起来，带了他的两个妻子，两个婢女和十一个孩子，过了雅波克河。他带着他们过了河，又送过了他所有的一切。雅各伯独自一人留在后面，有一个人与他角力直到破晓。那人见自己不能胜他，就摸了雅各伯的大腿窝；雅各伯与他角力时，大腿窝就脱了节。那人说：让我走，因为天已破晓。雅各伯说：你不祝福我，我就不让你走。那人问他说：你叫什么名字？他答说：雅各伯。那人说：你的名字不再叫雅各伯，要叫以色列，因为你与天主与人角力都得胜了。雅各伯问他说：请告诉我你的名字。那人答说：为什么问我的名字？于是在那里祝福了他。雅各伯给那地方起名叫培尼耳，因为我面对面见了天主，我的灵魂得了安慰。}\BibleRef{创 32:22}\par

又用别的话论到同一个雅各伯，这样说：\par

\BibleQuote{雅各伯来到客纳罕地的路次，即贝特耳，他和所有跟随他的人。他在那里筑了一座祭坛，称那地方为贝特耳，因为天主在他逃避他哥哥厄撒乌的时候曾在那里显现给他。黎贝加的乳母德波辣死了，葬在贝特耳下面的一棵橡树旁；雅各伯称它为哀哭的橡树。天主再次在路次显现给雅各伯，当他从叙利亚的美索不达米亚出来时，他祝福了他。天主对他说：你的名字不再叫雅各伯，但以色列将是你的名字。\BibleRef{创 35:6}}\par

祂被称为天主，祂是天主，且永远是天主。\par

当大家都同意这些理由后，我继续说道：\par

\Quote{雅各伯从誓约井出发，走向哈兰。他来到一个地方，便在那里过夜，因为太阳已经下山；他从那地方捡了一块石头，放在头下。他在那地方睡了，梦见地上竖立一架梯子，梯顶直达天上；天主的使者在梯子上上去下来。主站在梯子上，说：我是主，你祖父亚巴郎的天主，依撒格的天主；不要害怕：你所躺卧之地，我必赐给你和你的后裔；你的后裔必像地上的尘土，向西、南、北、东扩展；地上万民都要因你和你的后裔蒙受祝福。看，我与你同在，在你所行的路上保护你，并必领你回到这地方；因为我必不离开你，直到我实现了对你所说的一切。雅各伯从睡梦中醒来，说：主确实在这地方，我竟不知道。他就害怕，说：这地方何等可畏！这不是别的，乃是天主的住所，是天上的门。雅各伯清早起来，把放在头下的石头立作柱子，在柱顶浇上油；雅各伯称那地方为天主的住所，那城原名乌路斯。}\par

% structure: p0049 source=h2 target=subsection reason=relative-heading-rank; base=h2

\subsection{第五十九章 与圣父不同的天主与梅瑟交谈}

\Person{儒斯定}：请允许我进一步从\BibleBook{出谷}{纪}向你们指明，这同一位既是天使、天主、主、人，曾以人形显现给亚巴郎和依撒格，现在又从荆棘丛中的火焰里显现，并与梅瑟交谈。\par

他们回答说愿意愉快、耐心且热切地聆听后，我继续说道：\par

这些话记载在名为\BibleBook{出谷}{纪}的书中：\BibleQuote{过了许多日子，埃及王死了，以色列子民因劳役而叹息，}\BibleRef{出 2:23}等等，直到\BibleQuote{主，你们祖先的天主、亚巴郎的天主、依撒格的天主、雅各伯的天主显现给我说：我确实眷顾了你们，也眷顾了你们在埃及所受的苦难。}\BibleRef{出 3:16}\par

除了这些话，我继续说道：\par

诸位，你们是否察觉，梅瑟所称之为天使、在火焰中与他谈话的这位天主，向梅瑟宣称祂是亚巴郎、依撒格和雅各伯的天主？\par

% structure: p0055 source=h2 target=subsection reason=relative-heading-rank; base=h2

\subsection{第六十章 犹太人对那在荆棘丛中显现者的看法}

\Person{特黎风}：我们从你所引的段落中没有看出这一点，而只看出：在火焰中显现的是天使，但与梅瑟交谈的是天主；所以在那异象中实际上有两个人彼此相随，一个天使和天主。\par

\Person{儒斯定}：即使如此，我的朋友，在梅瑟所见异象中天使和天主在一起，然而，正如先前所引段落已向你们证明的，那对梅瑟说祂是亚巴郎的天主、依撒格的天主、雅各伯的天主的天主，并不是万有的创造者，而是那位已向你们证明曾显现给亚巴郎的那一位，祂服从事奉万有造物者的旨意，并同样在审判索多玛时执行祂的计谋；因此，即使如你们所说，有两个——一个天使和天主——，即使智力最低的人也不敢断言，万有之父和造物者，舍弃了一切天上之事，竟在人间一小块土地上显现可见。\par

\Person{特黎风}：既然先前已经证明，那位被称为天主和主、曾显现给亚巴郎的，从在天上的主那里领受了祂加给索多玛之地的惩罚，即使有一个天使伴随着那显现给梅瑟的天主，我们也会明白，那从荆棘丛中与梅瑟交谈的天主并不是万有的造物者，而是那位已经显示曾显现给亚巴郎、依撒格和雅各伯的；祂也被称为并被视为万有造物者天主的天使，因为祂向人宣布了父及万有造物者的命令。\par

\Person{犹斯定}：现在，\Person{特黎风}，我确实要表明，在梅瑟的异象中，正是那被称为\Quote{天使}而又被称为\Quote{天主}的独一者，显现给梅瑟并与他交谈。因为圣经如此说：\BibleQuote{上主的天使从荆棘中的火焰显现给他；他看见荆棘被火烧着，却没有烧毁。梅瑟说：我要过去看看这大异象，为什么荆棘没有烧毁。上主见他过去观看，便从荆棘中叫他。}\BibleRef{出 3:2} 同样，圣经称那在梦中显现给雅各伯的天使为\Quote{天使}，随后又说那在梦中显现的天使向他说：\BibleQuote{我是那在你逃避你哥哥厄撒乌时向你显现的天主。}\BibleRef{创 35:7} 并且，在亚巴郎时代索多玛所受的审判中，圣经说：主（那住在天上的主）施加了惩罚。同样地，圣经在此宣布上主的天使显现给梅瑟，随后又宣称他是主和天主，指的是同一位。通过许多已引用的证言，圣经宣称他是侍奉那超越世界、其上没有其他的[天主]的使者。\par

% structure: p0060 source=h2 target=subsection reason=relative-heading-rank; base=h2

\subsection{第六十一章 智慧由圣父所生，如火生火}

\Person{犹斯定}：朋友们，我要从圣经给你们另一个证据：天主在一切受造物之前生了一个元始，[祂是]一个出自祂的理性能力[由祂所生]。这能力被圣神有时称为主的荣耀，有时称为圣子，有时称为智慧，有时称为天使，然后称为天主，然后又称为主和\OriginalTerm{圣言}{Logos}；在另一场合，当祂以人形显现给农的儿子若苏厄时，祂自称元帅。因为祂可以被称为所有这些名字，因为祂奉行圣父的旨意，并因圣父以意志行动生了祂。正如我们中间发生的事：当我们发出一个言语时，我们就生了这个言语；但并非通过切割而减少我们内[存留]的言语。也如同我们在火中所见的：当火点燃[另一团火]时，它并不减少，而是保持相同；被它点燃的火同样似乎自立存在，并不减少那点燃它的火。智慧之圣言——祂本身是万有之父所生的天主、圣言、智慧、德能和生者的荣耀——当祂借撒罗满说以下的话时，将为我作证：\par

\Quote{如果我每日向你们宣告将要发生的事，我将回忆从永远的事，并加以回顾。主造了我，作为祂道路的开始，为祂的作为。从永远，在起初，祂就奠定了我，在祂造地之前，在祂造深渊之前，在泉水涌出之前，在群山奠定之前。在一切丘陵之前，祂生了我。天主造了原野、旷野和天下最高的居所。当祂预备诸天时，我与祂同在；当祂在风中安置祂的宝座时，当祂使高云坚固、深渊的泉源安全时，当祂奠定大地的根基时，我与祂一同安排。我是祂所喜悦的；每日每时，我因祂的容颜而喜乐，因为祂喜悦人居之世的完成，并喜悦人子。因此，我儿，现在要听从我。听从我的人是有福的，那持守我道路的凡人，每日在我门口守望，在我入口的门柱旁观察。因为我的出口是生命之出口，[我的]旨意已由主预备。但那些得罪我的，伤害自己的灵魂；那些恨我的，喜爱死亡。}\par

% structure: p0063 source=h2 target=subsection reason=relative-heading-rank; base=h2

\subsection{第六十二章 \Quote{让我们造人}一语与\BibleBook{}{箴言}的见证相符}

\Emph{儒斯定：}诸位朋友，梅瑟所记载的天主的话，在向我们指出祂所指示的那一位时，也表明天主在造人时怀着同样的意图，用以下的话说：\BibleQuote{让我们按我们的肖像和模样造人，让他们管理海里的鱼、天空的飞鸟、牲畜、全地以及地上所有爬行的生物。天主于是造了人：按天主的肖像造了他；造了他们男和女。天主祝福了他们，说：你们要生育繁殖，充满大地，治理大地。}为了使你们不改变刚才所引用的话的意思，也不重复你们老师的主张——要么说天主对自己说\Quote{让我们造}，就像我们准备做某事时常对自己说\Quote{让我们做}；要么说天主对元素（即地和其他类似物质，我们认为人是由这些元素形成的）说\Quote{让我们造}——我要再次引用梅瑟自己所叙述的话，从这些话我们可以明确无误地知道（天主）与一位在数目上与祂不同的、也是理性的存有在交谈。这些话是：\BibleQuote{天主说：看，亚当已成为我们中的一个，知道善恶。}\BibleRef{创 3:22}因此，当（梅瑟）说\BibleQuote{如同我们中的一个}时，他宣告了彼此关联的若干位格，且他们至少是两个。因为我不愿说你们中间流传的那个异端教义是真实的，也不愿说其教师能证明（天主）对天使说话，或者人的身体是天使的作品。但这个真正由圣父所生的子，在一切受造物之前就与圣父同在，圣父与祂交谈；正如圣经通过撒罗满所表明的，撒罗满称之为智慧的那一位，是作为万物的开始和由天主所生的子，而天主也在农的儿子若苏厄的启示中宣告了同样的事。因此，请听若苏厄书中的以下内容，使我的话显明给你们：\BibleQuote{当若苏厄靠近耶里哥时，他举目观看，看见一个人站在他对面。若苏厄走近他，说：你是帮我们，还是帮我们的敌人？他对若苏厄说：我是上主万军的统帅，现在我来。若苏厄就俯伏在地，说：我主，你要你仆人做什么？上主的统帅对若苏厄说：将你脚上的鞋脱下，因为你站的地方是圣地。耶里哥城门紧闭，无人出入。上主对若苏厄说：看，我把耶里哥和它的君王以及勇士都交在你手中。}\par

% structure: p0065 source=h2 target=subsection reason=relative-heading-rank; base=h2

\subsection{证明这位天主曾降生成人}

\Person{特里丰}：我的朋友，这一点已经通过许多论证有力地向我证明了。那么，还需要证明祂按照圣父的旨意，由童贞女降生成人，并被钉十字架而死。同时也请清楚地证明，祂此后复活并升天了。\par

\Emph{儒斯定：}朋友们，这一点在我之前引用过的预言话语中也已经表明；为了你们，我将重新回忆并解释这些话语，努力使你们在这件事上也同意我。那么，依撒意亚所记载的这段话：\BibleQuote{谁能述说祂的后裔？因为祂的生命从地上被夺去，}\BibleRef{依 53:8}——它岂不指向一位不是由人所生、而是因百姓的过犯由天主交付死亡的那一位吗？——至于祂的血，梅瑟（如我先前所说）在比喻中说，祂要在葡萄的血中洗自己的衣裳；因为祂的血不是从人的种子而来，而是从天主的旨意而来。此外，达味所说的话：\BibleQuote{在祢圣洁的光辉中，我从母腹中生祢，在晨星之前。上主起了誓，决不反悔：祢要按照默基瑟德的品位，永为司祭，}——这岂不向你们宣告祂从太初就已存在，而万有的天主圣父意愿祂由人的母腹所生吗？还有在其他已经引用过的话中（他说）：\BibleQuote{祢的御座，天主，永远常存；祢王国的权杖是正直的权杖。祢喜爱正义，憎恨不义；因此天主，祢的天主，用喜乐油傅了祢，胜过祢的同伴。祂以没药、香料和肉桂从祢的衣袍，从象牙宫中，使祢喜乐。君王的女儿都来尊敬祢。王后佩戴着金饰站在祢右边。女儿，请听，请看，请侧耳倾听，忘记你的民族和你的父家；君王将恋慕你的美丽：因为他是你的主，你应朝拜祂。}因此，这些话明确作证：那建立这一切的一位，见证祂堪受朝拜，作为天主和基督。此外，天主的话对那些相信祂的人说话，视他们为一个灵魂、一个会堂、一个教会，如同对女儿说话；它同样以以下话语明确宣告，这些话也教导我们忘记旧日的祖先习俗，它们这样说道：\BibleQuote{女儿，请听，请看，请侧耳倾听；忘记你的民族和你的父家；君王将恋慕你的美丽：因为他是你的主，你应朝拜祂。}\par

% structure: p0068 source=h2 target=subsection reason=relative-heading-rank; base=h2

\subsection{\Person{儒斯定}向否认需要这位基督的犹太人提出其他证据}

\Person{特里丰}：让祂被你们外邦人承认为主、基督和天主，正如圣经所宣告的；但我们这些事奉那位创造了这位（基督）的天主的人，不需要宣认或朝拜祂。\par

\Signature{由斯定} 若我像你一样好争辩而轻率，特黎丰，我便不会再与你继续谈论，因为你已准备好不理会所说的话，只想找些刁钻的答复；但现在，因我畏惧天主的审判，对于你们民族中是否有人能藉万军上主的恩宠得救，我并不冒然定论。因此，即使你行事不当，我仍会继续答复你提出的任何命题和任何反驳；事实上，我对各民族的每一个人也都如此，凡愿意与我一同探讨或向我询问此事者，我都同样对待。所以，你若留意我先前引用的圣经，便已明白，你们民族中得救的人乃是藉此（人）得救，并分享祂的份；你定然不会问我此事。我将再次重述先前引用的达味的话，并恳请你理解，不要行事不当，彼此怂恿只作反驳。达味的话如下：\BibleQuote{上主为王了；愿万民发怒：祂坐于革鲁宾之上；愿大地震荡。上主在熙雍为大，祂超乎万民之上。愿他们称谢你伟大的名，因它可畏而神圣；君王的尊荣喜爱审断。你预备了正直；你在雅各伯中施行了审断和公义。你们要尊崇上主我们的天主，敬拜祂脚下的脚凳；因为祂是圣的。梅瑟和亚郎在祂司祭之中，撒慕尔在呼求祂名的人之中；他们呼求上主，祂就应允他们。在云柱中祂向他们说话；因为他们遵守了祂的法度和祂所赐的诫命。}另从先前引用的达味的话，即你们愚昧地以为指撒罗满而说的话，因题为\Quote{为撒罗满作}，也可以证明它们并非指撒罗满，而是指此（基督）在太阳之前已存在，你们民族中得救的人都要藉祂得救。这些话是：\BibleQuote{天主，请将你的审断赐给君王，将你的公义赐给君王的儿子。祂将以公义审断你的百姓，以审断你贫穷的民众。大山将给百姓带来和平，小山带来公义。祂要审断百姓中的穷人，拯救穷苦之子，压制诽谤者：祂要与太阳共存，在月亮之前，直到万代。}等等，直到\BibleQuote{祂的名在太阳之前永存，地上万民都要因祂蒙福。万民都要称祂有福。上主以色列的天主当受称颂，惟有祂行奇妙的事；祂荣耀的名当受称颂，直到永远。愿全地充满祂的荣耀。阿们，阿们。}你们也记得达味所说的其他话，就是我先前提到过的，如何宣告祂要从至高的天上降下，又要回到原处，叫你们认出祂是从上降下的天主，又是生活在人间的真人；并宣告祂还要显现，那些刺透祂的人要看见祂，并为祂哀悼。这些话是：\BibleQuote{诸天述说天主的荣耀，穹苍宣扬祂手的化工。日复一日发出言语，夜复一夜显示知识；他们没有言语，没有词句，也听不到他们的声音。他们的声音传遍普世，他们的言语传到地极。祂在太阳中安设了帐幕；祂如同新郎走出洞房，又如勇士欢然奔路：祂从最高的天上出发，又回到最高的天上，没有一物能隐藏不被祂的热力所照。}\par

% structure: p0071 source=h2 target=subsection reason=relative-heading-rank; base=h2

\subsection{Chapter 65. 犹太人反对说天主不将光荣给予他人。由斯定解释这段经文。}

\Signature{特黎丰} 因这许多圣经而动摇，我不知该如何解释依撒意亚所写的经文，其中天主说祂不将光荣给予他人，如此说：\BibleQuote{我是上主天主；这是我的名；我决不将我的光荣给予另一个，也不将我的德能给予另一个。}\BibleRef{依 42:8}\par

\Signature{由斯定} 特黎丰，你若说这些话时，心地单纯、毫无恶意，既不重复上文也不补充下文，那还可原谅；但若你因以为能对此经文置疑，好叫我说圣经自相矛盾，那你就错了。但我决不敢作此猜想或如此说；若真有看似这样的经文被提出，且有人以此为借口说它与（别的）矛盾，我既深信没有一处圣经与另一处相矛盾，宁可承认我不理解所记的内容，并努力说服那些以为圣经自相矛盾的人，宁可与我同想。你提出难题的用意如何，天主知道。但我要提醒你该经文说些什么，好叫你从这（经文）本身也看出，天主只将光荣给予祂的基督。先生们，我要取一些短小的段落，即与特黎丰所说相连的，并前后连贯的章节。因我不重述其他部分的，而只取那些连贯在一起的。你们也请留心听。这些话是：\par

\BibleQuote{主天主，那创造诸天、铺张穹苍、奠立大地和地上万物、赐气息给地上的百姓、赐神灵给行在地上者的这样说：我上主天主，因公义召叫了你，我必握住你的手，保护你，立你作人民的盟约，作外邦人的光明，为开启盲人的眼，从牢狱中领出被囚的人，从监牢中领出坐在黑暗中的人。我是上主天主，这是我的名；我决不将我的光荣让给另一位，也不将我的荣誉让给雕像。看，先前的事已经成就；新的事我现在宣布，在它们发生之前，我已告诉给你们。你们应向上主高唱新歌，他的主权从地极开始。你们这些下海的人，以及在海上常航行的人，海岛上和海上的居民，都应歌唱。旷野及其中的村庄，以及刻达尔居民的帐幕，都应欢呼；住在磐石上的居民，应高声喊叫，从山顶上呼喊：他们应将光荣归于天主，在岛屿上传扬他的荣誉。万军的上主天主必出征，他要彻底消灭战争，激发热诚，向敌人勇敢呐喊。\BibleRef{依 42:5}}\par

当我重复了这段话之后，我接着说：\par

\Signature{犹斯定：}朋友们，你们是否觉察到，天主说他要将光荣赐给他所立为外邦人之光的那一位，而不赐给别的；并不像特里丰所说的，天主将光荣留给自己？\par

\Signature{特里丰：}我们也觉察到了；所以请继续讲接下来的话。\par

% structure: p0078 source=h2 target=subsection reason=relative-heading-rank; base=h2

\subsection{第六十六章 他从依撒意亚证明天主是由童贞女所生}

我便从先前停顿之处接续话题，证明他是由童贞女所生，并且依撒意亚早已预言了他的童贞女诞生，于是我又引用了同一段预言。内容如下：\Quote{上主又对阿哈次说：\InnerQuote{你向上主你的天主要一个征兆；或求诸深渊，或求诸高天。}阿哈次说：\InnerQuote{我不要求，我不试探上主。}依撒意亚说：\InnerQuote{达味家啊，你们听吧！你们使人厌烦还算小事，还要使我的天主厌烦吗？因此，上主要亲自给你们一个征兆：看，有位童贞女要怀孕生子，给他起名叫厄玛奴耳。他要吃奶酪和蜂蜜，直到他晓得弃恶择善。因为，在这孩子还不晓得弃恶择善以前，他必在亚述王面前接受大马士革的财富和撒玛黎雅的掠物。那因两个君王而令你们难以忍受的土地必将荒芜。但上主要使你们、使你的百姓、你的父家，自厄弗辣因离开犹大以来，未曾有过的日子临到你们。}}\par

\Signature{犹斯定：}如今，所有人都清楚，在亚巴郎血统的肉身后裔中，除了我们的这位基督，没有一个人是由童贞女所生，也没有人被称为由童贞女所生。\par

% structure: p0081 source=h2 target=subsection reason=relative-heading-rank; base=h2

\subsection{第六十七章 特里丰将耶稣与珀耳修斯相比；并宁愿说他是因遵行法律而被选为基督。犹斯定谈及法律是过去的}

\Signature{特里丰：}圣经并没有说\BibleQuote{看，有位童贞女要怀孕生子}，而是说\InnerQuote{看，有位年轻女子要怀孕生子}，等等，正如你所引用的。但整个预言是指希则克雅，并且根据这预言的条款，它已经应验在他身上。此外，在那些被称为希腊人的神话中，记载珀耳修斯是由一位名叫达娜厄的童贞女所生；在他们当中被称为宙斯的那位曾以金雨的形式降临到她身上。当你做出与他们相似的声称时，你应当感到羞耻；你更应该说这位耶稣是由人而生的人子。如果你能证明他从圣经中是基督，并且由于他遵行法律、品行完美，因而配得上被选为基督的尊荣，那还可以；但不要冒险谈论那些怪异之事，免得你像希腊人一样被指责为胡说八道。\par

\Signature{犹斯定：}特里丰，我愿意说服你，以及所有人在内：即使你用更恶劣的言辞嘲笑戏弄，你也不能动摇我的既定志向；我总要引用那些你认为可以用来证明你观点的言语，结合圣经的见证，来证明我所说的话。然而，你企图推翻我们之间已达成共识的那些事，这并不公平诚实；我们已同意，某些诫命是因你们人民心硬而由梅瑟设立的。因为你说，他因遵行法律而被选并成为基督（如果他被证明是这样的话）。\par

\Signature{特里丰：}你向我们承认，他受过割礼，并遵守了梅瑟所规定的其他法律礼仪。\par

\Signature{犹斯定：}我承认过，并且现在也承认；然而我承认他忍受这一切，并非因这些而称义，而是完成他的圣父——万有的创造者、主和天主——愿意他去完成的救世工程。因为我承认他忍受了十字架刑和死亡，以及你们的民族加在他身上的降生成人和诸多苦难。但是，既然你又不同意你刚才所同意的，特里丰，请回答我：那些生活在梅瑟之前的义人圣祖们，他们并未遵守任何那些经上表明在梅瑟之后才开始设立的礼仪，他们是否得救，并获得了真福的产业？\par

\Signature{特里丰：}圣经迫使我承认这一点。\par

\Signature{犹斯定：}同样，我再问你：天主命令你们的祖先奉献供物和祭献，是因为他需要它们，还是因为他们心硬、易于偶像崇拜呢？\par

\Person{特黎丰}：后者也照样是圣经迫使我们承认的。\par

\Person{尤斯定}：同样地，圣经岂没有预言说，天主应许要订立一个新盟约，不同于在曷勒布山上所立的那一个吗？\par

\Person{特黎丰}：这事也早已预言过了。\par

\Person{尤斯定}：旧的盟约岂不是在恐惧与战栗中赐给你们祖先的吗？以致他们不能聆听天主的声音。\par

\Person{特黎丰}：他承认了这一点。\par

\Person{尤斯定}：那么怎样呢？天主应许要另立一个新盟约，不像那旧盟约；他说，这新盟约要在没有恐惧、战栗和闪电的情况下赐下，它要显示出哪些命令和行为是天主知道是永恒的、适合万民的，以及祂因你们人民心硬而赐下了哪些诫命，正如祂藉先知所宣告的。\par

\Person{特黎丰}：对于这一点，凡是爱好真理而非喜爱争辩的人，也必然赞成。\par

\Person{尤斯定}：我不知道你怎么说那些非常爱争论的人，因为你本人就多次明显表现出这种态度，常常反驳你以前同意过的事。\par

% structure: p0096 source=h2 target=subsection reason=relative-heading-rank; base=h2

\subsection{他抱怨特黎丰的固执；回答他的异议；指摘犹太人的不信实}

\Person{特黎丰}：你在试图证明一件难以相信、几乎不可能的事，即天主竟屈尊降生成为人。\par

\Person{尤斯定}：如果我是用人间的教理或论证来证明这一点，你本不该容忍我。但如果我频频引用圣经，而且如此多关于此点的经文，并请求你领悟它们，那么你在认识天主的旨意和心意方面是心硬的。但如果你愿意永远如此，我也不会受丝毫损害；我将永远持守我遇见你之前所持的同样见解，离开你。\par

\Person{特黎丰}：你看，我的朋友，你是经过许多劳苦努力才掌握了这些真理的。因此我们也要仔细审察我们所遇到的一切，以便同意那些圣经迫使我们相信的事。\par

\Person{尤斯定}：我并非要求你不在探究所问之事时全力努力；但我要求你，在无话可说时，不要反驳你曾承认的那些事。\par

\Person{特黎丰}：我们将尽力这样做。\par

\Person{尤斯定}：除了我刚才向你提出的问题，我还想再提出一些；因为藉着这些问题，我将努力使谈话迅速结束。\par

\Person{特黎丰}：请问吧。\par

\Person{尤斯定}：你认为在圣经中，除了万有的创造者和基督之外，还有谁被说成是堪受敬拜、被称为主和天主的吗？基督已被许多圣经证明成了人。\par

\Person{特黎丰}：我们怎能承认这一点？因为我们曾经进行过如此大的探讨，关于除了父以外是否还有另一位。\par

\Person{尤斯定}：我也必须问你这一点，好让我知道你是否与你刚才所承认的意见不同。\par

\Person{特黎丰}：不是的，先生。\par

\Person{尤斯定}：既然你确实承认这些事，而且圣经又说：\BibleQuote{谁能述说祂的世代？}那么你难道不应认为祂不是人类的种子吗？\par

\Person{特黎丰}：那么，圣言如何对\Person{达味}说：天主必从你的腰中为自己立一个儿子，建立祂的王国，将祂安置在荣耀的宝座上呢？\par

\Person{尤斯定}：\Person{特黎丰}，如果\Person{依撒意亚}所说的预言：\BibleQuote{看，童贞女将怀孕}不是针对达味家，而是针对十二支派的另一个家，那么这事或许有些困难；但由于这预言是指着达味家说的，依撒意亚已经解释了天主以奥秘的方式对达味所说的话将如何实现。但朋友们，你们或许不知道：有许多话是以暗晦、比喻、奥秘的方式写成的，还有许多象征性的行动，后来生活在那些说话或行事之人之后的先知解释了它们。\par

\Person{特黎丰}：确实如此。\par

\Person{尤斯定}：因此，如果我能证明依撒意亚的这个预言是指着我们的基督，而不是如你们所说的希则克雅，那么在这件事上，我岂不也要迫使你们不相信你们的老师吗？他们竟敢断言你们与埃及王仆托肋米一起的七十位长老所给出的解释在某些方面是不正确的。因为对于圣经中某些明显暴露他们愚蠢和虚妄见解的语句，他们竟敢断言并非如此记载。但对于另一些他们觉得可以曲解并与人类行为调和的语句，他们则说不是指我们的这位耶稣基督，而是指他们乐意解释的那个对象。例如，他们教导你们说，我们现在讨论的这段圣经是指希则克雅——而我答应过要证明他们是错的。既然他们被迫承认，我们向他们提及的某些圣经——那些明确证明基督要受苦、被敬拜、被称为天主的圣经，并且我已向你们引述过的——确实是指基督，但他们又敢断言这人不是基督。然而他们承认他将要来临，受苦、统治、被敬拜并成为天主；而我也要同样证明这看法是可笑和愚蠢的。但由于我被要求先回答你们打趣所说的话，我就先回答，然后再答复下面的话。\par

\SourceRecord{Translated by Marcus Dods and George Reith.；Ante-Nicene Fathers；Vol. 1.；Edited by Alexander Roberts, James Donaldson, and A. Cleveland Coxe.；Buffalo, NY: Christian Literature Publishing Co.,；1885.；<http://www.newadvent.org/fathers/01285.htm>.}

% structure: p0001 source=h1 target=section reason=page-title

\section{与\Person{特里丰}对话（第69-88章）}

% structure: p0002 source=h2 target=subsection reason=relative-heading-rank; base=h2

\subsection{第69章 魔鬼既然效法真理，就编造了关于\Person{巴克斯}、\Person{赫拉克勒斯}和\Person{埃斯库拉皮乌斯}的寓言}

\Person{儒斯定}：那么，\Person{特里丰}，请确信，我因那被称为魔鬼者据称在希腊人中所行的仿冒品——正如在埃及由术士们所行的，以及在\Person{厄里亚}时代由假先知们所行的——而牢固地确立了对圣经的认识和信德。因为当他们讲述\Person{朱庇特}之子\Person{巴克斯}是由\Person{朱庇特}与\Person{塞墨勒}结合所生，并且他是葡萄藤的发明者；当他们叙述他被撕碎、死去、复活并升天；当他们将酒引入他的奥迹时，我岂不看出（魔鬼）模仿了由族长\Person{雅各伯}预言并由\Person{梅瑟}记载的预言吗？当他们说\Person{赫拉克勒斯}强壮，走遍天下，由\Person{阿尔克墨涅}所生，死后升天，我岂不看出那论\Person{基督}的话——\BibleQuote{强壮如巨人奔跑他的赛程}——也被同样模仿了？当他（魔鬼）推出\Person{埃斯库拉皮乌斯}作为起死回生和医治百病者时，我岂能不说在这件事上他也模仿了关于\Person{基督}的预言？但由于我还没有向你们引用这样的圣经——它说\Person{基督}将做这些事——我必然要提醒你们一处：由此你们可以明白，对于那些不认识天主的人，即外邦人，他们\BibleQuote{有眼看不见，有心不明白}，崇拜木偶，圣经预言他们将会放弃这些虚妄，并寄望于这位\Person{基督}。经上这样记载：\par

\BibleQuote{欢欣吧，干渴的旷野；愿旷野欢乐，如百合花开；约旦的荒野都要开花欢乐；黎巴嫩的光荣加于它，加尔默耳的美誉。我的人民将看见主的崇高，天主的荣耀。你们这些怠惰的手和软弱膝要坚强。心灵软弱的人要振作：坚强，不要害怕。看，我们的天主将给予，且将给予报应的审判。他必来拯救我们。那时瞎子的眼要开启，聋子的耳要听见。那时瘸子要跳跃如鹿，口吃者的舌头要清楚：因为水在旷野中涌出，在干渴之地有溪谷；干地要成为水池，在干渴之地要有水泉涌出。}\BibleRef{依 35:1}\par

从\Person{天主}那里涌出的活水之泉，在那缺乏\Person{天主}知识的土地，即外邦人的土地上，就是这位\Person{基督}，祂也曾在你们的民族中出现，用祂的话医治了那些从母胎中身体残废、耳聋、瘸腿的人，使他们跳跃、听见和看见。并且使死人复活，叫他们存活，祂的行为迫使当时的人认出祂。但尽管他们看见了这样的工作，却断言这是法术。因为他们竟敢称祂为术士和欺世盗名之人。然而祂行了这些事，并说服了那些注定要相信祂的人；因为即使任何人身体有缺陷，却遵守祂所传授的教义，祂将在祂的第二次来临时将他完全健全地举起，在他成为不死、不朽和无忧之后。\par

% structure: p0006 source=h2 target=subsection reason=relative-heading-rank; base=h2

\subsection{第70章 \Person{密特拉}的奥迹同样被曲解自\Person{达尼尔}和\Person{依撒意亚}的预言}

\Person{儒斯定}：当那些记录\Person{密特拉}奥迹的人说他由岩石所生，并称那些相信他的人入教的地方为洞穴时，我岂不看出这里他们模仿了\Person{达尼尔}的预言——一块非人手凿出的石头从大山中被切割出来——并且他们也同样试图模仿\Person{依撒意亚}的全部话语？因为他们编造了同样的正义之言也由他们引用。但我必须向你们重复所指的\Person{依撒意亚}的话，好使你们从中知道事情就是这样。它们是：\BibleQuote{听，远方的人，我所做的事；近处的人要知道我的大能。熙雍的罪人被移开；战栗要抓住不敬虔的人。谁将向你们宣告那永远的地方？那人行正义，说正直的话，憎恨罪恶和不义，使手不沾贿赂，塞耳不听流血的审判，闭眼不看邪恶：他必住在坚固岩穴的高处。必有饼给他，他的水是确实的。你必看见君王带着荣耀，你的眼要远望。你的灵魂要殷勤追求敬畏主。文士在哪里？谋士在哪里？那数算被养育的人——大小人民——的在哪儿？他们不与他们商议，也不知晓声音的深处，以致听不见。那人民已变得贬值，听的人中没有理解。}\BibleRef{依 33:13}现在很明显，在这个预言中（所指）是我们的\Person{基督}给我们的饼，为纪念祂为了相信祂的人而成为肉身，也为他们而受苦；以及祂给我们的杯，为纪念祂自己的血，怀着感恩。这个预言证明我们将看见这位君王带着荣耀；而预言的措辞清楚地宣告，那预知要相信祂的人民，预知要殷勤追求敬畏主。此外，这些圣经同样明确地说，那些被认为是通晓圣经著作并听见预言的人，没有理解。当我听见，\Person{特里丰}，\Person{帕尔修斯}是由一位贞女所生时，我明白那欺骗人的蛇也伪造了这一点。\par

% structure: p0008 source=h2 target=subsection reason=relative-heading-rank; base=h2

\subsection{第71章 犹太人拒绝\WorkTitle{七十贤士译本}的解释，而且他们还从中删除了某些段落}

\Person{儒斯定}：但我远不信任你们的教师，他们拒绝承认与埃及王仆托肋买同在的七十位长老所作的译本是正确的；他们企图另造一种译本。我希望你们注意，他们已完全删除了许多由与仆托肋买同在的那七十位长老翻译的经文，而正是通过这些经文，这位被钉十字架的人被明确地证明是天主和人，并被钉十字架，并死亡；但我知道你们全民族否认这一点，所以我不讨论这些点，而是继续用你们仍然承认的那些段落进行辩论。因为你们同意我提出来引起你们注意的那些话，除了你们反驳\Quote{看，童贞女将怀孕}这句话，并说应当读作\Quote{看，年轻女子将怀孕}。我曾应许要证明这预言所指的不是像你们所学的希则克雅，而是我的这位基督：现在我要开始证明。\par

\Person{特黎丰}：我们首先要求你告诉我们一些你声称已被完全删除的经文。\par

% structure: p0011 source=h2 target=subsection reason=relative-heading-rank; base=h2

\subsection{第七十二章 犹太人从厄斯德拉和耶肋米亚中删除的段落}

\Person{儒斯定}：我照你的意思做。那么，从厄斯德拉关于逾越节律法的言论中，他们删除了以下内容：\Quote{厄斯德拉对百姓说：\InnerQuote{这个逾越节是我们的救主和我们的避难所。如果你们明白了，并且心里领受了，我们将祂放在一个标记上使祂谦卑，此后希望于祂，那么这地方将永不被遗弃——万军的天主说。但如果你们不信任祂，不听从祂的宣告，你们将成为外邦人的笑柄。}}又从耶肋米亚的言论中，他们删除了以下内容：\Quote{我好像一只被牵到宰杀之地的羔羊；他们设谋害我，说：\InnerQuote{来吧，我们把木头放在祂的饼上，把祂从活人之地除掉，使祂的名字不再被记念。}}\BibleRef{耶 11:19}既然耶肋米亚言论中的这段话仍写在犹太人会堂中的一些圣经抄本里（因为被删去的时间不长），并且从这些话可以证明犹太人曾商议关于基督本人，要钉死祂并杀害祂，祂自己也被宣告如同羊被牵到宰杀之地，正如依撒意亚所预言的，并且在这里被描绘成一只无害的羔羊；但他们面对这些话感到困惑，便投身于亵渎。此外，从同一个耶肋米亚的言论中，这些也被删除了：\Quote{上主天主记念祂死在坟墓中的以色列百姓；祂降下来向他们宣讲祂自己的救恩。}\par

% structure: p0013 source=h2 target=subsection reason=relative-heading-rank; base=h2

\subsection{第七十三章 ［这些话］\BibleQuote{从木头上}已从第九十六篇圣咏中被删去}

\Person{儒斯定}：从第九十六篇圣咏中，他们删除了达味这句话的简短说法：\BibleQuote{从木头上。}因为当经文说：\BibleQuote{在外邦中传报，上主从木头上为王了}，他们留下了\BibleQuote{在外邦中传报，上主为王了。}你们民族中从未有人曾被说成作为天主和主在外邦中为王，唯独那位被钉十字架的，关于祂，圣神在同一篇圣咏中证实祂已复活，并从坟墓中得释放，宣告在列邦的神明中没有像祂的：因为它们都是魔鬼的偶像。但我要把整篇圣咏念给你们听，好让你们明白所说的内容。它是这样的：\BibleQuote{请向上主高唱新歌，普世大地，请向上主歌唱！请向上主歌唱，赞美祂的圣名，日日宣扬祂的救恩！请在列邦中传述祂的光荣，在万民中传述祂的奇事！因为上主是伟大的，极应受赞美，超越众神当受敬畏。因为列邦的众神都是魔鬼，而上主创造了诸天。尊荣和华美在祂面前，圣洁和威严在祂圣所。请将光荣和力量归于上主，将上主之名的光荣归于上主！请携带祭品，进入祂的庭院；请在上主圣殿内朝拜上主！愿全地在祂面前战栗：请在列邦中传报，上主为王了！祂奠定了世界，使不致摇荡，祂将以公平审判万民。愿诸天欢乐，大地踊跃，愿海和其中的一切澎湃！愿原野和其中的一切欢腾！愿林中的树木在上主面前欢乐，因为祂来了，祂来审判大地。祂将以正义审判世界，以真理审判万民。}\par

\Person{特黎丰}：百姓的首领是否如你所说删除了圣经的任何部分，只有天主知道；但这似乎难以置信。\par

\Person{儒斯定}：诚然，这似乎难以置信。因为这事比他们在土地上饱食玛纳后所造的金牛犊，或比将儿女祭献给魔鬼，或比杀害先知更为可怕。但在我看来，你似乎没有听过我所说的他们偷走的圣经。因为已经引用的那些，除了我们保留的和将要提出来的之外，已足够证明有争议的问题。\par

% structure: p0017 source=h2 target=subsection reason=relative-heading-rank; base=h2

\subsection{第七十四章 第九十六篇圣咏的开头被特黎丰归于圣父。但通过\BibleQuote{在外邦中传报，主…}等话，它指的是基督。}

\Person{特黎丰}：我们知道你引用这些是因为我们要求你。但在我看来，你最后引用的达味的这篇圣咏所指的，除了圣父和天地的创造者之外，并非其他任何人。而你却断言它指的是那位受难者，你也在热切地努力证明祂就是基督。\par

\Emph{尤斯定：}请你们注意，我要讲述圣神在这篇圣咏中所说的话；你们将会知道，我没有犯罪地说，我们也没有真的被迷惑；因为这样你们自己就能明白圣神所说的许多其他话。\newline{}\BibleQuote{你们应向主高唱新歌，普世大地，应向上主歌唱；应向上主歌唱，赞美祂的圣名，日复一日宣示祂的救恩，在万民中传扬祂的奇事。}祂命令所有知道这救恩奥迹的世上居民，即基督的苦难（祂借此拯救了他们），歌唱并赞颂万有之父天主，承认祂应受赞美和敬畏，祂是造天地的主，为人类成就了这救恩，祂也曾被钉死，并由祂（天主）认为配得统治全地。正如［清楚可见］于祂所说要领你们祖先进入的那地；［因为祂这样说：］\BibleRef{申 31:16}\BibleQuote{这百姓［将随从外神］，背弃我，破坏我当日与他们所立的盟约；那时我必离弃他们，掩面不顾他们；他们必被吞灭，许多灾祸患难必临到他们；他们在那日必说：因为我天主不在我中间，这些灾祸才临到我。在那日，我必因他们所行的一切恶事掩面不顾他们，就是因为他们转向了别的神。}\par

% structure: p0020 source=h2 target=subsection reason=relative-heading-rank; base=h2

\subsection{第七十五章 证明耶稣是\WorkTitle{出谷纪}中天主的名号}

\Emph{尤斯定：}此外，在\WorkTitle{出谷纪}中，我们也看到天主自己的名号——祂说没有启示给亚巴郎或雅各伯——就是耶稣，并由梅瑟以奥秘的方式宣告。经上这样记载：\BibleQuote{上主对梅瑟说：你告诉这百姓：看，我派遣我的使者在你前面，在路上保护你，领你到我所预备的地方。你们应谨慎听从他，服从他，不可违抗他。因为他决不会背弃你们，因为我的名号在他内。}\BibleRef{出 23:20}现在请明白，那领你们祖先进入那地的人被称为这名号耶稣，起初称为\OriginalTerm{敖瑟亚}{Auses}（\OriginalTerm{敖舍亚}{Oshea}）。\BibleRef{户 13:16}因为如果你们明白这事，你们也会知道那对梅瑟说\Quote{我的名号在他内}的祂，其名号就是耶稣。的确，祂也被称为以色列，雅各伯的名号也改为此。依撒意亚显示，那些被派遣传报天主信息的先知被称为祂的使者和宗徒。因为依撒意亚在某处说：\BibleQuote{请派遣我。}\BibleRef{依 6:8}而那改名的先知耶稣［若苏厄］是刚强伟大的，这是众所周知的。既然我们知道天主曾以如此多的形态向亚巴郎、雅各伯和梅瑟启示自己，我们为何困惑而不相信，按照万有圣父的旨意，祂可能由童贞女降生成人？尤其我们有这些圣经，从中可以清楚看出祂按圣父的旨意成了人。\par

% structure: p0022 source=h2 target=subsection reason=relative-heading-rank; base=h2

\subsection{第七十六章 从其他经文证明基督同样的威严与统治}

\Emph{尤斯定：}因为当达内尔谈到\Quote{一位相似人子者}领受了永远的王国时，他不是暗示这件事吗？因为他说，在说\Quote{相似人子者}时，祂显现了，并是人，但不是出于人的种子。他也在谈论这块不用手凿出的石头时，奥秘地宣告了同样的事。因为\Quote{不用手凿出}这句话表示这不是人的工作，而是万有圣父和天主的旨意的工作，祂带来了祂。当依撒意亚说：\Quote{谁能述说祂的后裔？}他是指祂的出身无法言传。凡是从人而生的，没有一个出身是无法言传的。当梅瑟说祂要在葡萄血中洗自己的衣服时，这不是指明我现在常对你们说的那模糊的预言，即祂有血，但不是从人来的；正如不是人，而是天主，生了葡萄的血？当依撒意亚称祂为\Quote{奇妙的谋士}的天使时，岂不是预表祂是那些真理的导师，即祂来［地上］时所教导的？因为唯独祂公开教导了那些奇妙的谋略，就是圣父为所有曾经和将要中悦祂的人，以及那些违背祂旨意的人（无论是人还是天使）所设计的，祂说：\BibleQuote{将有许多人从东方［和西方］来，同亚巴郎、依撒格、雅各伯一起坐席在天国里；本国的子民却被赶到外面的黑暗中。}\BibleRef{玛 8:11}又说：\BibleQuote{有许多人在那日要对我说：主啊，主啊，我们不是因你的名吃过、喝过、说过预言、赶过魔鬼吗？我要对他们说：离开我。}\BibleRef{玛 7:22}又在别的言语中，祂要谴责不配得救的人，祂说：\BibleQuote{离开我，进入那为魔鬼和他的使者所预备的永火里去。}\BibleRef{玛 25:41}又在别的言语中，祂说：\BibleQuote{我已经给你们权柄，可以践踏蛇、蝎子、\OriginalTerm{蜈蚣}{\GreekText{σκολόπενδρα}}和一切仇敌的权势。}现在我们这些信靠被本丢·彼拉多钉死的我们的主耶稣的人，当我们驱逐所有魔鬼和邪灵时，它们都服从我们。因为如果先知们隐约预告基督要受苦，之后成为万有的主，但没有人能明白这话，直到祂自己使宗徒们确信这些话是圣经中明确记载的。因为祂在被钉十字架前呼喊说：\BibleQuote{人子必须受许多苦，被经师和法利塞人弃绝，并被钉死，第三天要复活。}\BibleRef{路 9:22}达味也曾预言，祂要按圣父的旨意，在太阳和月亮之前从母腹出生，并显明祂是基督，是刚强的天主，应受敬拜。\par

% structure: p0024 source=h2 target=subsection reason=relative-heading-rank; base=h2

\subsection{第七十七章 他回头解释依撒意亚的预言}

\Emph{特黎丰：}我承认如此这般重大证据足以说服人；但希望你知道，我是为求你提供你常许诺要给我的证据。那么，请继续向我们阐明这一点，好让我们看到你如何证明那[段落]指着你的这位基督。因为我们声称这预言与希则克雅有关。\par

\Emph{犹斯定：}我愿照你所愿而行。但首先请你们自己向我说明，希则克雅如何能在未识呼父呼母之前，当着亚述王的面获取大马士革的权势和撒玛黎雅的掠物？因为你们不能按照你们希望的解说，认为希则克雅曾当着亚述王的面与大马士革和撒玛黎雅的居民作战。因为先知的话说：\BibleQuote{孩子尚未知呼父呼母以前，他要当着亚述王的面夺取大马士革的权势和撒玛黎雅的掠物。}假如先知之神没有附加这句话\BibleQuote{孩子尚未知呼父呼母以前，他要夺取大马士革的权势和撒玛黎雅的掠物}，而只是说\BibleQuote{她要生一个儿子，他必夺取大马士革的权势和撒玛黎雅的掠物}，那么你们可以说天主预知这事，便预言他将夺取这些东西。但现在预言附加了这句话：\BibleQuote{孩子尚未知呼父呼母以前，他要夺取大马士革的权势和撒玛黎雅的掠物。}而你们无法证明在犹太人中有谁曾发生过这样的事。但我们能证明这事发生在我们基督身上。因为在他诞生时，从阿剌伯来的贤士去朝拜他，他们先到当时统治你们土地的\Person{黑落德}那里，圣经因他的不虔和罪恶品格称他为亚述王。因为你们知道，圣神常用比喻和图像宣示这类事件，正如他曾对耶路撒冷全体人民所说：\BibleQuote{你的父亲是阿摩黎人，你的母亲是赫特人。}\BibleRef{则 16:3}\par

% structure: p0027 source=h2 target=subsection reason=relative-heading-rank; base=h2

\subsection{第七十八章　他证明这预言只与基督相符，基于随后所记载的话。}

\Emph{犹斯定：}当贤士从\Place{阿剌伯}来到\Person{黑落德}那里，说他们因天空出现的一颗星知道在你们国内有一位君王已诞生，并且来朝拜他时，\Person{黑落德}王从你们百姓的长老那里得知，先知论\Place{白冷}如此记载：\BibleQuote{你，白冷，地属犹大，在犹大首领中绝非最小；因为将由你出来一位领袖，祂将牧养我的百姓。}\BibleRef{米 5:2}于是，来自\Place{阿剌伯}的贤士来到\Place{白冷}，朝拜了那婴儿，并将礼物——黄金、乳香和没药献给他；但在\Place{白冷}朝拜婴儿后，因在启示中被警告，就没有回到\Person{黑落德}那里。\Person{玛利亚}的丈夫\Person{若瑟}，起初想休退他的未婚妻\Person{玛利亚}，以为她因与男人交合而怀孕，即因奸淫，却在异象中被命令不可休退他的妻子；那向他显现的天使告诉他，她腹中所怀的是因圣神而有的。于是他惧怕，没有休退她；但在\Person{季黎诺}手下第一次在\Place{犹大}登记户口时，他离开居住的\Place{纳匝肋}，上到所属的\Place{白冷}去登记，因为他的家族属犹大支派，那时他们居住在那地区。随后，他受命与\Person{玛利亚}一同前往\Place{埃及}，与婴儿留在那里，直到另一个启示警告他们返回\Place{犹大}。当婴儿在\Place{白冷}诞生时，因为\Person{若瑟}在那村庄找不到住处，便住在村庄附近的一个山洞里；他们在那里时，\Person{玛利亚}生了基督，将他放在马槽里，来自\Place{阿剌伯}的贤士在那里找到了他。\par

我已向你们重复了\Person{依撒意亚}所预告的预兆，那预兆预表了山洞；但为了今天与我们同来的人，我要再次提醒你们这段经文。\par

于是，我重复了我已写下的\Person{依撒意亚}的经文，并补充说，借着这些话，主持\OriginalTerm{密特拉}{Mithras}奥秘的人被魔鬼所激动，声称在他们称为山洞的地方，他们由他（魔鬼）而得到启蒙。\par

所以，黑落德见从阿剌伯来的贤士没有按他所要求的回到他那里，反而遵照给他们的命令由另一条路返回本国去了；而若瑟带着玛利亚和婴孩也已按照启示进入埃及；他不知道贤士们曾去朝拜的婴孩是谁，就下令将白冷当时所有的婴孩一概屠杀。耶肋米亚藉圣神预言此事将发生，说：\Quote{在辣玛听到了声音，哀号与极大的悲号，辣黑耳痛哭她的子女，她不愿受安慰，因为他们不在了。}\BibleRef{耶 31:15}因此，由于从辣玛——即阿剌伯（因为阿剌伯至今有一个地方叫辣玛）——发出的声音，哀号将临到雅各伯（又称以色列）的神圣先祖辣黑耳被埋葬的地方，即白冷；妇女们为自己被杀的儿女哭泣，因所遭遇的事而得不到安慰。因为依撒意亚那句话\Quote{他必夺去大马士革的权势和撒玛黎雅的掠物}，预言了那住在大马士革的恶魔的权势在基督一出生时就被祂制服；这已被证明发生了。因为贤士们原被那恶魔的权势捆绑，行各种恶事，但通过来朝拜基督，表明他们已经摆脱了那俘虏他们的统治；而圣经告诉我们这统治驻扎在大马士革。此外，那有罪不义的权势在比喻中也被恰当地称为撒玛黎雅。你们中没有人能否认大马士革过去和现在都在阿剌伯地区，虽然它现在属于所谓叙利亚-腓尼基。因此，诸位先生，你们应当从那些从天主领受恩宠的人——即我们基督徒——学习你们所没有领悟的事，而不应竭尽全力维护自己的道理，羞辱天主的道理。因此，这恩宠也转移给了我们，正如依撒意亚所说的：\Quote{这百姓亲近我，用嘴唇尊敬我，心却远离我；他们恭敬我也是枉然，教导人的命令和道理。因此，看，我要继续迁移这百姓，我要迁移他们；我要除去智慧人的智慧，使明达人的明达归于无有。}\BibleRef{依 29:13}\par

% structure: p0032 source=h2 target=subsection reason=relative-heading-rank; base=h2

\subsection{第79章 他证明反对特里丰的论点：恶天使背叛了天主}

特里丰有些生气，但仍尊重圣经，从他的表情可以看出来，他对我说：\par

\Emph{特里丰：}天主的言语是神圣的，但你的解释纯属编造，从你刚才所解释的可以清楚看出；甚至还是亵渎，因为你断言天使犯了罪并背叛了天主。\par

而我，希望他能听我说，就用更温和的语气回答，说：\par

\Emph{尤斯定：}先生，我钦佩您的这份敬虔；我祈祷您也能对那位天使被记载服事的天主持有同样的心情，正如达尼尔所说的：一位相似人子的被引领到万古常存者面前，一切王国的权柄都赐给了祂，直到永远。但是，先生，为了让您知道并非我们的鲁莽使我们采用这种您所指责的解释，我将从依撒意亚本人那里给您证据；因为他证实恶天使曾在埃及的塔尼斯居住，并且现在仍居住在那里。以下是他的话：\Quote{祸哉，叛逆的子女！上主这样说：你们商议，却不藉着我；结盟，却不藉我的神，以致罪上加罪；他们因下埃及而犯罪，却不求问我，好得法郎的援助，并受埃及的荫庇。因为法郎的荫庇必成为你们的羞辱，成为那些信赖埃及者的侮辱；因为塔尼斯中的首领是恶天使。他们为那不能帮助他们的人民劳碌是枉然的，终必成为羞辱和侮辱。}\BibleRef{依 30:1}而且，匝加利亚也告诉我们，正如您自己所提到的，魔鬼站在耶叔亚大司祭的右边，要反对他；主说：\Quote{拣选耶路撒冷的上主斥责你。}\BibleRef{匝 3:1}此外，约伯传中也记载，正如您自己所说，天使们来站在上主面前，魔鬼也同他们一起来。\BibleRef{约 1:6}并且我们有梅瑟在创世纪开头记载的：蛇欺骗了厄娃，并受了诅咒。我们知道在埃及有术士，他们模仿天主通过忠信仆人梅瑟显示的大能。您也知道达味曾说：\Quote{列邦的神祇是魔鬼。}\par

% structure: p0037 source=h2 target=subsection reason=relative-heading-rank; base=h2

\subsection{第80章 尤斯定关于千年统治的意见。几个公教徒拒绝它}

\Emph{特里丰：}我向您说过，先生，您非常渴望在所有事情上都稳妥，因为您紧握圣经。但是请告诉我，您真的承认这个地方——耶路撒冷——将被重建吗？您期望您的子民被聚集，与基督、先祖和先知们一同喜乐，既包括我们民族的人，也包括在你们的基督来临之前加入他们的其他皈依者？或者您让步了，承认这一点是为了在辩论中显出胜过我们的样子？\par

\Person{尤斯定}：我不是如此可怜的人，\Person{特里丰}，竟会说一套而想另一套。我先前曾向你承认，我和其他许多人都持此见解，且相信这样的事将会发生，正如你确知的那样；但另一方面，我也向你表示，许多属于纯正虔敬信德的人，是真基督徒，却持不同看法。此外，我曾向你指出，有些被称为基督徒的人，却是无神、不虔的异端者，他们教导各种亵渎、无神、愚昧的道理。但为让你知道我不只是在你面前这样说，我将尽我所能，把我们之间所论的一切论点写下来；在其中，我将记录我自己承认那些我向你承认的同样事情。因为我选择不追随人及人的道理，而是追随天主及祂所传授的道理。你若遇到一些自称基督徒的人，他们却不承认这一真理，竟敢亵渎\Person{亚巴郎}的天主、\Person{依撒格}的天主、\Person{雅各伯}的天主；他们声称没有死人的复活，且他们的灵魂死后升天；不要以为他们是基督徒，就如一个人若正确思量，就不会承认撒杜塞人，或类似的基尼斯特派、默黎斯特派、加里肋亚人、赫肋尼派、法利塞人、洗礼派是犹太人（我告诉你我的想法时，请勿急躁），他们只是被称为犹太人和\Person{亚巴郎}的子孙，以嘴唇尊敬天主，但心却远离祂，正如天主自己所宣告的。但我及其他在各方面持正统思想的基督徒确信，将有死人的复活，以及在耶路撒冷的一千年，那城那时将被建造、装饰并扩大，正如\Person{厄则克耳}、\Person{依撒意亚}及其他先知所宣告的。\par

% structure: p0040 source=h2 target=subsection reason=relative-heading-rank; base=h2

\subsection{第八十一章 他试图以\BibleBook{依撒意亚}{}和\BibleBook{}{默示录}证明此意见}

\Person{尤斯定}：关于这一千年时期，\Person{依撒意亚}曾如此说：\par

\BibleQuote{因为将要有新天新地，先前的事不再被记忆，也不再上人心；然而他们要在其中欢喜快乐，这些是我所创造的。因为，看，我要使耶路撒冷成为喜乐，使我的百姓成为欢欣；我要因耶路撒冷而喜乐，因我的百姓而欢欣。其中不再听到哭泣的声音，也不再听到哀号的声音。那里再没有夭折的婴儿，也没有不满寿数的老人；因为百岁的人算是青年，但百岁死去的罪人应受诅咒。他们要建造房屋，并自己居住；他们要栽种葡萄园，并吃其果实。他们建造的，别人不得居住；他们栽种的，别人不得吃。因为我百姓的日子将如树木的年数，他们劳苦的成果必丰盈。我的选民不徒然劳苦，不生孩子遭诅咒；因为他们是蒙主祝福的正义苗裔，他们的后代也同他们一样。将来，他们尚未呼求，我已应允；他们还在说话时，我就说：什么事？那时，豺狼与羔羊一同觅食，狮子像牛一样吃草，蛇以尘土为食。在我的圣山上，他们不再伤害或虐待，上主说。}\par

现在我们明白了，在这些话中使用的表达\BibleQuote{我百姓的日子将如树木的年数，他们劳苦的成果必丰盈}隐晦地预言了一千年。因为正如亚当被告知，在他吃树的那一天他必死，我们知道他并没有活满一千年。此外，我们感知到，\BibleQuote{主的日子如一千年}这一表达与此主题相关。此外，我们这里有一位名叫\Person{若望}的人，是基督的宗徒之一，藉着给他的启示，他预言那些相信我们基督的人，将在耶路撒冷居住一千年；此后，普遍的、一句话，永久的复活和所有的人的审判也将同样发生。正如我们的主也说：\BibleQuote{他们不娶也不嫁，却要与天使相等，成为复活的天主的儿女。}\BibleRef{路 20:35}\par

% structure: p0044 source=h2 target=subsection reason=relative-heading-rank; base=h2

\subsection{第八十二章 犹太人的先知恩赐被转移给基督徒}

\Person{Justin}: 因为先知的恩赐至今仍在我们中间。因此你们应当明白，[这些恩赐]从前在你们民族中，如今已转移给我们。正如在你们圣先知的时代有假先知一样，现在我们中间也有许多假教师，我们的主曾预先警告我们要提防他们；所以我们在任何方面都不欠缺，因为我们知道祂预知一切将发生在我们身上的事，就是祂从死者中复活并升天之后。因为祂说过，我们会被处死，并因祂的名而被憎恨；而且许多假先知和假基督会以祂的名出现，迷惑许多人：事情正是如此。因为许多人教导了不虔敬、亵渎和不圣洁的道理，假借祂的名编造；他们也教导，甚至至今仍在教导那些来自魔鬼不洁之灵的、放在他们心里的东西。因此我们非常希望你们被说服，不要被这样的人误导，因为我们知道，凡能说真理却不说的人，必受天主的审判，正如天主通过厄则克耳作证说：\BibleQuote{我立你作犹大家族的守望者。如果罪人犯罪，你不警告他，他必死在罪中；但我要向你追讨他的血债。如果你警告他，你就无罪。}因此我们出于恐惧，非常热切地渴望按照圣经与人交谈，而不是出于贪财、贪荣或贪乐。因为没有人能在这类[恶习]上指控我们。我们也不愿像你们人民的首领那样生活，天主斥责他们说：\BibleQuote{你们的首领是盗贼的同伙，喜爱贿赂，追逐报酬。}\BibleRef{依 1:23} 如果你们知道我们中间有人如此，不要因他们的缘故亵渎圣经和基督，也不要竭力歪曲解释。\par

% structure: p0046 source=h2 target=subsection reason=relative-heading-rank; base=h2

\subsection{第八十三章 证明《圣咏》中\Quote{上主对我主说}等句不适合希则克雅}

\Person{Justin}: 因为你们的教师竟敢将这句经文\BibleQuote{上主对我主说：你坐在我右边，等我使你的仇敌作你的脚凳}应用于希则克雅，仿佛当亚述王派人来威胁他时，他被邀请坐在圣殿的右边，而依撒意亚告诉他不要害怕。我们知道并承认依撒意亚所说的事发生了：亚述王在希则克雅的日子停止了攻打耶路撒冷，上主的天使击杀了亚述军队约十八万五千人。但这圣咏显然不是指他。因为经上这样写着：\par

\BibleQuote{上主对我主说：你坐在我右边，等我使你的仇敌作你的脚凳。祂要从耶路撒冷伸出权杖，并在你仇敌中掌权。在圣者的荣光中，在晨星之前，我已生了你。上主起了誓，决不反悔：你按照默基瑟德的品位，永为司祭。}\par

那么，谁不承认希则克雅不是按照默基瑟德的品位永为司祭呢？谁不知道他不是耶路撒冷的救赎者呢？谁不知道他并没有向耶路撒冷伸出权杖，也没有在仇敌中掌权，而是天主在他在哀伤和痛苦中为他击退了敌人呢？而我们的耶稣，虽尚未在荣耀中来临，却已向耶路撒冷伸出了权杖，即那对一切曾被魔鬼辖制的万民的召唤和悔改之道，正如达味所说：\BibleQuote{万国的神祇都是魔鬼。}祂强有力的道已使许多人离弃他们素来事奉的魔鬼，并藉此相信全能的天主，因为万国的神祇都是魔鬼。我们先前也提到，\BibleQuote{在圣者的荣光中，在晨星之前，我已从母腹中生了你}这句话是对基督说的。\par

% structure: p0050 source=h2 target=subsection reason=relative-heading-rank; base=h2

\subsection{第八十四章 证明\Quote{看，一位贞女}等预言唯独适合基督}

\Person{Justin}: 此外，\BibleQuote{看，贞女将怀孕生子}这段预言是针对祂说的。因为如果依撒意亚所指的那位不是由贞女所生，那么圣神为何宣告：\BibleQuote{看，上主自己要给你们一个征兆：看，贞女将怀孕生子}？因为如果祂也要像所有长子一样由性交而生，天主为什么要说祂要给出一个并不常见的征兆呢？真正构成征兆且能取信于人的，乃是万物的首生者要藉贞女的子宫降生成人，成为一个婴孩；这正是祂藉着预言之神以多种方式预见并预言了的，正如我曾向你们反复说明的；这样，当事情发生时，众人就知道这是全能者权能与旨意运作的结果，正如厄娃是从亚当的一根肋骨造成，所有的生物在起初都是藉天主的圣言所造。但你们在这些事上竟敢歪曲你们的先辈与埃及王仆托勒密一同翻译的经文，声称经文并非如他们所解释的那样，而是说\BibleQuote{看，那年轻女子将怀孕}，好像一个妇女性交生子就是大事；这确实是所有年轻女子（除不育者外）都做到的；但即使是不育的，天主若愿意，也能使她们生育。因为撒慕尔的母亲原是不育的，却因天主的旨意生了孩子；圣祖亚巴郎的妻子也是这样；还有生若翰洗者的依撒伯尔，以及其他类似的事。所以你们不要以为天主不能做任何祂愿意的事。尤其是当这事已被预言时，你们不要试图歪曲或误解预言，因为你们只会伤害自己，而不会伤害天主。\par

% structure: p0052 source=h2 target=subsection reason=relative-heading-rank; base=h2

\subsection{第八十五章 基于《圣咏》第二十四篇及基督对魔鬼的权柄，证明基督是万军的上主}

儒斯定：此外，你们中有些人竟敢将那句预言解释为指希则克雅，另一些人则解释为指撒罗满，那预言说：\BibleQuote{众城门，你们要抬起头来；永久的门，你们要被举起，光荣的君王要进来。}但可以证明，这句预言既不指后者，也不指前者，总之不指你们任何一位君王，而唯独指我们的这位基督，祂无美貌可夸，也无荣耀，正如依撒意亚、达味以及全部圣经所说；祂是万军之主，由赋予祂尊荣的圣父的旨意所立；祂也是从死者中复活、升天的那一位，如圣咏和其余圣经所宣告的，称祂为万军之主。对于这一点，如果你们愿意，很容易被眼前发生的事说服。因为每个魔鬼，只要奉这位天主子——祂是一切受造物的首生者，由童贞女降生成人，受难，被你们的民族在般雀·比拉多手下钉在十字架上，死亡，从死者中复活，升天——的名字被驱赶，就必被制服。但如果你们奉你们中间任何一位——无论是君王、义人、先知或先祖——的名驱魔，魔鬼不会服从你们。但如果你们中有人奉亚巴郎的天主、依撒格的天主、雅各伯的天主之名驱魔，或许它会服从你们。确实，我先前说过，你们的驱魔者行驱魔时多行诡诈，甚至像外邦人一样，使用熏香和咒语。但那些就是天使与掌权者，达味预言的话语命令它们抬起城门，好让从死者中复活的那一位，耶稣基督，万军之主，按照圣父的旨意进入；达味的这句话同样表明了这一点。为了那些昨天没有与我们在一起的人，我今天再次提醒你们注意；而且为了他们，我还总结了许多我昨天说过的话。现在，尽管我已多次重复，但我对你们说这些话，我知道并非荒谬。因为看到太阳、月亮和其他的星体持续保持同一轨道，带来不同的季节；又看到计算者被问及二乘二是多少，因他屡次说是四，仍不停止再说一遍是四；同样，其他一些确信无疑的事，也照样不断被提及和承认；然而，一个以先知圣经为论据的人，却要离开并停止不断引用同一圣经，只因认为他能拿出比圣经更好的东西，这是何等荒谬。那么，我借以证明天主启示说天上既有天使又有万军的经文是这样的：\BibleQuote{请赞美主自高天：请赞美祂于高处。众天使，请赞美祂；众军旅，请赞美祂。}\par

姆纳塞阿斯（第二天与他们同来的人之一）：我们非常高兴你愿意为我们重述同样的话。\par

儒斯定：朋友们，请听那促使我这样做的圣经。耶稣命令我们甚至要爱仇敌，正如依撒意亚在许多章节中所预言的，其中也包含我们重生的奥秘，事实上，也包含所有期待基督在耶路撒冷显现、并藉其行为努力讨祂喜悦之人的重生。以下是依撒意亚所说的话：\BibleQuote{凡敬畏主话语的人，请听主的话。你们的弟兄，那些恨你们、弃绝你们的人，说：愿主的名受显扬。祂已为你们的喜乐显现，他们却要蒙羞。有喧哗声从城中发出，有声音从圣殿发出，是主对骄傲者施行报应的声音。她尚未临产，痛苦尚未到来，就生了一个男孩。谁听说过这样的事？谁见过这样的事？大地一日之内岂能生产？一个国家岂能在一时诞生？因为熙雍一临产，就生下了她的儿女。我甚至使那不生产的有了这种期待，主说。看，我使那生育的和那不生育的都生育了，主说。耶路撒冷，请欢乐，举行盛会，所有爱慕她的人，请欢欣。所有为她哀悼的人，请喜乐，使你们得以吸吮、满足于她安慰的乳房，使你们吸奶后因祂荣耀的进入而喜悦。}\BibleRef{依 66:5}\par

% structure: p0056 source=h2 target=subsection reason=relative-heading-rank; base=h2

\subsection{第八十六章 旧约中基督借以统治的十字架的木头有多种预像}

\Emph{\Person{贾斯定}:} 那么就请听，圣经所宣告的这位在受难之后将再次带着荣耀来临的人，如何既由被称为种植在乐园中的生命树所象征，也由那些将要临于所有义人的事件所预表。\Person{梅瑟}被差遣带着一根棍杖来拯救人民；他手持此杖在人民前头，分开了海。以此杖他看见磐石涌出水来；当他把一棵树扔进\Place{玛辣}苦水中，水就变甜了。\Person{雅各伯}将木棍放入水槽，使舅父的羊受孕，从而得到羊羔。同一位\Person{雅各伯}夸口说，他藉着棍杖渡过了河。他说他曾看见一架梯子，圣经记载天主站在其上。但我们已从圣经证明那并非圣父。\Person{雅各伯}在同一地方把油倒在石头上，那向他显现的天主亲自作证，他傅油立柱是为那向他显现的天主。那石头象征性地预示基督，我们也已从多段圣经证明；那香膏，无论是油，或是\Emph{\OriginalTerm{香料}{stacte}}，或是任何其他复合的芬芳香脂，都指向祂，我们也已证明，因为经上说：\BibleQuote{因此，天主，你的天主，以喜乐油傅了你，胜过你的同伴（参：\BibleRef{咏45:7}）}。的确，所有君王和受傅者都从祂获得君王和受傅者的名分；正如祂亲自从圣父领受了君王、基督、司铎、天使以及祂所拥有或曾拥有的其他称号。\Person{亚郎}的棍杖开花，表明他是大司铎。\Person{依撒意亚}预言将从叶瑟的根发出枝条，那就是基督。\Person{达味}说义人\BibleQuote{像种在水渠旁的树，按时结果，枝叶不枯（参：\BibleRef{咏1:3}）}。又说义人如棕榈树繁茂。天主曾从一棵树向\Person{亚巴郎}显现，如经上所记，在\Place{玛默勒}橡树旁。人民过\Place{约旦河}后，找到了七十棵柳树和十二股泉水。\Person{达味}确认天主以棍杖和牧杖安慰他。\Person{厄里叟}将一根木棍抛入\Place{约旦河}，捞起了先知弟子们用来砍树建造房屋（他们愿在其中阅读研习天主的法律和诫命）的斧头铁头；同样，我们的基督被钉在木架上，又用水洗净我们，救赎了我们——我们曾沉沦于所犯的最重大过犯中——并使我们成为祈祷与朝拜之所。此外，有一根棍杖藉着伟大的奥迹指出了\Person{塔玛尔}诸子之父\Person{犹大}。\par

% structure: p0058 source=h2 target=subsection reason=relative-heading-rank; base=h2

\subsection{第87章。\Person{特黎丰}反对说：\Quote{将安息在祂身上}等语。\Person{贾斯定}解释之。}

\Emph{\Person{特黎丰}:} 现在，不要以为我问这些问题是想推翻你的说法，而是我希望就我现在询问的要点得到信息。那么，请告诉我，当圣经藉着\Person{依撒意亚}宣称：\BibleQuote{从叶瑟的根将发出一条枝条；从叶瑟的根将开出一朵花；天主的神将安息在祂身上，智慧和聪敏的神，超见和刚毅的神，知识和孝爱的神：敬畏上主的神将充满祂（参：\BibleRef{依11:1-2}）}（现在你向我承认这指的是基督，并且你主张祂是先在的天主，并藉着天主的旨意降生成人，由童贞女而生为人），那么，这位被依撒意亚圣经所列举的圣神能力充满的人，好像祂缺乏它们似的，怎么能证明祂是先在的呢？\par

\Emph{\Person{贾斯定}:} 你问得极为慎重和明智，因为确实似乎有一个难题；但请听我所说，使你也能明白其中的理由。圣经说这些列举的圣神能力降临在祂身上，不是因为祂需要它们，而是因为它们要安息在祂内，即要在祂内得以实现，以致在你们民族中不再有按照古代惯例的先知：这一点你清楚看到了。因为在祂之后，你们中再也没有先知兴起。现在，为使你（知道）你们的先知，每个从天主领受一种或两种能力，实行并说了我们从圣经所学的事，请听我下面的话。\Person{撒罗满}拥有智慧之神，\Person{达内尔}拥有聪敏和超见之神，\Person{梅瑟}拥有刚毅和孝爱之神，\Person{厄里亚}拥有敬畏之神，\Person{依撒意亚}拥有知识之神；其他先知也是如此：每人拥有一种能力，或一种与另一种交替相连；还有\Person{耶肋米亚}，以及十二位先知，和\Person{达味}，总之，你们中的其余先知。因此，当\Emph{祂}来了之后，在祂为人类完成的这个救恩时代的时期，这些恩赐需要从你们中停止，祂便安息了，即停止了；这些恩赐在祂内得到安息后，又如预言的，成为从祂神能力的恩宠出发，祂按照祂认为每人相称的，赐予那些相信祂的人。我已经说过，并再次说，曾预言在祂升天后，这事将由祂完成。因此经上说：\BibleQuote{祂升上高天，掳获了俘虏，将恩赐赐予世人（参：\BibleRef{咏68:18}）}。又在另一预言中说：\BibleQuote{此后，我要将我的神倾注在一切血肉之人身上，我的仆人和我的婢女们都要说预言（参：\BibleRef{岳2:28}）}。\par

% structure: p0061 source=h2 target=subsection reason=relative-heading-rank; base=h2

\subsection{第88章。基督受圣神，非因缺乏。}

\Emph{犹斯定：}现在，我们中间可以看到拥有天主的圣神之恩的男女；因此，\Person{依撒意亚}所列举的能力要临到祂这预言就应验了，不是因为祂需要能力，而是因为那些能力不会在祂之后延续。让这成为你的一个凭证，就是我告诉你的由来自\Place{阿剌伯}的贤士所做的事，他们一见婴孩出生就前来朝拜祂；因为祂在出生时就已拥有祂的能力；祂像所有人一样长大，采用适当的方法，为每个发展阶段分配合适的（需要），并由各种食物滋养，等候了约三十年，直到\Person{若翰}作为祂来临的先行者出现在祂面前，并以圣洗的方式走在祂前面，正如我已表明的。然后，当\Person{耶稣}去到\Person{若翰}施洗的约但河，祂一踏入水中，约但河里就燃起了火；当祂从水中出来时，圣神像鸽子一样降在祂身上，正如我们这位基督的宗徒们所记载的。现在我们都知道，祂到这河来并非因为祂需要圣洗，或需要圣神像鸽子一样降下；正如祂甘愿降生和被钉十字架，不是因为祂需要这些事，而是为了人类，这人类从\Person{亚当}起就处于死亡和蛇的诡诈之下，而且每个人都犯了个人的罪过。因为天主愿意天使和人都拥有自由意志，并能自主地做祂赋予各人能力去做的事，祂如此造了他们，使他们如果选择中悦祂的事，祂就使他们免于死亡和惩罚；但如果他们作恶，祂就照祂认为合适的方式处罚各人。因为祂骑驴进入\Place{耶路撒冷}——我们已表明这曾被预言过——并不是这行为使祂成为基督，而是给人们提供了一个证据，证明祂就是基督；正如在\Person{若翰}的时代，人们需要证据，以便知道谁是基督一样。因为当\Person{若翰}留在约但河边，宣讲悔改的圣洗，只束着皮带，穿着骆驼毛的衣服，只吃蝗虫和野蜜，人们以为他就是基督；但他向他们喊说：\Quote{我不是基督，而是呼喊者的声音；因为有比我更强的那一位要来，我不配提祂的鞋带。}\BibleRef{依 1:27}当\Person{耶稣}来到约但河时，祂被认为是木匠\Person{若瑟}的儿子；祂没有俊美，正如圣经所预言的；祂被视为木匠（因为祂在人间时常常做木工，制作犁和轭，藉此教导正义和积极生活的象征）；然而那时，圣神，为了人的缘故，正如我先前所说的，以鸽子的形状降在祂身上，同时从天上传来声音，这也是\Person{达味}在说话时曾发出的，他代表基督，说出圣父要对祂说的话：\Quote{祢是我的儿子，我今天生了祢。}圣父说，祂的出生是为了人，在他们认识祂的时候：\Quote{祢是我的儿子，我今天生了祢。}\par

\SourceRecord{Translated by Marcus Dods and George Reith.；Ante-Nicene Fathers；Vol. 1.；Edited by Alexander Roberts, James Donaldson, and A. Cleveland Coxe.；Buffalo, NY: Christian Literature Publishing Co.,；1885.；<http://www.newadvent.org/fathers/01286.htm>.}

% structure: p0001 source=h1 target=section reason=page-title

\section{与特里丰对话录（第89-108章）}

% structure: p0002 source=h2 target=subsection reason=relative-heading-rank; base=h2

\subsection{第八十九章 唯独十字架因诅咒令特里丰反感，然而它证明耶稣就是基督}

\Emph{特里丰：}请放心，我们全民族都等待基督；我们也承认你所引的所有经文都指向祂。此外，我也承认，耶稣这名——农的儿子若苏厄以此称呼——强烈地使我倾向于这个观点。但是，基督是否应如此耻辱地被钉十字架，我们对此有所怀疑。因为凡被钉十字架的人，按照法律被称为可诅咒的，所以我在这点上极不相信。确实，经文明明宣告基督必须受苦；但我们希望知道，你是否能向我们证明，祂是否必须受法律所诅咒的苦。\par

\Emph{犹斯定：}如果基督不必受苦，并且先知没有预言祂将因人民的罪被处死、受辱、被鞭打、被列于罪犯之中、如羊被牵去宰杀（先知说无人能述说祂的世代），那么你们就有理由感到惊讶。但如果这些是祂的特征，并向所有人表明祂，我们除了满怀信心地相信祂，还能做什么呢？凡明白先知著作的人，只要一听到祂被钉十字架，难道不会说这就是祂，而不是别人吗？\par

% structure: p0005 source=h2 target=subsection reason=relative-heading-rank; base=h2

\subsection{第九十章 梅瑟伸开的双手预表十字架}

\Emph{特里丰：}那么，通过经文引导我们，让我们也受你说服；因为我们知道祂应当受苦，被牵如羊。但向我们证明，祂是否必须被钉十字架，并因法律所诅咒的死而如此耻辱和羞辱地死去。因为我们甚至无法让自己想到这一点。\par

\Emph{犹斯定：}你知道先知们所说和所行的事，都是通过比喻和预象掩盖起来的，正如你向我们承认的那样；所以他们所说的大多数事情并不容易为所有人理解，因为他们用这些手段隐藏了真理，以使那些渴望发现和学习它的人能够付出很多努力才能获得。\par

\Emph{特里丰的同伴：}我们承认这一点。\par

\Emph{犹斯定：}那么，请听下面的内容；因为梅瑟首先通过他所做的记号，展示了基督这看似诅咒的事。\par

\Emph{特里丰：}你说的是什么记号？\par

\Emph{犹斯定：}当人民与阿玛肋克作战时，农的儿子，名叫耶稣（若苏厄）的人率领战斗，梅瑟自己向天主祈祷，双手伸开，胡尔和亚郎整日扶着他的手，以免他疲倦时手垂下来。因为如果他放弃这模仿十字架的记号的任何部分，人民就会被打败，正如梅瑟著作所记载的；但如果他保持这种姿势，阿玛肋克就相应地被打败，而获胜者是通过十字架获胜的。因为人民强盛，并不是因为梅瑟如此祈祷，而是因为那位名叫耶稣（若苏厄）的人在前线作战时，他自己做了十字架的记号。因为你们谁不知道，伴随哀哭和眼泪、身体俯伏或屈膝的祈祷最能平息天主的怒气？但无论是他还是任何人，都没有以这种姿势坐在石头上祈祷。甚至石头也没有象征基督，正如我已经表明的。\par

% structure: p0012 source=h2 target=subsection reason=relative-heading-rank; base=h2

\subsection{第九十一章 十字架在若瑟的祝福和举起的蛇中预言}

\Emph{犹斯定：}天主通过梅瑟以另一种方式显示了十字架奥迹的力量，当祂在祝福若瑟时所说的话：\BibleQuote{\Quote{从主的祝福中是祂的土地；因为有天上的时令、甘露、地下的深渊泉源、太阳的季节性果实、月份的聚合、永恒山岳的高处、丘陵的高处、永远流淌的河流、地之肥美的果实；愿那在荆棘丛中显现者所悦纳的临到若瑟的头上和冠冕上。愿他在兄弟中得荣耀；他的美丽如同公牛的头胎；他的角是独角兽的角；他要用这些角从地极到地极推赶万民。}\BibleRef{申 33:13}}现在，没有人能说或证明独角兽的角代表任何其他事实或形象，除了那描绘十字架的预象。因为一根木梁垂直放置，其最高端向上形成一只角，另一根木梁横置于其上，两端像角一样出现在两侧，连接于同一只角。而固定在中央的部分，被钉十字架的人悬挂其上，也像一只角突出来；它看起来也像一只角与其他角连接固定。而\InnerQuote{他用这些角从地极到地极推赶万民}这句话，表明如今在各民族中的事实。因为来自所有民族的一些人，通过此奥迹的力量，被如此推动，即心中被刺透，从虚妄的偶像和魔鬼转向事奉天主。但同样的形象也显明为不信者的毁灭和定罪；正如以色列民出埃及时，透过梅瑟伸手的预象和农的儿子所称呼的耶稣（若苏厄）之名，阿玛肋克被击败，以色列得胜。而为了对抗咬伤以色列的蛇所竖立的预象和记号，似乎是为了拯救那些相信此后藉着被钉十字架的那一位，死亡被宣告临到蛇，而救恩却临到那些被蛇咬伤并投奔那派遣祂的儿子到世界被钉十字架者的人。因为梅瑟的预言之神并没有教导我们相信蛇，因为祂从一开始就向我们表明蛇被天主诅咒；并在依撒意亚中告诉我们，蛇要作为仇敌被大力的剑杀死，这剑就是基督。\par

% structure: p0014 source=h2 target=subsection reason=relative-heading-rank; base=h2

\subsection{第九十二章 除非通过天主的极大恩宠理解圣经，否则天主似乎并非始终教导同一义德}

\Emph{犹斯定：}除非一个人藉着天主的大恩宠获得能力，理解先知们所说和所行的事，否则，若能复述那些言语或事迹的外表，却不能解释其中的道理，对他毫无益处。这些事岂不显然在许多未理解它们的人眼中显得可鄙吗？因为如果有人想问，为何在厄诺客、诺厄和他的儿子们，以及其他类似处境中的人，既未受割损也未守安息日，却中悦了天主，而天主却藉着其他领袖，并在经过许多世代后颁布法律，要求那些生活在亚巴郎和梅瑟之间的人藉割损而成义，并要求那些生活在梅瑟之后的人藉割损和其他规条——即安息日、祭献、奠祭和供物——成义；[天主将受诋毁]，除非你像我已说的那样，证明那预知的天主知道，你的民族将应被驱逐出耶路撒冷，并且无人被允许进入其中。（因为你除了肉身的割损之外，并无其他分别，正如我先前所说的。因为亚巴郎被天主称为义人，并非因割损，而是因信德。因为在他受割损之前，经上有这话论及他：\Quote{亚巴郎信了天主，天主就以此算为他的正义。}\BibleRef{创 15:6} 因此，我们这些在肉身未受割损的人，藉着基督信靠天主，并拥有那对我们这些获得它有益处的割损——即心的割损——我们希望在在天主面前显为正义和蒙悦纳：因为我们已经藉着先知的话，得到祂的证言。）[此外，天主将受诋毁，除非你证明]你们被命令守安息日、献供物，并且主容许有一个被称为天主之名的处所，为的是，如已说的那样，你们不至于因崇拜偶像和忘记天主而不虔敬和无神，正如你们似乎始终显明的那样。（现在，天主为此原因命令安息日和供物的规条，我在先前的话中已证明；但为了今天来的人，我愿几乎重述全部。）因为若非如此，天主将受诋毁，被视为无预知，且未教导所有人认识和行同样的正义行为（因为在梅瑟之前似乎已有许多世代的人类）；并且圣经是不真实的，它肯定：\Quote{天主是真实和正义的，祂的一切道路都是判断，在祂内没有不义。}但既然圣经是真实的，天主始终愿意像你们这样，不愚昧也不自爱，以便你们能获得基督的救恩，祂中悦了天主，并从天主那里得到证言，正如我已说的那样，引用圣先知预言的话作为证明。\par

% structure: p0016 source=h2 target=subsection reason=relative-heading-rank; base=h2

\subsection{第九十三章：同样的正义赋予众人。基督将其总括为两条诫命}

\Emph{犹斯定：}因为[天主]将永远且普遍公正的事，以及一切正义，摆在每种人类种族面前；每个种族都知道奸淫、淫乱、杀人以及类似的事是罪恶；尽管他们都行这些事，但他们并非不知道，每当他们如此行时，他们行事不义，除非那些被不洁之神附身，并被教育、恶习和罪恶制度所败坏，且丧失（或更确切地说，熄灭和压制了）其自然观念的人。因为我们可以看到，这些人不愿承受他们加于别人的同样的事，并以敌对的良心彼此指责他们所犯的行为。因此，我认为我们的主和救主耶稣基督说得好，祂将一切正义和虔诚总括为两条诫命。它们是：\Quote{你当全心、全意、全力爱主你的天主，并爱近人如你自己。}\BibleRef{玛 22:37} 因为那全心、全力爱天主的人，充满敬畏天主之心，就不会尊敬别的神；并且既然天主愿意如此，他就会尊敬那同一位主和天主所爱的天使。而那爱近人如己的人，会愿他得到自己所愿的同样好事，没有人会愿自己遭恶事。因此，爱近人的人会祈祷和努力，使近人拥有与自己同样的益处。人的近人，无非是同样有情感有理性的存在——人。所以，既然一切正义分为两支，即关乎天主和关乎人，那么，凡按照圣经所说，全心、全力爱主天主并爱近人如己的人，就真是义人。但你们从未显明对天主、对先知或对自己有友谊或爱，反倒显而易见，你们常被发现是崇拜偶像者和义人的凶手，以致你们甚至下手对付基督；直到今天，你们仍固守你们的邪恶，咒骂那些证明这位被你们钉死的人就是基督的人。不仅如此，你们还认为祂是因敌对天主且受天主咒诅而被钉死，这想法是你们最不合理心智的产物。因为你们本有方法从梅瑟所给的记号中明白此人是基督，却不愿意；此外，你们以为我们毫无论据，便随意提出你们脑中想到的问题，而你们自己一遇到某些坚定持守的基督徒，就找不出论据了。\par

% structure: p0018 source=h2 target=subsection reason=relative-heading-rank; base=h2

\subsection{第九十四章：挂在树上的人为何受咒诅}

\Emph{犹斯定：} 请告诉我，难道不是天主藉梅瑟命令不可制造任何天上、地上任何物件的形象或肖像，却又使梅瑟在旷野里铸造了铜蛇，立它作为一个记号，使被蛇咬的人因仰望而得救？然而，天主自己并不为不义。因为藉此，正如我先前所说，祂揭示了那奥秘——祂宣告要击碎那导致亚当犯命的蛇的权势，并使那些因这记号而相信祂的人（这记号所预表的是那位将要被钉十字架的），从蛇的毒牙——即恶行、拜偶像和其他不义之事——中获得救恩。除非如此理解，请给我一个理由，为何梅瑟立起铜蛇作为记号，并吩咐被咬的人仰望它，受伤者便得痊愈；况且梅瑟自己曾命令不可制造任何肖像。\par

\Emph{另一位在第二天来的人：} 你说得对，我们无法给出理由。因为我多次就此事询问过老师们，他们中没有一人给我理由。因此，请继续你的讲述；我们正留心听你揭示这奥秘，正因这奥秘，先知们的教导受到不实的诽谤。\par

\Emph{犹斯定：} 正如天主命令藉铜蛇制造记号，而祂自己是无可指责的；同样，虽然法律中记载了对被钉十字架之人的诅咒，但天主的基督身上并无诅咒，相反，凡犯了应受诅咒之事的人都藉着祂而得救。\BibleRef{迦 3:13}\par

% structure: p0022 source=h2 target=subsection reason=relative-heading-rank; base=h2

\subsection{基督亲自担当了我们应得的诅咒}

\Emph{犹斯定：} 因为整个人类都将被发现处于诅咒之下。梅瑟法律中记载：\BibleQuote{凡不持守写在法律书上的一切事去遵行的，是可咒骂的。}\BibleRef{申 27:26} 没有人能完全遵行一切，你们也不敢否认这一点；只是有些人比别人更严格遵守所颁布的诫命。但如果那些在律法之下的人因未遵守一切诫命而显为受诅咒，那么，所有外邦人，他们拜偶像、引诱青年、犯其他罪行，岂不更显为受诅咒吗？那么，如果万有之父愿意祂的基督为整个人类家庭担当所有人的诅咒，祂知道在祂被钉死之后，祂必复活祂，你们为何要争论祂呢？祂是顺从父的旨意而忍受这些事的，你们却仿佛祂是受诅咒的，为何不为你们自己哀哭呢？因为虽然父使祂为人类家庭忍受这些，但你们的行为并非出于顺从天主的旨意。你们杀害先知时，并没有行虔诚。你们中不可有人说：如果祂父愿意祂如此受苦，为使人因祂的创伤而得愈合，我们就没做错。如果你们确实悔改自己的罪，承认祂是基督，并遵守祂的诫命，那么你们可以这样说；因为正如我先前所说，罪过得赦免将属于你们。但如果你们诅咒祂和相信祂的人，并在有权势时杀害他们，那么，作为不义和犯罪的人，心硬无知，你们怎可免于被要求交账，因为你们下手加害于祂？\par

% structure: p0024 source=h2 target=subsection reason=relative-heading-rank; base=h2

\subsection{那诅咒是对犹太人将行之事的预言}

\Emph{犹斯定：} 因为法律上的那句话：\BibleQuote{凡挂在树上的，是可咒骂的。}\BibleRef{申 21:23} 确认了我们对被钉十字架的基督的盼望，不是因为被钉者是天主所诅咒的，而是因为天主预言了你们和像你们这样的人将要行的事——你们不认识祂是那个先于万有存在者，是永恒的天主司铎、君王和基督。你们清楚看到这事已经应验。因为你们在会堂里诅咒一切因祂而被称为基督徒的人；其他民族也实际执行这诅咒，处死那些仅仅承认自己是基督徒的人。我们对所有人说：你们是我们的弟兄；但不如承认天主的真理。虽然他们和你们都不被我们说服，反而极力迫使我们否认基督之名，但我们宁愿选择并接受死亡，确信天主通过基督所应许的一切美善，必将赏赐给我们。此外，我们还为你们祈祷，愿基督怜悯你们。因为祂教导我们也要为仇敌祈祷，说：\BibleQuote{爱你们的仇敌，你们要仁慈，如同你们的天父那样仁慈。}\BibleRef{路 6:35} 因为我们看到全能的天主仁慈宽厚，使祂的太阳升起，光照忘恩者和义人；降雨给圣洁者和恶人。祂教导我们，祂将审判他们所有人。\par

% structure: p0026 source=h2 target=subsection reason=relative-heading-rank; base=h2

\subsection{关于基督十字架的其他预言}

\Person{Justin}：先知梅瑟在胡尔和亚郎扶持他双手时，直到黄昏都保持这种姿态，这并非没有意义。因为主确实在木架上几乎直到黄昏，他们在傍晚埋葬了祂；然后第三天祂复活了。达味这样宣告说：\BibleQuote{我扬声呼求上主，祂从祂的圣山上应允了我。我躺下，睡去；我醒了，因为上主扶持了我。}依撒意亚也论及祂将如何死亡说：\BibleQuote{我向一个悖逆顶嘴的百姓伸开我的双手，他们行走不善之道。}又说祂将复活，依撒意亚亲自说：\BibleQuote{祂的坟茔被从中间移去，我要使富人赔偿祂的死。}\BibleRef{依 53:9} 此外，达味在第二十一篇圣咏中又以奥秘比喻论及苦难和十字架说：\BibleQuote{他们刺透了我的手和脚；他们数算了我的所有骨骸。他们注视并观看我；他们分了我的衣裳，为我的长衣拈阄。}因为当他们钉祂在十字架上时，钉入钉子，刺透了祂的手和脚；那些钉祂的人分了祂的衣裳，各自拈阄取其所欲，并按阄的结果领取。而你们坚持这篇圣咏不是指基督；因为你们在各方面都是瞎眼的，不明白在你们的民族中，凡被称为君王或基督的人，从未在活着时被刺透手脚，或像这样神秘地——即被钉十字架——死亡，唯独这位耶稣除外。\par

% structure: p0028 source=h2 target=subsection reason=relative-heading-rank; base=h2

\subsection{第九十八章　基督在圣咏第二十二篇中的预言}

\Person{Justin}：我要重复整篇圣咏，为使你们听到祂对圣父的敬畏，以及祂如何将一切归给祂，并祈求祂救祂脱离此死；同时在这圣咏中宣告那些起来反对祂的人是谁，并表明祂确实成为了能受苦的人。内容如下：\par

\BibleQuote{我的天主，我的天主，求祢垂听我：祢为何舍弃了我？我的得救远离我，我犯罪的言语。我的天主，我白日呼求祢，祢不俯听；夜间呼求，我并非无知。但祢，以色列的赞颂，居住在圣所中。我们的祖先信赖祢，他们信赖，祢便解救他们。他们呼求祢，得了解救；他们信赖祢，没有蒙羞。但我是一条虫，不是个人；是人的耻辱，百姓的轻蔑。凡看见我的人都嗤笑我；他们撇着嘴，摇着头说：他信赖上主，上主救他吧！上主喜爱他，就拯救他吧！祢曾使我从母腹中诞生，使我依靠祢的胸怀；我从母腹就被投靠祢；祢是我从母胎中就有的天主。求祢不要远离我，因为苦难临近，无人帮助。许多公牛围绕我，巴珊的公牛包围我。他们向我张开大口，如同咆哮的狮子。我如水被倒出，我的骨头都脱了节。我的心像蜡在我腹中熔化。我的精力枯干如瓦片，我的舌头贴在我的上颚；祢将我置于死亡的尘土中。因为许多狗围绕我，恶党包围我。他们刺透了我的手和我的脚，我数算了我的所有骨头。他们注视并观看我；他们分了我的衣裳，为我的长衣拈阄。但是上主，求祢不要远离我，快来扶助我；救我的灵魂脱离刀剑，救我的独生子脱离狗爪。救我脱离狮子的口，并使我由野牛的角得到卑微。我要向我的弟兄宣扬祢的名，在教会中我要赞美祢。敬畏上主的人，请赞美祂！雅各伯的后裔，请光荣祂！以色列的后裔，请敬畏祂！}\par

% structure: p0031 source=h2 target=subsection reason=relative-heading-rank; base=h2

\subsection{第九十九章　圣咏开头是基督临终的话}

\Person{Justin}：现在我要向你们证明，整篇圣咏是这样关乎基督的，我再用我将解释的话来说明。开头所说的\BibleQuote{我的天主，我的天主，祢为何舍弃了我？}从一开始就预告了在基督时期要说的话。因为当祂被钉十字架时，祂说：\BibleQuote{我的天主，我的天主，祢为何舍弃了我？}接着的话：\BibleQuote{我的得救远离我，我犯罪的言语。我的天主，我白日呼求祢，祢不俯听；夜间呼求，我并非无知。}这些以及祂将要作的事都已预告。因为在祂要被钉十字架的那一天，祂带了三位门徒到那称为橄榄山的山上，这山位于耶路撒冷圣殿对面，祂这样祈祷说：\BibleQuote{父啊！如果可以，请叫这杯离开我。}\BibleRef{玛 26:39} 祂又祈祷说：\BibleQuote{不要照我的意愿，而要照祢的意愿。}以此表明祂确实成了能受苦的人。但免得有人说祂那时不知道祂必须受苦，祂随即在圣咏中加上：\BibleQuote{我并非无知。}正如天主问亚当他在哪里，或问加音亚伯在哪里时并没有无知；而是为了说服各人是什么样的人，并为了让我们借着圣经的记录知道一切：同样，基督声明无知不在祂这一边，而在他们那一边，他们以为祂不是基督，幻想他们能杀死祂，并以为祂像某个凡人一样会留在阴府。\par

% structure: p0033 source=h2 target=subsection reason=relative-heading-rank; base=h2

\subsection{第一百章　基督在何种意义上被称为雅各伯、以色列和人子}

\Person{儒斯定}：接着下文——\BibleQuote{但是祢，以色列的赞美，居住在圣所中}——宣告祂要行堪受赞美与惊叹的事，即将在被钉十字架后第三天从死者中复活；这是祂从圣父所获得的。因为我已经表明基督被称为雅各伯和以色列；并且我已经证明，不仅是在若瑟和犹大的祝福中奥秘地宣告了关于祂的事，而且在福音中也记载祂说：\BibleQuote{一切由我父交付于我。}又说：\BibleQuote{除了圣子，没有人认识圣父；除了圣父，也没有人认识圣子，以及圣子所愿意启示的人。}\BibleRef{玛 11:27}因此，祂藉着恩宠，从圣经中向我们启示了我们所领悟的一切，使我们认识祂是天主的首生子，在万物之先；同样，祂也是族长的儿子，因为祂藉着出自他们家族的童贞女取了肉身，并屈尊降生成一个无美貌、受轻慢、受苦的人。因此，在祂的话语中，祂谈到将来的苦难时说：\BibleQuote{人子必须受许多苦，被法利塞人和经师弃绝，被钉十字架，并在第三日复活。}\BibleRef{玛 16:21}祂这样说祂是人子，或是因为祂由童贞女出生，而童贞女正如我所说，出自达味、雅各伯、依撒格和亚巴郎的家族；或是因为亚当是祂和最初列举的那些人的父亲，玛利亚就是从他们而出的。因为我们知道，女人的父亲也是她们女儿所生之子的父亲。因为基督称祂的一位门徒——以前名叫西满——为伯多禄，因为他藉着圣父的启示认出祂是基督、天主子；并且既然我们在宗徒的回忆录中记载祂是天主子，我们也称祂为圣子，我们明白祂在万物之先由圣父藉着祂的能力和意志而生（因为先知们的著作以各种方式称祂为智慧、日子、东升、刀剑、石头、棍杖、雅各伯和以色列）；并且祂由童贞女降生成人，为了使由蛇而来的悖逆在其起源的同一方式中受到毁灭。因为厄娃，本是童贞女和无玷的，怀了蛇的话，就生出了悖逆和死亡。但是童贞玛利亚接受了信德和喜乐，当加俾额尔天使向她报告喜讯，说主的神要降临在她身上，至高者的能力要荫庇她；因此她所生的圣者是天主子；她回答说：\BibleQuote{愿照你的话成就于我吧。}\BibleRef{路 1:38}藉着她，祂出生了，我们已经证明许多圣经指涉祂，并且天主藉着祂毁灭蛇以及那些像它的天使和人类，却为那些悔改自己的邪恶并相信祂的人带来从死亡中的拯救。\par

% structure: p0035 source=h2 target=subsection reason=relative-heading-rank; base=h2

\subsection{101、基督将一切归于圣父}

\Person{儒斯定}：圣咏接着的话是这样，祂说：\BibleQuote{我们的祖先信靠了祢；他们信靠，祢就拯救了他们。他们呼求祢，没有蒙羞。但我是一条虫，不是一个人；受人羞辱，被百姓轻看。}这表明祂承认他们是祂的祖先，他们信靠天主并蒙祂拯救，他们也是童贞女的祖先，藉着她祂出生并降生成人；祂预言祂将蒙同一天主拯救，但不夸耀凭自己的意志或能力完成任何事。因为当祂在地上时，祂也是这样行事，祂回答那个称呼祂为\Quote{善师}的人说：\BibleQuote{你为什么称我为善？只有一位是善的，我天上的父。}但当祂说：\BibleQuote{我是一条虫，不是一个人；受人羞辱，被百姓轻看}时，祂预言了那些确实存在和发生在祂身上的事。因为我们这些信靠祂的人到处受羞辱，\Quote{被百姓轻看}；因为，被你们的民族弃绝和侮辱，祂承受了你们所计划反对祂的那些侮辱。下文说：\BibleQuote{凡看见我的人都嗤笑我；他们嘴唇说话，摇头说：他倚靠上主，上主就救他吧！既然他喜悦他。}这也预言了将发生在祂身上的事。因为那些看见祂被钉十字架的人，每个人都摇头，撇着嘴，彼此扭着鼻子，嘲笑着说那些在祂宗徒回忆录中记载的话：\BibleQuote{祂说祂是天主子；让他下来吧！让天主拯救他吧。}\par

% structure: p0037 source=h2 target=subsection reason=relative-heading-rank; base=h2

\subsection{102、关于基督出生时所发生事件的预言。为何天主允许它}

\Person{Justin}:接下来的内容——\Quote{从我母亲乳房的希望。我在母胎中，就投靠了祢；我自幼就依靠祢，祢是我从母腹中就有的天主：因为没有援助者。许多牛犊围困了我，肥壮的公牛环绕了我。他们向我张口，如同凶残吼叫的狮子。我的骨头都脱了节，像水一样流散。我的心变成了蜡，在我腹中熔化。我的力量干枯如瓦片，我的舌头紧贴在我的喉咙}——预言了将要发生的事；因为那句话：\Quote{从我母亲乳房的希望}（是这样解释的）。祂一在白冷出生，正如我先前所说，黑落德王从阿拉伯贤士那里得知祂的事，就设谋杀害祂，而因天主的命令，若瑟带着祂和玛利亚起身去了埃及。因为圣父已定意，祂所生的那一位应被处死，但要在祂长大成人并宣扬那出自祂的言语之后。但如果你们中有人对我们说：\Quote{天主何不直接处死黑落德？}我预先回答：\Quote{天主何不从一开始就消灭那蛇，使它不存在，反而说：\BibleQuote{我要把仇恨放在他同女人之间，以及他的后裔同她的后裔之间？}\BibleRef{创 3:15}祂岂不能立即造出许多人？然而，既然祂知道这是好的，祂创造了天使和人类，赋予他们自由去行正义，并指定了时期，祂知道在这时期内给他们操练自由意志是好的；又因祂同样知道这是好的，祂设立了普遍和特殊的审判；但每一个人的自由意志却得以保全。}因此圣经在塔的毁灭，语言的分散和改变时说：\BibleQuote{上主说：看，这百姓是一个民族，都说同样的语言；他们已经开始这样做；现在他们企图做的，没有一样能阻挠他们。}\BibleRef{创 11:6}而那句话：\Quote{我的力量干枯如瓦片，我的舌头紧贴在我的喉咙}，也是预言祂将按圣父的旨意所行的事。因为祂坚强的话的能力——祂常用这能力驳倒法利塞人和经师，简言之，一切盘问祂的你们民族的教师——曾有如一股充沛强劲的泉源，其水源被截断那样止息，当祂保持沉默，选择在比拉多面前不回答任何人时；正如祂的宗徒们的回忆录所宣告的，为使依撒意亚所记载的能结出有效的果实，那里写着：\BibleQuote{上主赐给我一个舌头，使我知道何时应该说。}\BibleRef{依 50:4}再者，当祂说：\Quote{祢是我的天主；求祢不要远离我，}祂教导众人应当寄望于创造万物的天主，并只向祂寻求救恩和帮助；而不像其他人那样认为救恩可以靠出身、财富、力量或智慧获得。而你们的行径一向如此：你们曾造了牛犊，并且一直显得忘恩负义，杀害义人，并以自己的出身自豪。因为如果天主子明说，祂得救［既不］因为祂是儿子，也不因为祂强壮或智慧，而是没有天主祂不能得救，即使祂是无罪的，正如依撒意亚用言语所表明的，甚至祂在言语上也没有犯罪（因为祂口里没有不义或诡诈），那么你们或那些不以此望德而期望得救的人，怎能自以为没有自欺呢？\par

% structure: p0039 source=h2 target=subsection reason=relative-heading-rank; base=h2

\subsection{第一〇三章 法利塞人是公牛：咆哮的狮子是黑落德或魔鬼}

\Person{犹斯定}：那么，圣咏接下来所说的——\BibleQuote{因为苦难近了，无人帮助我。许多公牛犊围绕我，肥壮的公牛四面围困我。他们向我张口，如同抓食吼叫的狮子。我的骨头都脱节，像水一样泼散}——同样预知了发生在他身上的事件。因为在那夜，你们民族中由法利塞人、经师和教师所差遣的一些人，从橄榄山来攻击他，圣经所称的抵触和过早毁灭的公牛犊围困了他。而\BibleQuote{肥壮的公牛四面围困我}这句话，是他预先论到那些行为与公牛犊相似的人，当他被带到你们的教师面前时。圣经称他们为公牛，因为我们知道公牛是公牛犊的起源。所以，正如公牛是公牛犊的生育者，你们的教师乃是他们的儿女前往橄榄山捉拿他并带到他面前的原因。\BibleQuote{无人帮助我}这句话也指示了所发生的事，因为连一个作为无辜者帮助他的人也没有。\BibleQuote{他们向我张口，如同吼叫的狮子}这句话，指认了当时犹太人的王，被称为黑落德，他是继承那在基督诞生时屠杀白冷所有婴孩的黑落德的人，因为他想象那从阿拉伯来的贤士所说的人必在他们中间，但他不知道那比一切更强者的旨意，他曾命令若瑟和玛利亚带着婴孩往埃及去，并在那里停留，直到再次得到启示返回本国。他们就在那里停留，直到那屠杀白冷婴孩的黑落德死了，阿赫劳斯继承了他。而他在基督来到他父所赐予的十字架救世计划之前就死了。当黑落德接续阿赫劳斯，领受分给他的权柄时，比拉多把被捆绑的耶稣作为礼物送给他；天主预知此事会发生，曾这样说：\BibleQuote{他们将他带到亚述，献给君王。}\BibleRef{欧 10:6} 或者他藉着那向着他吼叫的狮子意指魔鬼；摩西称之为蛇，但在约伯书和匝加利亚书中，他被称为魔鬼，而耶稣称他为\OriginalTerm{撒殚}{Satanas}，表明他因所行的事获得了一个复合的名称。因为在希伯来语和叙利亚语中，\OriginalTerm{叛徒}{Sata}意为叛徒；\OriginalTerm{蛇}{Nas}这个词，按希伯来术语的解释，被称为\Emph{蛇}，两者结合便产生了\OriginalTerm{撒殚}{Satanas}这个词。因为当耶稣从约旦河上来，有声音对他说\BibleQuote{你是我的儿子，我今天生了你}时，据宗徒们的回忆录记载，这魔鬼曾来试探他，甚至对他说：\BibleQuote{你若俯伏朝拜我}；基督回答说：\BibleQuote{退到我后面去，撒殚！你应朝拜主，你的天主，惟独事奉他。}\BibleRef{玛 4:9} 因为他既欺骗了亚当，也希望能对基督施以某种伤害。此外，\BibleQuote{我的骨头都脱节，像水一样泼散；我的心像蜡熔化在肚肠中}这句话，预知了那一夜当人们到橄榄山攻击捉拿他时所发生的事。因为在我所说的由他的宗徒及他们的跟随者所写的回忆录中记载[说]，当他祈祷并说\BibleQuote{若是可能，就让这杯离开我}时，他的汗滴如同血滴下落，\BibleRef{路 22:44} 他的心和骨头都战栗；他的心像蜡在肚肠中熔化：为使我们明白，父愿意他的子为了我们真正经受这样的苦难，而不会说，他既然是天主子，就不感觉那发生在他身上、加于他的事。此外，\BibleQuote{我的精力像陶片一样枯干，我的舌头贴在上颚}这句话，正如我以前所说，预知了那沉默，当时他谴责了你们所有的教师为愚昧，却一句也不回答。\par

% structure: p0041 source=h2 target=subsection reason=relative-heading-rank; base=h2

\subsection{第一〇四章。圣咏预言基督死亡的境况}

\Person{犹斯定}：\BibleQuote{你使我落入死亡的尘土中；因为许多犬类围困了我，恶党围绕我。他们刺穿了我的手和我的脚。他们数点了我的所有骨头。他们观看，注视着我。他们分了我的外衣，为我的内衣拈阄}——正如我以前所说，这是预知了恶党（他称之为犬类和猎手）将定他死罪，并说那些猎捕他的人是聚集起来并不断努力定他罪的。据他宗徒们的回忆录记载，这确实发生了。我已指出，在他被钉十字架后，钉死他的人分了他的外衣。\par

% structure: p0043 source=h2 target=subsection reason=relative-heading-rank; base=h2

\subsection{第一〇五章。圣咏也预言十字架上的苦难及基督在世最后祈祷的主题}

\Person{犹斯定}：\BibleBook{}{圣咏}随后的话——\BibleQuote{但是祢，主，求祢不要从祢那里撤回对我的援助；求祢垂顾来帮助我。求祢从刀剑下拯救我的灵魂，从狗爪下拯救我的独生子；求祢从狮口中拯救我，从独角兽的角下拯救我的卑微}——这也是对所发生事件的启示和预言。因为我已经证明，祂是万物之父的独生子，由祂以独特方式生为圣言和德能，后来经由童贞玛利亚成为人，正如我们从回忆录所学到的。此外，祂将死于钉十字架也以类似方式预言了。因为那段经文\BibleQuote{求祢从刀剑下拯救我的灵魂，从狗爪下拯救我的独生子；求祢从狮口中拯救我，从独角兽的角下拯救我的卑微}指明了祂所受的苦难，即钉十字架。因为\Quote{独角兽的角}我已向你们解释过，只是十字架的预像。而祈求祂的灵魂从刀剑、狮口和狗爪下获救，是祈求无人能攫取祂的灵魂；这样，当我们生命终结时，我们也可向天主恳求同样的请求，祂能阻挡每一个无耻的邪魔，不让祂们带走我们的灵魂。至于灵魂仍然存活，我已从撒乌耳要求那女巫召上撒慕尔的灵魂这一事实向你们证明。同样，所有类似的义人和先知的灵魂都曾落入这等权势的掌控之下，这确实可从那女巫的具体事例中推知。因此，天主也通过祂的圣子教导我们，为了谁似乎成就了这些事，我们应时常努力认真，并在死时祈求我们的灵魂不致落入任何这等权势之手。因为当基督在十字架上交付祂的灵魂时，祂说：\BibleQuote{父啊，我把我的灵魂交在祢手中}\BibleRef{路 23:46}，这也是我从回忆录中学到的。因为祂劝诫祂的门徒要超过法利塞人的生活方式，警告说，如果他们不这样做，他们可以肯定自己不会得救；这些话记载在回忆录中：\BibleQuote{除非你们的义德超过经师和法利塞人的义德，你们决不能进入天国}\BibleRef{玛 5:20}。\par

% structure: p0045 source=h2 target=subsection reason=relative-heading-rank; base=h2

\subsection{第106章 基督的复活在圣咏的结尾得到预言}

\Person{犹斯定}：\BibleBook{}{圣咏}的其余部分表明，祂知道祂的父会赐予祂所祈求的一切，并使祂从死者中复活；祂敦促所有敬畏天主的人赞美祂，因为祂藉着被钉十字架者的奥迹，怜悯了所有信从的人；祂站在祂的弟兄——宗徒们当中（他们曾因祂被钉十字架而逃离祂，后经祂复活，并亲自说服他们，在受难前祂曾向他们提及祂必须受这些苦，且这些事已被先知预先宣布），在和他们一起生活时，向天主歌唱赞美，正如宗徒的回忆录所显明的。话是这样说的：\BibleQuote{我要向我的弟兄宣扬祢的圣名，在聚会中我要赞美祢。你们敬畏主的人，赞美祂；雅各伯的后裔，全体光荣祂！以色列的后裔，全体敬畏祂！}至于说祂将一位宗徒的名字改为伯多禄，并且记载在关于祂的回忆录中，说这事确实发生了，同样，祂也将另外两兄弟——载伯德的儿子们——的名字改为波阿乃尔革，即雷霆之子；这是宣告一个事实：正是祂使雅各伯被称为以色列，使欧瑟亚被称为耶稣（若苏厄），凡从埃及出来而幸存的人，在若苏厄的名下被领入应许给祖先的福地。至于祂要从亚巴郎的后裔中兴起如同星辰，梅瑟早已预告说：\BibleQuote{由雅各伯将出现一颗星，由以色列将兴起一权杖}\BibleRef{户 24:17}；另一部圣经又说：\Quote{看哪，有一个人，祂的名字是东方。}因此，在祂诞生时，天上出现了一颗星，正如祂宗徒的回忆录所记载的，来自阿拉伯的贤士认出这个标记，前来朝拜了祂。\par

% structure: p0047 source=h2 target=subsection reason=relative-heading-rank; base=h2

\subsection{第107章 从约纳的历史教导同一道理}

\Person{犹斯定}: 祂在受十字架酷刑后的第三天必要复活，在回忆录中记载着：你们民族中有些人质问祂说：\BibleQuote{请给我们显示一个征兆。}祂回答他们说：\BibleQuote{邪恶淫乱的世代寻求征兆，除了约纳的征兆外，再没有征兆给他们。} 由于祂说这话是隐晦的，听众应当理解祂是在受十字架酷刑后的第三天要复活。祂也显示你们世代比\Place{尼尼微}城更邪恶更淫乱；因为当\Person{约纳}在大鱼腹中第三天被吐出来后，向\Place{尼尼微}人宣讲：三天后（另一些版本作四十天）他们都要灭亡，于是他们全体（人兽）披上麻衣，以真诚的忏悔从心中远离不义，相信天主对所有离弃邪恶的人是仁慈宽厚的；以致那城的君王自己也与贵族一同披上麻衣，禁食祈祷，从而获得所求，使那城不被倾覆。但当\Person{约纳}忧愁，因为他所宣告的（第四十）第三天城并未倾覆时，天主藉着一株从地上为他突然长出的蓖麻（他坐在其下遮荫，这蓖麻他没有栽种也没有浇灌，却突然长出来给他遮荫），又藉着蓖麻枯萎（\Person{约纳}为此忧愁）的另一安排，天主指责他不该因\Place{尼尼微}城未倾覆而不当恼怒，说：\BibleQuote{你怜惜了那蓖麻，你并没有为它劳苦，也没有使它生长；它一夜发生，一夜枯死。何况\Place{尼尼微}这座大城，其中不能分辨左手右手的有十二万多人，并有许多牲畜，我岂能不爱惜呢？}\par

% structure: p0049 source=h2 target=subsection reason=relative-heading-rank; base=h2

\subsection{第一〇八章 基督的复活没有使犹太人悔改。但他们却派人到全世界去控诉基督}

\Person{犹斯定}: 尽管你们民族中所有人都知道\Person{约纳}生命中的事件，并且基督在你们中间说过要给你们约纳的征兆，劝告你们至少在祂从死者中复活后要悔改自己的恶行，在\Person{天主}面前哀恸如同\Person{尼尼微}人，以免你们的民族和城市被攻取和毁灭——正如它们已被毁灭那样——然而，你们在得知祂从死者中复活后，不仅没有悔改，反而如我先前所说，派遣了精心挑选并受任命的人到全世界去宣扬，说从一位耶稣（一个加里肋亚人的骗子，我们把他钉了十字架）那里兴起了一个不虔敬、无法无天的异端；祂的门徒在夜间把祂从十字架上解下后埋葬的坟墓中偷走了祂，现在藉着宣称祂已从死者中复活并升到天上而欺骗人。此外，你们指控祂教导了那些不虔敬、无法无天、不圣洁的教理——就是你们所提到的那些——以定那些承认祂是基督、是出自\Person{天主}的导师和\Person{圣子}之人的罪。除此之外，即使你们的城市被攻取、土地被蹂躏，你们仍不悔改，反而胆敢咒骂祂和所有信祂的人。然而，我们并不恨你们或那些因你们而对我们有成见的人；我们反而祈祷，愿现在你们所有人都能悔改，并从\Person{天主}——是万物慈悲、忍耐的圣父——获得怜悯。\par

\SourceRecord{Translated by Marcus Dods and George Reith.；Ante-Nicene Fathers；Vol. 1.；Edited by Alexander Roberts, James Donaldson, and A. Cleveland Coxe.；Buffalo, NY: Christian Literature Publishing Co.,；1885.；<http://www.newadvent.org/fathers/01287.htm>.}

% structure: p0001 source=h1 target=section reason=page-title

\section{与特黎丰对话（第109至124章）}

% structure: p0002 source=h2 target=subsection reason=relative-heading-rank; base=h2

\subsection{第一百零九章：外邦人的归化已由米该亚预言}

\Emph{犹斯定：}但为证明外邦人将悔改他们以往在错误生活中所行的邪恶，当他们听到由耶路撒冷出发的宗徒们所宣讲、并通过他们而学到的道理时，请允许我向你们引用十二位小先知之一米该亚的预言中的一句简短陈述。如下：\BibleQuote{到末日，上主的圣殿山必要矗立在群山之上，超乎一切山岳；万民都要向它涌来。将有许多民族前去，说：\InnerQuote{来！我们攀登上主的圣山，往雅各伯天主的殿里去！}他必将他的道路指示给我们，我们也必在他指引的路上行走；因为法律将出自熙雍，上主的话将出自耶路撒冷。他将统治万民，为远方许多强国定断是非；他们必要把刀剑铸成锄头，把枪矛制成镰刀；民族与民族不再持刀相向，人也不再学习战斗；各人坐在自己葡萄树下和无花果树下，无人再扰乱他们，因为万军的上主亲口说了。万民各奉自己神祇之名而行，但我们却永远奉上主我们天主之名而行。到那一天，我要聚集跛足的，集合被放逐的和我所惩治的；我要使跛足的成为遗民，使被放逐的成为强盛的民族；上主将在熙雍山作他们的君王，从今世直到永远。}\par

% structure: p0004 source=h2 target=subsection reason=relative-heading-rank; base=h2

\subsection{第一百一十章：预言的一部分已应在基督徒身上，其余将于基督第二次来临时应验}

\Emph{犹斯定：}现在我知道你们的老师，各位先生，承认这段经文全部指向基督；我也知道他们坚持他还没有来；或者如果说他来了，他们断言不知道他是谁；但当他显现并荣耀时，那时才知道他是谁。然后，他们说，这段经文所提到的事件将会发生，就好像预言的言语还没有结出任何果实一样。哦，无理性的人！不明白所有这些经文所证明的，即宣告了基督的两次来临：一次，他被描述为受苦、无荣耀、受辱、被钉十字架；另一次，他将带着荣耀从天而降，那时那背道之人，说出反对至高者的奇怪言语，将敢于在地上对我们基督徒行非法的事，我们这些从法律和从耶路撒冷藉耶稣的宗徒们所传出的话学习了真天主敬拜的人，已逃往雅各伯的天主和以色列的天主那里避难；我们这些充满战争、互相残杀和一切邪恶的人，在全地上各自将我们的战争武器——刀剑打成犁头，枪矛打成农具——并且我们培养虔敬、正义、仁爱、信德和望德，这些我们是从圣父自己通过那位被钉十字架者得来的；各人坐在自己的葡萄树下，即每人拥有自己已婚的妻子。因为你们知道先知的话说：\BibleQuote{你的妻子好比结实累累的葡萄树。}现在明显的是，没有人能恐吓或征服我们这些在全地相信耶稣的人。因为很明显，尽管被斩首、被钉十字架、被抛给野兽、锁链、火刑和所有其他折磨，我们并不放弃我们的宣认；相反，这种事发生得越多，其他人就越多人成为信徒，并藉耶稣的名字敬拜天主。因为正如有人砍掉葡萄树的结果部分，它会重新生长，并长出其他茂盛结果的枝条；同样的事也发生在我们身上。因为天主和救主基督所栽种的葡萄树就是他的子民。但预言的其余部分将在他的第二次来临时应验。因为\Quote{那受压迫的[并被驱逐的]}的表达，即从世界[暗示]，就你们和所有其他人所能做到的而言，每个基督徒不仅被驱逐出自己的财产，甚至被驱逐出整个世界；因为你们不允许任何基督徒活着。但你们说同样的命运也落到了你们自己的民族身上。现在，如果你们在战败后被驱逐，你们确实遭受了这样的待遇，正如所有圣经所见证的；但我们，在认识天主的真理之后，并没有做过这样的[邪恶]行为，却被天主见证，我们与那最正义、唯一无玷和无罪的基督一同被带走离开大地。因为依撒意亚呼喊：\BibleQuote{看，义人怎样丧亡，没有人放在心上；虔诚的人被收去，没有人注意；义人被收去，是脱离了灾祸。}\BibleRef{依 57:1}\par

% structure: p0006 source=h2 target=subsection reason=relative-heading-rank; base=h2

\subsection{第一百一十一章：两只山羊预示两次来临；有关第一次来临的其他预像，外邦人藉基督的血获得释放}

\Emph{儒斯定：}就是以象征所宣告的，在梅瑟时代就已经预示了这位基督的两次来临，正如我之前提到的，这从在禁食期间作为祭献的公山羊的象征中显明了。再者，通过梅瑟和若苏厄所做之事，同样的事也以象征的方式被宣告和预先告诉了我们。因为其中一个（梅瑟）伸开双手，在山上一直待到傍晚，他的双手被支撑着；这除了十字架的预像之外，不预示任何其他事物。而另一个，他的名字被改为耶稣（若苏厄），带领战斗，以色列人得胜了。这事发生在两位圣人和天主的先知身上，好叫你明白，他们中哪一位也不能同时承担两个奥秘：我指的是十字架的预像和名字的预像。因为唯独祂的力量是、曾是、并将是，祂的名字使一切权势战栗，因为它们因被祂毁灭而极其痛苦。因此，我们受苦并被钉十字架的基督并非被法律诅咒，而是显明唯独祂要拯救那些不背离祂信德的人。逾越节的血涂在每个人的门框和门楣上，拯救了在埃及得救的人，当时埃及人的长子被击杀。因为逾越节就是基督，祂后来被祭献，正如依撒意亚所说的：\BibleQuote{祂被牵去宰杀，如同羊羔。}\BibleRef{依 53:7}并且经上记载，在逾越节的日子你们逮捕了祂，也在逾越节期间把祂钉在十字架上。正如逾越节的血拯救了在埃及的人，同样，基督的血也将拯救那些相信的人脱离死亡。那么，如果这个记号没有在门上，天主会被欺骗吗？我不是这个意思；而是肯定祂预先宣告了人类通过基督的血将要得到的救恩。因为朱红线绳的记号，是由农的儿子若苏厄派往耶里哥的探子给妓女辣哈布的，告诉她把线绳系在窗户上，她是从那里放他们下去躲避敌人的，这也显示了基督的血的预像，通过这血，那些从前在各民族中曾是妓女和不义的人得救了，获得了罪的赦免，并且不再继续犯罪。\par

% structure: p0008 source=h2 target=subsection reason=relative-heading-rank; base=h2

\subsection{第一一二章 犹太人贫弱地解释这些预象，只注意琐事}

\Emph{儒斯定：}但是你们，以低下（和属世）的方式解释这些事，归咎于天主很大的软弱，如果你们只是这样听着，而不探究所说言语的力量。因为就连梅瑟也会因此被视为违犯者：他命令不可制造天上、地上或海里的任何形状的东西；然后他自己却造了一条铜蛇，挂在竿上，吩咐被咬的人去看它；他们一看就得了救。那么，那条蛇（我已经说过，天主在起初就诅咒了它，并用巨大的剑将它斩断，正如依撒意亚所说的，\BibleRef{依 27:1}）当时被理解为拯救了百姓吗？我们要按你们教师愚蠢的理解接受这些事，而不视为预象吗？我们岂不更应该将这竿与被钉十字架的耶稣的相似之处相联系，因为梅瑟伸开双手，连同那被称为耶稣（若苏厄）的人，为你们的人民取得了胜利吗？因为这样，我们就不会对立法者所做的事感到困惑了——他没有离弃天主，却说服百姓寄望于一个野兽，而正是借着这个野兽，犯罪与不服从才有了开端。这事是由那位有福的先知以极大的智慧和奥秘说出来的；所有先知所说的和所做的，没有一件是可以合理指责的，只要人拥有在他们里面的知识。但如果你们的教师只向你们解释，为什么在这个段落中说母骆驼，而不在那个段落；或者为什么在供物中使用这么多细面和这么多油；并且是以低下卑劣的方式做这些，而他们却从不敢讲论或解释那些伟大且值得探究的要点，或命令你们不要听我们解释这些要点，不要与我们交谈，难道他们不该听到我们的主耶稣基督对他们说的话：\BibleQuote{粉饰的坟墓，外面好看，里面却装满死人的骨头；你们将薄荷献上十分之一，却吞下了骆驼；瞎眼的向导！}因此，如果你们不藐视那些自高自大、希望被称为辣彼、辣彼之人的教义，并以如此的诚恳和智慧来面对预言的话语，如同从你们自己的人民那里遭受与先知们同样的苦难，你们就不可能从先知著作中获得任何益处。\par

% structure: p0010 source=h2 target=subsection reason=relative-heading-rank; base=h2

\subsection{第一一三章 若苏厄是基督的预像}

\Emph{耶稣：}我的意思是：耶稣（若苏厄），正如我多次指出的，当他被派遣去侦探客纳罕地时，梅瑟给他起名叫耶稣（若苏厄）。他为何这样做，你们既不询问，也不感到困惑，亦不严格查考。因此，基督从你们的眼前溜走了；你们虽阅读，却不明白；即使现在，你们听见耶稣就是我们的基督，你们也不认为这名号是无目的或偶然赋予他的。但你们却就亚巴郎的原名上为何添加了一个\Quote{\OriginalTerm{\GreekText{α}}{\GreekText{α}}}而进行神学讨论；至于撒辣的名字上为何添加了一个\Quote{\OriginalTerm{\GreekText{ρ}}{\GreekText{ρ}}}，你们也使用类似的高调辩论。但你们为何不类似地查考农（努）的儿子欧舍亚、他父亲所给的名字，为何被改为耶稣（若苏厄）呢？但不仅他的名字被改了，而且他被立为梅瑟的继承人，是他同时代从埃及出来的人中唯一一个，他带领残存的人民进入圣地；并且是他，而非梅瑟，带领百姓进入圣地，并将那地分给与他一同进入的人为产业；同样，耶稣基督也将再次聚集百姓的分散，并将美地分给各人，尽管方式不同。因为前者给了他们暂时的产业，因为他既不是天主的基督，也不是天主子；但后者，在神圣的复活之后，将赐给我们永恒的产业。前者，在他被命名耶稣（若苏厄）之后，并从他（圣神）那里领受了能力，命令太阳停住。因为我已证明，是耶稣基督向梅瑟、亚巴郎以及所有其他圣祖显现并交谈，履行圣父的旨意；他后来由童贞玛利亚降生为人，并永远活着。因为后者是那一位，圣父将藉着他并为了他更新天地；他是在耶路撒冷永远发光的光；他是按照默基瑟德的品位为撒冷王的，是至高者永恒的大司祭。前者被说成曾用石刀第二次给百姓行割损（这乃是耶稣基督亲自用石刀给我们行割损的标记，即从石头及其它材料所造的偶像中割损我们），并在各处用石刀将那些受过割损的人从未受割损的人中，即从世界的错误中，聚集起来；这里所说的石刀，就是主耶稣基督的话。因为我已证明，基督被先知在比喻中称为磐石和岩石。因此，石刀我们理解为他的话，藉着这些话，许多在错误中的人从无割损中被割损，即心得割损；这是天主藉着耶稣从那时期开始命令那些从亚巴郎领受割损的人受割损的，说耶稣（若苏厄）要用石刀第二次给那些进入那圣地的人行割损。\par

% structure: p0012 source=h2 target=subsection reason=relative-heading-rank; base=h2

\subsection{第114章。一些辨别关于基督的言论的规则。犹太人的割损与基督徒所领受的割损大不相同。}

\Emph{犹斯定：}因为圣神有时使某些作为未来预象的事被清楚实行；有时祂说出关于将要发生的事的话，如同当时正在发生或已经发生。除非阅读的人领悟这技巧，否则就不能适当地追随先知的话。为举例，我将重复一些先知段落，使你可以理解我所说的。当祂藉依撒意亚说：\BibleQuote{他被牵去待宰，如同一只羊，又如羔羊在剪毛者前无声}\BibleRef{依 53:7}，祂说话如同那苦难已经发生。当祂又说：\BibleQuote{我张开了双手，迎接一个违抗和反叛的民族}\BibleRef{依 65:2}；以及当祂说：\BibleQuote{主，谁相信了我们的报道？}\BibleRef{依 53:1}——这些话如同在宣布已经发生的事件。因为我已经指出，基督常被在比喻中称为磐石，在比喻性的言语中称为雅各伯和以色列。再者，当祂说：\BibleQuote{我要仰观诸天，你手指的工程}，除非我明白祂使用言语的方法，我就不能通晓地理解，反而是像你们的经师们一样，认为万有的圣父、无始的天主有手、脚、手指和灵魂，如同一个复合体；因此他们教导说，是圣父自己显现给亚巴郎和雅各伯。因此我们是有福的，我们已用石刀第二次受了割损。因为你们第一次的割损是用铁器实行的，你们仍然心硬；但我们的割损是第二次，在你们的之后设立，用尖锐的石头，即由不用手凿的角石的宗徒所宣讲的话，将我们从偶像崇拜和一切邪恶中割损。我们的心就这样从邪恶中被割损，以致我们乐于为那善磐石的名字而死，这磐石使活水涌出，赐给那些藉着他爱慕万有圣父的人心，并赐给那些愿意饮用生命之水的人。但我说这些事，你们却不了解；因为你们不明白所预言基督要做的事，也不相信我们这些提醒你们注意所记载之事的人。因为耶肋米亚这样呼喊：\BibleQuote{祸哉！因为你们离弃了活泉，却为自己挖了破裂的水池，不能存水。难道因为我在你们眼前给了耶路撒冷休书，熙雍山就要变为旷野吗？}\BibleRef{耶 2:13}\par

% structure: p0014 source=h2 target=subsection reason=relative-heading-rank; base=h2

\subsection{第115章。关于基督徒在匝加利亚中的预言。犹太人在辩论中的恶毒方式。}

\Person{尤斯定}：但你们应当相信匝加利亚，他以比喻显示基督的奥秘，并隐秘地宣告它。以下是他的话：\BibleQuote{熙雍女子，欢乐喜乐罢！因为，看，我来，我要住在你们中间，上主说。在那一天，许多民族要归附上主。他们要作我的子民，我要住在你们中间；你们要知道万军的上主派遣了我到你们这里来。上主必在圣地上承受犹大作为他的产业，他必再拣选耶路撒冷。凡有血肉的，在上主面前都要战栗，因为他从他的圣云中奋起。他指给我看大司铎耶稣（Joshua）站在上主的天使面前；撒殚站在他右边控告他。上主对撒殚说：拣选耶路撒冷的上主斥责你。看，这不是从火中抽出来的一根柴吗？}\BibleRef{匝 2:10}\par

\Person{特黎丰}正欲回答反驳我时，我说：\par

请稍待，先听我说：因为我将要给出的解释并非你们所想的那样，好像巴比伦地方（你们的民族曾在那里为囚徒）没有一个名叫耶稣（Joshua）的司铎。但即使我这样说了，我也已经指出：如果你们民族中有个名叫耶稣（Joshua）的司铎，先知在他的神视中仍未亲眼看见他，正如他也并未亲眼看见撒殚或上主的天使一样，而是在神魂超拔中，当启示临到他时。但我现在说：正如圣经说农的儿子（Nun）以耶稣（Joshua）这名字行了大事和奇事，预先宣告了我们的主所要成就的事；同样，我现在继续指出：在巴比伦你们民族中，在司铎耶稣（Joshua）的日子里，所赐下的启示，是宣告了我们的司铎——他是天主，是基督，是万有之父天主的圣子——所要完成的事。\par

事实上，我感到惊奇，为何刚才你保持沉默，当我说话时，并且为何你没有打断我，当我说农的儿子（Nun）是同时代出埃及的人中唯一与那被描述为比那一代年轻的人一同进入圣地的人时。因为你们像苍蝇一样成群扑上伤口。因为即使有人说了千万句好话，若有一句小话令你们不悦，因它不够明了或准确，你们就不顾许多好话，只抓住那点小话，并热切地将其定为不敬和该死的。为的是，当你们受天主同样审判时，你们要为你们的极大狂妄——无论是恶行还是你们歪曲真理所得的恶劣解释——承担更重的罪责。因为你们用什么审判审判人，你们也理当受同样的审判。\par

% structure: p0019 source=h2 target=subsection reason=relative-heading-rank; base=h2

\subsection{说明此预言如何适用于基督徒}

\Person{尤斯定}：但为使你们明了神圣的耶稣基督的启示，我重新开始我的论述，并断言那启示正是为我们这些相信基督——这位大司铎，即被钉十字架者——而赐下的；我们虽曾生活在淫乱和各种污秽的交谈中，却因我们耶稣的恩宠，照祂父的旨意，脱去了我们身上所沾染的所有那些污秽的邪恶。虽然撒殚常在我们身边抵抗我们，并渴望引诱所有人归向他，但上主的天使，即通过耶稣基督赐予我们的天主的德能，斥责他，他便离开我们。我们便如同从火中抽出，当从我们过去的罪中被净化，并从磨难和火炼（撒殚及其同谋以此考验我们）中被拯救；天主子耶稣应许再次从中拯救我们，并赐给我们预备好的衣服，如果我们遵守祂的诫命；祂并承诺为我们预备永恒的国度。因为正如那被先知称为司铎的耶稣（Joshua），显然穿着污秽的衣服，因为据说他娶了一个妓女为妻，并被称为从火中抽出的柴，因为当抵抗他的撒殚被斥责时，他获得了罪的赦免；同样，我们这些因耶稣之名，同心合意相信万有的创造者天主的人，因祂首生圣子的名，脱去了污秽的衣服，即我们的罪；并被祂召叫的话语所强烈点燃，我们是天主真正的大司铎的种族，正如天主自己作证说，在外邦人中，各处都有人向祂呈献悦纳而纯洁的祭献。但天主除了通过祂的司铎，不接收任何人的祭献。\par

% structure: p0021 source=h2 target=subsection reason=relative-heading-rank; base=h2

\subsection{玛拉基亚关于基督徒祭献的预言。不能理解为指散居犹太人的祈祷。}

\Emph{\Person{犹斯定}:} 因此，天主预见了我们通过这名所献的一切祭献，即\Person{耶稣基督}命令我们奉献的，就是在饼与杯的圣体圣事中的祭献，这是基督徒在普天下各处所呈献的，祂见证这些祭献是祂所悦纳的。但祂完全拒绝了你们和你们那些司铎所呈献的祭献，说：\BibleQuote{我不从你们手中收纳你们的祭品；因为从日出到日落，我的名在万民中受尊崇（祂说）；但你们却亵渎了它。}\BibleRef{拉 1:10} 然而，即使在现在，你们由于好辩，竟断言天主不接受那时住在耶路撒冷、被称为以色列人的那些人所献的祭献；反而说祂悦纳那散居的国民中个人的祈祷，并称他们的祈祷为祭献。我也承认，祈祷和感恩由配得的人奉献时，是天主唯一完美而悦纳的祭献。因为基督徒只承担奉献这样的祭献，并通过他们在固体和液体食物中所纪念的，藉此记起圣子所受的苦难；你们民族的大司祭和经师使祂的名在普天下受亵渎和咒骂。但你们以\Person{耶稣}之名所归给所有成为基督徒的人的那些污秽衣裳，天主宣告将要从我们身上除去，那时祂要使众人从死者中复活，并指定一些人在那永恒不朽的王国中成为不朽、不死、无忧的；而将另一些人送去受永不熄灭的火的刑罚。至于你们和你们的经师们，你们曲解圣经的话，以为是论到当时散居的你们民族，并声称他们在各处所献的祈祷和祭献是纯洁而悦纳的，要知道你们在说谎，并千方百计欺骗自己：因为首先，即使现在你们的民族也没有从日出延伸到日落，有些民族中从来没有你们种族的人居住过。因为没有一个种族的人，无论是蛮族、希腊人，或无论叫什么名字的游牧者、漂泊者、住帐篷的牧人，其中不通过被钉死的\Person{耶稣}之名献上祈祷和感恩。而且，正如圣经所示，玛拉基亚写下这些话的时候，你们如今遍及全地的流散还没有发生。\par

% structure: p0023 source=h2 target=subsection reason=relative-heading-rank; base=h2

\subsection{第118章 劝人在基督来临前悔改；基督徒因信基督，远比犹太人更为虔诚}

\Emph{\Person{犹斯定}:} 所以你们更应当放弃好辩之心，在大审判的日子来到之前悔改；那时，你们支派中所有刺透这\Person{基督}的人都要哀哭，正如我已表明圣经所预言的。我也解释了主所起的誓，\BibleQuote{按照默基瑟德的品位}，以及这预言的意义；至于\Person{依撒意亚}的预言说：\BibleQuote{他的坟墓被夺去了，}\BibleRef{依 53:8} 我已经说过，是指\Person{基督}将来的埋葬和复活；并且我也多次指出，这\Person{基督}就是所有生者和死者的审判者。\Person{纳堂}对\Person{达味}论及祂时，也接着说：\BibleQuote{我要作祂的父亲，祂要作我的儿子；我的仁慈必不离开祂，如我离开祂以前的人一样；我要在祂的家中和祂的王国中永远建立祂。}\BibleRef{撒下 7:14} \Person{厄则克耳}说：\BibleQuote{在祂的殿中，除了祂以外，不再有别的王子。}\BibleRef{则 44:3} 因为祂是被选的司祭和永远的君王，就是\Person{基督}，祂是天主之子；不要以为\Person{依撒意亚}或其他先知预言在祂第二次来临时，会在祭台上献上血祭或奠祭，而是指真正的、属灵的赞美和感恩。我们信靠祂并非徒然，也没有被那些教导我们这些教理的人引入歧途；这一切都是由于天主奇妙的预知，为使我们因着新而永久的盟约的召唤，即\Person{基督}的召唤，被证明比你们更有智慧和更敬畏天主；你们自以为爱天主和有智慧，其实不然。\Person{依撒意亚}惊叹于此，说：\BibleQuote{君王们都要掩口无言：因为那些未曾听见有关祂的事的人，将要看见；那些未曾听过的人，将要明白。主，谁相信了我们的报道？上主的手臂向谁显示了呢？}\BibleRef{依 52:15} \Person{特黎丰}，我重复这些，在允许的范围内，为了今天同你一起来的人的缘故，只是简短扼要地尽力而为。\par

\Emph{\Person{特黎丰}:} 你做得很好；虽然你重复同样的事情颇长，但请放心，我和我的同伴都很乐意听。\par

% structure: p0026 source=h2 target=subsection reason=relative-heading-rank; base=h2

\subsection{第119章 基督徒是应许给亚巴郎的圣洁子民。他们像亚巴郎一样蒙召}

\Person{犹斯定}: 诸位先生，你们以为，如果我们没有因祂所喜悦的那位的旨意而蒙受分辨之恩，我们岂能明白圣经中这些事吗？为了使梅瑟的话应验：\Quote{他们以异[神]激起我的嫉妒，以他们的可憎之物惹我发怒。他们祭拜了素不认识的魔鬼，是新近出现的神祇，是你们祖先未曾知道的。你遗弃了生你的天主，忘记了养你的天主。主看见了，便起了嫉妒，因祂儿女的怒气而发怒，说：我要转面不顾他们，指示他们终将遭遇何事；因为这实在是一个乖僻世代，没有信德的儿女。他们以那不是神的，激起我的嫉妒，以他们的偶像惹我发怒；我也要以那不是子民的，激起他们的嫉妒，以愚昧的国民惹他们的怒气。因为在我怒中有火燃起，烧至阴府，吞灭大地及其出产，烧着山岳的基础；我要将灾祸加在他们身上。}\BibleRef{申 32:16} 在那位义者被处死以后，我们作为另一个子民兴旺起来，像新而繁盛的麦穗发出；正如先知所说：\Quote{在那一天，许多民族要归附主，成为祂的子民，并要住在全地之中。}\BibleRef{匝 2:11} 但是我们不仅是一个子民，而且是一个圣洁的子民，正如我们已经指出的。\Quote{他们要称他们为圣洁的子民，是主救赎的。}\BibleRef{依 62:12} 因此，我们不是一个可轻视的子民，也不是野蛮的种族，如同卡里亚人和弗里吉亚人那样；反而天主拣选了我们，并向那些不寻求祂的人显现了。祂说：\Quote{看啊，我对那不呼求我名的民族说，我是天主。}\BibleRef{依 65:1} 因为这正是天主古时向亚巴郎所应许的那个民族，祂曾宣称要使他成为多民之父；但这并非指阿剌伯人、埃及人或厄东人，因为依市玛耳成了一个强大民族的父亲，厄撒乌也是如此；并且如今阿孟人人数众多。此外，诺厄是亚巴郎的父亲，事实上也是全人类的父亲；而其他人是其他人的祖先。那么，基督赐给亚巴郎的更大恩宠是什么？就是：祂以相似的召唤用祂的声音呼叫他，吩咐他离开所居住的地方。祂也用那声音召唤了我们，我们已经离开了过去的生活方式，不再像地上其他居民那样在邪恶中度日；并且我们与亚巴郎一同承受圣地，当我们在无尽的永恒中获得产业时，因着同样的信德，我们成为亚巴郎的子女。因为他相信天主的声音，这就算为他的正义；同样，我们相信天主藉基督的宗徒所讲、又藉先知传给我们的话，甚至至死弃绝了世上的一切。因此，祂应许给他一个有同样信德的民族，敬畏天主、正义、并使父喜悦；但你们却不是，\Quote{你们是没有信德的。}\par

% structure: p0028 source=h2 target=subsection reason=relative-heading-rank; base=h2

\subsection{第120章 基督徒曾被应许给\Person{依撒格}、\Person{雅各伯}和\Person{犹大}}

\Person{犹斯定}：请注意，天主向依撒格和雅各伯所作的许诺也是同样的。因为祂这样对依撒格说：\Quote{地上万民都要因你的后裔获得祝福。}\BibleRef{创 26:4} 对雅各伯说：\Quote{地上万族都要因你和你的后裔获得祝福。}\BibleRef{创 28:14} 祂没有说这是给厄撒乌或勒乌本，也不是给任何人，只给了那些基督将按天主旨意、藉着童贞玛利亚而由之兴起的人。但如果你考虑犹大的祝福，你就会明白我所说的话。因为后裔从雅各伯分开，经由犹大、培勒兹、叶瑟和大卫而来。这正是一个象征，表示你们民族中有些人将是亚巴郎的子女，并且也将在基督的份中；而其他人，虽然确实是亚巴郎的子女，却像海边的沙，又多又无数，却不结果实，只喝海水。你们民族中有一大群人被证明属于这一类，他们吸收了苦涩和无神的道理，却拒绝了天主的圣言。因此，祂在关于犹大的段落中说：\Quote{权杖不离犹大，司权者不离他双股之间，直到那为他所存留的来到；祂将是万民的期望。}\BibleRef{创 49:10} 显然，这话不是指犹大，而是指基督。因为我们万民所期盼的，不是犹大，而是耶稣，祂曾带领你们的祖先出离埃及。因为预言指的是基督的来临：\Quote{直到那为祂存留的来到，祂将是万民的期望。}因此，耶稣来了，正如我们已详细说明的，并且祂还要再次显现在云之上；你们亵渎祂的名字，并竭力使祂的名字在全地被亵渎。诸位，我本可与你们争论关于你们所解释的读法，说经上写着\Quote{直到那为祂存留的事来到}，尽管七十贤士译本并非如此解释，而是\Quote{直到那为祂存留的来到}。但由于下文表明是指基督（因为它是\Quote{祂将是万民的期望}），我不与你们作纯文字上的争论，正如我未曾试图从你们不承认的圣经章节中，即我从耶肋米亚先知、厄斯德拉和大卫的话中所引用的，来证明基督；而是从你们至今仍承认的章节中，而你们如果你们的老师明白，我确信他们早就删掉了，就像他们删掉关于依撒意亚的死一样，你们曾用木锯把他锯成两半。这正是一个奥秘的预表，表明基督将要你们民族一分为二，并使那些配得尊荣的人与圣祖和先知一起被提升到永远的国度；但祂说，祂将把其他人送到不灭之火的惩罚中，与所有民族中类似不顺服和不悔改的人一起。\Quote{因为}祂说，\Quote{他们要来自西方和东方，与亚巴郎、依撒格和雅各伯一起坐席在天国里；但国度的子民将被赶到外面的黑暗中。}我说这些，除了讲真理之外，毫无顾忌，也不受任何人强迫，即使我立刻被你们撕碎。因为我写信给凯撒时，绝没有考虑我本族（即撒玛黎雅人）的任何意见，而是声明他们错了，因为他们信从本族的术士西满，他们说西满是超越一切权力、权威和能力的天主。\par

% structure: p0030 source=h2 target=subsection reason=relative-heading-rank; base=h2

\subsection{第一二一章 从外邦人相信耶稣的事实，证明祂就是基督}

他们保持沉默，我便接着说：\par

\Person{犹斯定}：朋友们，[圣经]通过大卫论及这位基督时，不再说\Quote{因他的后裔}万民要受祝福，而是\Quote{因他}。所以这里说：\Quote{祂的名将万古常新，超越太阳；万民都要因祂蒙受祝福。}但如果万民都在基督内受祝福，而我们这些万民都信祂，那么祂确实就是基督，而我们就是那些因祂受祝福的人。天主从前赐下太阳作为崇拜的对象，正如经上所记，但从未有人见过因信太阳而忍受死亡；然而为了耶稣之名，你可以看到来自万族的人忍受并继续忍受一切痛苦，而不否认祂。因为祂的真理和智慧之言比太阳的光芒更炽热、更光明，渗入人心和思想的深处。因此圣经也说：\Quote{祂的名将超越太阳。}此外，匝加利亚说：\Quote{祂的名字是\InnerQuote{东方}。}\BibleRef{匝 6:12} 论到同一位，他说：\Quote{各支派都要哀悼。}\BibleRef{匝 12:12} 但如果祂在第一次来临（那是无尊威、无美貌且极其卑微的）时已经如此照耀、如此大能，以致没有一个民族不知道祂，并且各地的人都已悔改自己旧有的邪恶生活方式，以至于邪魔都屈服于祂的名字，一切权势和国度都畏惧祂的名超过畏惧所有死人，那么，当祂在荣耀的来临时，岂不必要彻底消灭所有恨祂、不义地背离祂的人，而赐安息给祂自己的人，以他们所期待的一切作为赏报吗？因此，我们得以听见、明白并藉着这位基督得救，并认识由圣父启示的一切[真理]。所以祂对祂说：\Quote{祢称为我的仆人，为复兴雅各伯的支派，领回以色列的遗民，还是小事；我还要使祢成为万民的光，使祢的救恩直达地极。}\BibleRef{依 49:6}\par

% structure: p0033 source=h2 target=subsection reason=relative-heading-rank; base=h2

\subsection{第一二二章 犹太人无理地将此理解为关于归依者}

\Person{犹斯定}：你们以为这些话是指外方人和归依者说的，但实际上是指我们这些被耶稣光照的人。因为基督甚至也曾为他们作证；但如今你们却成了加倍的地狱之子，正如祂自己所说的。\BibleRef{玛 23:15} 因此，先知所记载的，不是论及那些人，而是论及我们，经上关于我们说：\InnerQuote{我要领瞎子走他们所不知道的路，使他们走他们未曾认识的路径。我是见证，上主天主说，并是我的仆人所拣选的。}\BibleRef{依 42:16} 那么，基督为谁作证呢？显然是为那些相信的人。但归依者不仅不信，反而比你们加倍地亵渎祂的名，并企图折磨和杀害我们这些信祂的人；因他们在各方面都要与你们相似。在另一处祂又呼喊道：\InnerQuote{我，上主，凭公义召叫了祢，我必握住祢的手，保护祢，立祢作人民的盟约，外邦人的光明，为开启盲人的眼睛，从牢狱中领出囚犯。}\BibleRef{依 42:6} 各位先生，这些话确实也是论及基督，并关系到被光照的外邦人；或者你们又要说，祂是向他们论及法律和归依者？\par

这时，那些第二天来的人中有人好像身在剧场般喊叫起来：\par

\Signature{特黎丰的同伴}：但这难道不是指法律和那些被法律光照的人吗？这些就是归依者啊。\par

\Person{犹斯定}：（看着特黎丰）不是的，因为如果法律能光照外邦人和拥有它的人，那还需要新盟约做什么？但既然天主预先宣告祂要订立一个新盟约，一条永久的法律和诫命，这些就不能理解为旧的法律和它的归依者，而是指基督和祂的归依者，即我们这些外邦人，祂光照了我们，正如祂在某处说：\InnerQuote{在悦纳的时候，我应允了祢；在救恩的日子，我帮助了祢；我立了祢作人民的盟约，为复兴土地，为继承荒芜之地。}\BibleRef{依 49:8} 那么，基督的产业是什么？岂不是万民？天主的盟约是什么？岂不是基督？正如祂在另一处说：\InnerQuote{祢是我的儿子，我今日生祢。祢向我请求，我必将万民赐祢作产业，将地极赐祢作基业。}\par

% structure: p0038 source=h2 target=subsection reason=relative-heading-rank; base=h2

\subsection{\Emph{第一二三章　犹太人荒谬的诠释　基督徒才是真以色列}}

\Person{犹斯定}：因此，既然所有这些后来的预言都指着基督和万民，你们也该相信前面的预言同样是指祂和他们的。因为归依者不需要盟约，既然有一条同样的法律加于所有受割礼的人身上，经上论及他们这样说：\InnerQuote{外方人也要与他们联合，并要与雅各伯家联合。}\BibleRef{依 14:1} 并且，归依者受了割礼以便加入百姓，就成为像他们一样的人；而我们虽被称为一个百姓，却是外邦人，因为我们没有受割礼。此外，你们若以为归依者的眼睛要被开启而你们自己的却不，你们被视为瞎子和聋子而他们却被光照，那是可笑的。如果你们说法律已赐给万民，而你们却不知道，那就更显荒谬。因为你们本该敬畏天主的义怒，不至于成为无法无天、飘荡的儿子；你们本该非常害怕常常听见天主说：\InnerQuote{没有信德的儿子。谁眼瞎，不是我的仆人呢？谁耳聋，不是统治他们的人呢？天主的仆人也被弄瞎了。你们看了又看，却没有观察；耳朵开了，却没有听见。}天主对你们的称赞是光荣的吗？天主的见证适合祂的仆人吗？你们虽然屡次听见这些话，却不羞愧。你们不因天主的威胁而战栗，因为你们是愚顽硬心的百姓。\InnerQuote{因此，看哪，我要继续对待这百姓，}上主说，\InnerQuote{我要除掉他们，并除掉智慧人的智慧，隐藏明达人的明达。}\BibleRef{依 29:14} 这也是当得的：因为你们既不智慧也不明达，反而狡猾诡诈；只善于作恶，却全然不能认识天主隐藏的旨意，或主的忠信盟约，或找出永恒的道路。\InnerQuote{因此，上主说，我要在以色列和犹大播种人种和兽种。}\BibleRef{耶 31:27} 至于另一个以色列，祂藉依撒意亚这样说：\InnerQuote{到那日，在亚述和埃及之间必有一个第三以色列，在万军的上主所祝福的地上蒙福，说：我的百姓在埃及和亚述是有福的，以色列是我的产业。}既然天主祝福这百姓，称他们为以色列，并宣布他们是祂的产业，你们怎不懊悔你们对自己的欺骗（好像只有你们是以色列），并咒骂天主所祝福的百姓呢？因为祂论及耶路撒冷和其四邻时，这样补充说：\InnerQuote{我要使人在你身上增多，就是我的百姓以色列；他们要继承你，你必成为他们的产业；你不再使他们丧失子女。}\BibleRef{则 36:12}\par

\Person{特黎丰}：那么，你们是以色列吗？祂这些话是对你们说的吗？\par

\Person{犹斯定}：如果我们没有就这些主题进行过长时间的讨论，我可能会怀疑你们是否出于无知而问这个问题；但既然我们已经通过论证达成结论并得到你们的同意，我不相信你们对我刚说的话不了解，或只是想再次挑起争论，而是你们要我在这些人面前展示同样的证明。\par

看到他眼神中同意的表情，我便继续说道：\par

\Emph{Justin:} 再者，在\BibleBook{依撒意亚}{}中，如果你们有耳可听，天主以比喻谈论基督，称他为\Person{雅各伯}和\Person{以色列}。他这样说：\BibleQuote{雅各伯是我的仆人，我必扶持他；以色列是我所拣选的，我必将我的圣神赐予他，他必将审判带给外邦人。他不争辩，不呼喊，也没有人听见他在街市上的声音。压伤的芦苇，他不折断；将熄的灯心，他不吹灭；但他必将审判引向真理。他要发光，不被折断，直到他在世上设立了审判。外邦人将因他的名而信赖。}\BibleRef{依 42:1} 因此，正如从一个人\Person{雅各伯}（又名以色列）起，你们整个民族被称为雅各伯和以色列；同样，我们来自基督，他生我们归于天主，像\Person{雅各伯}、\Person{以色列}、\Person{犹大}、\Person{若瑟}、\Person{达味}一样，被称为并确实是天主的真儿子，并且遵守基督的诫命。\par

% structure: p0044 source=h2 target=subsection reason=relative-heading-rank; base=h2

\subsection{第124章 基督徒是天主的儿子}

当我看到他们因为我称我们是天主的儿子而困扰时，我预料到他们会发问：\par

\Emph{Justin:} 先生们，请听圣神如何论及这民族，说他们都是至高者的儿子；并且这位基督会亲临他们的集会，审判所有的人。这话是达味说的，按照你们的版本，是这样：\BibleQuote{天主站在诸神的会中；他在诸神中施行审判。你们审判不义要到几时？你们偏袒恶人要到几时？应为孤儿和穷人伸冤，为谦卑和贫乏的人施行正义。解救贫乏的人，从恶人手中拯救穷人。他们无知，也不明白，在黑暗中行走；大地的根基都要动摇。我曾说：你们是神，都是至高者的儿子。但你们要像世人一样死亡，像众首领之一一样跌倒。天主啊，求你起来审判大地，因为你将万国归你所有。}\BibleRef{咏 82} 但在七十贤士译本中，写的是：\BibleQuote{看，你们要像世人一样死亡，像众首领之一一样跌倒。}\BibleRef{咏 82} 为要显明人的悖逆——我指的是亚当和厄娃——以及众首领之一的坠落，即那被称为蛇的，因他欺骗了厄娃而遭受巨大倾覆。但我的论述不打算触及这一点，而是要向你们证明，圣神责备人，因为人原本像天主一样，没有痛苦和死亡，如果他们遵守了他的诫命，并配得上称为他的儿子，但他们却变得像亚当和厄娃，为自己招致死亡。至于这篇圣咏的解释，你们可以按自己意愿理解，但无论如何，它表明所有人都配成为\Quote{神}，并有能力成为至高者的儿子；并且每个人都将像亚当和厄娃一样被审判和定罪。现在我已经详细证明基督被称为天主。\par

\SourceRecord{Translated by Marcus Dods and George Reith.；Ante-Nicene Fathers；Vol. 1.；Edited by Alexander Roberts, James Donaldson, and A. Cleveland Coxe.；Buffalo, NY: Christian Literature Publishing Co.,；1885.；<http://www.newadvent.org/fathers/01288.htm>.}

% structure: p0001 source=h1 target=section reason=page-title

\section{\WorkTitle{与特黎丰对话（第一二五至一四二章）}}

% structure: p0002 source=h2 target=subsection reason=relative-heading-rank; base=h2

\subsection{第一二五章 他解释以色列一词的力量，以及它如何适用于基督}

\Person{Justin}: 各位先生，我希望从你们那里学习以色列这个名字的力量。\par

他们沉默时，我继续说道：\par

\Person{Justin}: 我知道的事，就告诉你们；因为我认为，知道而不说，是不对的；或者，我怀疑你们明明知道，却出于嫉妒或自愿的无知而自欺，却还要不断挂虑，也是不对的。但我单纯而坦诚地说出一切，正如我的主所说的：\BibleQuote{有一个撒种的人出去撒种；有的落在路旁，有的落在荆棘中，有的落在石头上，有的落在好地上。}\BibleRef{玛 13:3}\newline{}所以我必须说话，希望在某个地方找到好地；因为我的这位主，作为强壮有力者，来向所有人讨回自己的东西，如果他知道主人有能力并且来讨回自己的东西，而他将东西交给每一个银行，没有因任何原因挖掘，就不会斥责他的管家。因此，以色列这个名字的意义是：一个胜过能力的人；因为\OriginalTerm{胜者}{Isra}是胜过的人，而\OriginalTerm{能力}{El}是能力。而且，基督在成为人时如此行，已由雅各伯与向他显现的那位摔跤的奥秘所预示，因为他侍奉圣父的旨意，然而他仍然是天主，因为他是所有受造物的首生者。因为他成为人时，正如我之前所说，魔鬼来到他面前——即被称为蛇和撒殚的那个能力——试探他，并试图通过要求他朝拜自己来使他跌倒。但他摧毁并战胜了魔鬼，证明他是邪恶的，因为他要求被当作天主来朝拜，违反圣经；他是背离天主旨意的叛徒。因为他回答魔鬼说：\BibleQuote{经上记载：你要朝拜主，你的天主，唯独事奉他。}\BibleRef{玛 4:10}\newline{}于是魔鬼被战胜和定罪，那时就离开了。但是，由于我们的基督将要被麻木，即因痛苦和苦难的经历，他通过触摸雅各伯的大腿并使其萎缩预先表明了这一点。但以色列原本是他的名字，当他用自己的名祝福有福的雅各伯时，他将雅各伯的名字改为以色列，由此宣告所有通过他逃往圣父那里的人构成有福的以色列。但你们丝毫不理解这些，也不准备理解，因为你们是按肉身的种子，是雅各伯的子孙，就期望自己必然得救。但你们在这些事上自欺，我已经用许多话证明了。\par

% structure: p0006 source=h2 target=subsection reason=relative-heading-rank; base=h2

\subsection{第一二六章 基督按双重本性的各种名号；显示他是天主，并向圣祖显现}

\Person{Justin}: 但你们，特黎丰，如果知道那位有时被称为大谋略的天使、厄则克耳称为人、达尼尔称为人子、依撒意亚称为婴儿、达味称为基督和应受朝拜的天主、许多人称为基督和石头、撒罗满称为智慧、梅瑟称为若瑟、犹大和星、匝加利亚称为东方、依撒意亚再次称为受苦者、雅各伯和以色列、以及棍杖、花朵、基石、天主之子——你们就不会亵渎现在已经来临、出生、受苦、升天的那一位；他将再次来临，那时你们的十二支派将要哀悼。因为如果你们理解了先知所写的，就不会否认他是天主，是唯一、无始、不可言说的天主的儿子。因为梅瑟在\BibleBook{出谷}{纪}中某处写道：\BibleQuote{主对梅瑟说，对他说，我是主，我曾显现给亚巴郎、依撒格和雅各伯，作他们的天主；但我没有把我的名字启示给他们，我与他们立定了我的盟约。}同样地，他又说：\BibleQuote{一个人与雅各伯摔跤。}并断言那是天主；叙述雅各伯说：\BibleQuote{我面对面见了天主，我的生命得以保存。}而记载称他与那位摔跤、向他显现并祝福他的地方为\OriginalTerm{天主的容貌}{Peniel}。梅瑟也说天主在玛默勒橡树旁显现给亚巴郎，那时他正坐在帐棚门口，在中午。接着他继续说：\BibleQuote{他举目观看，见有三个人站在他面前；他一看见，就从帐棚门口跑去迎接他们。}\BibleRef{创 18:2}\newline{}过了一会儿，其中一人应许给亚巴郎一个儿子：\BibleQuote{撒辣为什么笑，说：我果真能生育吗？我已经老了？在天主岂有难事吗？到了约定的日期，我要回到你这里，按着生命的时候，撒辣要生一个儿子。}他们就离开亚巴郎走了。他又这样提到他们：\BibleQuote{那些人从那里起身，面向索多玛。}\BibleRef{创 18:16}\newline{}接着，对亚巴郎说话的那位，那曾是的和现在仍是的，又说：\BibleQuote{我岂能向亚巴郎，我的仆人，隐瞒我要做的事呢？}\BibleRef{创 18:17}\par

梅瑟著作中的后续内容，我引用并解释了。\par

由此我证明，那被描述为天主、显现给亚巴郎、依撒格、雅各伯和其他圣祖的那一位，是在圣父和主的权下受命的，并事奉他的旨意。\par

然后我继续说以前没有说过的话：\par

\Person{Justin}: 所以，当人民想吃肉，梅瑟对那位（在那里也称为天使）失去信心时，那一位曾应许天主会让他们饱足，他是天主又是天使，由圣父派遣，被描述为说并做这些事。因为圣经这样说：\BibleQuote{主对梅瑟说：难道主的手不够吗？现在你要知道我的话是否隐瞒你。}\BibleRef{户 11:23}\newline{}又，用其他话说：\BibleQuote{但主对我说：你不得过这约但河；主你的天主，他走在你们前面，他必剪除外邦人。}\par

% structure: p0012 source=h2 target=subsection reason=relative-heading-rank; base=h2

\subsection{第一二七章 这些圣经章节不适用于圣父，而是适用于圣言}

\Person{犹斯定}：立法者和先知记载了这些及类似的言论；我想我已经充分说明，无论何处天主说：\Quote{天主由亚巴郎面前离开了}，\BibleRef{创 18:22} 或\Quote{主对梅瑟说}，\BibleRef{出 6:29} 以及\Quote{主下来要看看世人所造的城和塔}，\BibleRef{创 11:5} 或当\Quote{天主把诺厄关在方舟内}时，\BibleRef{创 7:16} 你切勿以为未受生的天主亲自降临或从某处上升。因为那不可言说的父及万有的主，既未到过任何地方，也不行走，也不睡觉，也不起来，而是在祂自己的处所（无论那在何处），以不可思议的大能迅速观看与听闻；祂没有眼目也没有耳朵，却看见一切，知道一切，我们中无人能逃过祂的鉴察；祂在整个世界中不受移动或局限于一隅，因祂在造世之前已存在。那么，祂如何能同任何人谈话，或被任何人看见，或出现在大地极小的部分呢？当时在西乃山上的人连从祂那里派来的那位的光荣都看不透；梅瑟本人也不能进入他所支搭的会幕，因为那会幕充满了天主的荣耀；而撒罗满将约柜运入他在耶路撒冷为约柜所建的殿宇时，司祭也站立不住。因此，亚巴郎、依撒格、雅各伯或任何其他人，都未看见父及万有的不可言说的主，也未看见基督；但他们看见了那按祂旨意成为祂圣子的那一位，祂是天主，又因奉行祂的旨意而称为天使；祂也乐意由童贞女降生为人；祂在荆棘丛中与梅瑟交谈时也是一团火。除非我们这样理解圣经，否则必然得出结论：当梅瑟所记载的事发生时，万有的父和主并不在天上：\Quote{主从天主那里降下硫磺与火，落在索多玛和哈摩辣}，\BibleRef{创 19:24} 又如达味所说：\Quote{城门，请高举你们的门楣；古老的门户，请高抬起来，因为光荣的君王要进来}；以及祂说：\Quote{上主对我主说：你坐在我右边，等我把你的仇敌放在你脚下}。\par

% structure: p0014 source=h2 target=subsection reason=relative-heading-rank; base=h2

\subsection{第一二八章 圣言被派遣，并非如同无生命的权能，而是作为由父的实体所生的位格}

\Person{犹斯定}：基督既是主，又是天主圣子，先前曾以权能显现为人、天使，并如荆棘丛中那样显现于火的荣耀中，也同样在审判索多玛时显现，这已由前面所说的充分证明了。\par

于是我再次重复先前从《出谷纪》所引用的全部内容，即关于荆棘丛中的显现和若苏厄（耶稣）的命名，并继续说：\par

诸位，请不要以为我屡次重复这些话是多此一举；但我知道有些人想抢先反驳，说那从万有之父那里派遣、向梅瑟、亚巴郎或雅各伯显现的权能，之所以称为天使，是因为祂来到人间（因祂将父的命令传达给人）；称为光荣，是因为祂有时以不可承受的异象显现；称为人、人子，是因为祂按父的意愿以这样的形像显现；他们称祂为\Quote{圣言}，因为祂将父的消息带给人类；但他们坚持这权能是与父不可分、不可离的，正如他们所说太阳在地上的光与天上的太阳不可分、不可离一样；当太阳落下时，光也一同落下；所以父，他们说，当祂愿意时，使祂的权能发出，当祂愿意时，又将其收回。他们以此方式教导说，祂创造了天使。但事实证明，有些天使是永远存在的，绝不会再化为那种由之而出的形像。至于这权能，先知的话称之为天主（也已充分证明），且称为天使，并非像太阳的光那样仅在名称上有所区别，而是在数目上确实不同，我在前面已简要讨论过；当时我断言，这权能是由父按其能力和旨意所生，却不是由于分割，好像父的本质被分开一样；如同一切被分割的东西在分割前后并不相同；为举例说明，我取从火中引燃的火为例，我们看到火星与火有区别，而能点燃许多火的原火并不减少，反而保持不变。\par

% structure: p0018 source=h2 target=subsection reason=relative-heading-rank; base=h2

\subsection{第一二九章 从圣经其他段落得到证实}

\Person{犹斯定}：现在我要再次引用我曾为证明此点所说的话。当圣经说：\Quote{主从天主那里降下火}，先知的话表明有两位：一位在地上，经上说祂下来要察看索多玛的呼叫；另一位在天上，祂也是地上之主的\Quote{主}，正如祂是父和天主，是地上之主的能力和祂之为\Quote{主}及\Quote{天主}的缘由。再者，当圣经记载天主在起初说：\Quote{看，亚当成了我们中的一个}，\BibleRef{创 3:22} 这\Quote{我们中的一个}一词也是指数目的事，这些话不容许像那些既不能陈述也不能理解真理的诡辩者所强加于其上的比喻意义。在《智慧篇》中也写着：\Quote{若我要向你们述说日常的事，我将从起初追忆它们。上主在祂造化的开始，在太初就造了我，为祂的工程。从永远，在起初，在造地以前，在深渊产生以前，在泉水涌出以前，在群山奠定以前，在丘陵之前，我已受生。}\par

我重复这些话后，补充说：\par

我的听众，如果你们留心，就会察觉圣经宣告这\Quote{后裔}是父在一切受造物之前所生的；而那所生的与那生者，任何一人都承认在数目上是不同的。\par

% structure: p0022 source=h2 target=subsection reason=relative-heading-rank; base=h2

\subsection{第一三〇章 他回到外邦人的归化，并显示这已预言之}

当所有人都表示赞同后，我说：\par

\Person{尤斯定}：我现在要引用一些先前未曾引述的经文。这些经文记载于忠信的仆人梅瑟的比喻中，如下：\Quote{诸天哪，应当与祂一同欢乐，愿天主的众天使都敬拜祂。}\BibleRef{申 32:43}\par

我又补充了该段接下来的部分：\par

\Quote{外邦人哪，当与祂的百姓一同欢乐，愿天主的众天使在祂里面得坚固；因为祂要报应祂儿子的血，必要报应，并以报应惩罚祂的仇敌，报应那些恨祂的人；上主要洁净祂百姓的地。}藉这些话，祂宣告我们这些外邦人要同祂的百姓——即亚巴郎、依撒格、雅各伯、众先知，以及一切蒙天主喜悦的这百姓——一同欢乐，正如我们先前已达成一致的。但我们将不接受你们整个民族；因为我们从依撒意亚知道\BibleRef{依 66:24}，那些违命之人的肢体必被虫和不灭的火所吞噬，永存不朽；以致他们成了一切血肉的景象。\par

除此之外，诸位，我还想从梅瑟亲口的话中再补充一些经文，好叫你们明白，天主从起初就按宗族和语言分散了所有人；并从所有宗族中选取了你们的宗族——一个无用、悖逆、无信德的世代；并显明那些从万邦中被拣选的人，藉着基督——祂也称他为雅各伯，并起名为以色列——遵行了祂的旨意。这些人，正如我先前详细说过的，必是雅各伯和以色列。因为当祂说\Quote{外邦人哪，当与祂的百姓一同欢乐}时，祂赐给他们同样的产业，却没有以同一名称称呼他们；但当祂说他们作为外邦人与祂的百姓一同欢乐时，祂称他们为外邦人，为要责备你们。因为正如你们藉偶像崇拜惹祂发怒，祂也同样使那些拜偶像的人配得认识祂的旨意，并继承祂的产业。\par

% structure: p0028 source=h2 target=subsection reason=relative-heading-rank; base=h2

\subsection{第一三一章　归向基督的外邦人比犹太人更忠于天主}

\Person{尤斯定}：但我要引用那段揭示天主如何分散万民的经文。如下：\Quote{问你父亲，他必指示你；问你的长者，他们也必告诉你；至高者将万民分开，如同分散亚当子孙时，按以色列子孙的数目立定万民的疆界；上主的份就是祂的百姓雅各伯，以色列是祂产业的阄分。}\par

七十贤士翻译为\Quote{祂按天主天使的数目立定万民的疆界}。但我的论点并不因此有丝毫削弱，所以我采纳你们的解释。你们自己若愿承认真理，就必须承认：我们这些藉着被轻看和可耻的十字架奥迹蒙天主召叫的人（为了宣认此奥迹、服从它，并为了我们的虔敬，我们受魔鬼和魔鬼的军旅借着你们所提供的援助，遭受了甚至致死的刑罚），忍受一切折磨也不肯在言语上否认基督——我们藉着祂蒙召，进入由父预备的救恩——我们比你们更忠于天主。你们曾以强有力的大手和极大的荣耀的眷顾从埃及被救赎出来：海为你们分开，留下干路，在那里天主将追赶你们的人——装备极其精良、战车辉煌——尽都击杀，使为你们而开的路重归大海；又有光柱照耀你们，使你们比世上任何民族更拥有独特、永不熄灭、永不沉没的光；天主为你们降下吗哪作为食粮，适合天上的天使，使你们无须预备食物；玛拉的水变为甘甜；又有被钉十字架之祂的记号，既在咬你们的蛇的事上（我方才已提及，天主在适当时期之前预先预示这些奥迹，为将恩宠赐予你们——你们却永远被证明是忘恩负义的），也在梅瑟伸手和敖舍亚被称为耶稣（若苏厄）的预表中，当你们与阿玛肋克争战时：关于这事，天主命令将此事记载下来，并将耶稣的名字存记在你们心中，说祂是要从天下涂抹阿玛肋克纪念的那一位。如今很显然，阿玛肋克的纪念在农的儿子之后仍然存留；但祂藉着被钉十字架的耶稣显明，那些预象是祂将要经历之事的预告：魔鬼将被毁灭，并且惧怕祂的名；一切执政的、掌权的都要惧怕祂；从万民中相信祂的人要显明为敬畏天主、爱好和平的人。我方才引述的事实，特黎丰，也指明这点。再者，当你们想吃肉时，有极多的鹌鹑赐给你们，多到数不胜数；又有水从磐石中涌出；有云彩跟随你们，遮蔽炎热，抵御寒冷，宣告另一种新天的样式和意义；你们的鞋带没有断，鞋子没有破，衣服也没有穿坏，甚至孩子们的衣裳也随着他们长大。\par

% structure: p0031 source=h2 target=subsection reason=relative-heading-rank; base=h2

\subsection{第一三二章　耶稣的名字在旧约中有何等大能}

\Person{尤斯定}：然而在这之后，你们铸了一只牛犊，并非常热衷与外邦女子行淫，事奉偶像。再者，当土地以如此大的权能赐给你们，以至于你们目睹太阳因那位名叫耶稣（若苏厄）的人的吩咐停在天空中，三十六小时不落下，以及所有其他按时机为你们所行的奇迹时——其中有一件我现在似乎愿意再提一下，因为它有助于使你们认识我们确知是基督、天主圣子的那位耶稣，祂曾被钉十字架、复活、升天，并要再来审判所有人，直到亚当本人。\par

那么，你们知道，当约柜被阿市多得敌人夺去，一种可怕且无法治愈的疾病在他们中间爆发时，他们决定将约柜放在一辆车上，套上刚产犊的母牛，藉此试验他们是否因约柜而受天主能力的打击，以及天主是否愿意将约柜送回原处。他们这样做了之后，无人牵引的母牛没有去约柜被取走的地方，而是到了一个人的田地，那人的名字叫\Person{奥舍亚}，与后来改名为\Person{耶稣（若苏厄）}（如前所述）的人同名，他也带领百姓进入那地，并将地分给他们；母牛到了这些田地后便停在那里，以此向你们表明，它们是受那名字的能力引导；正如从前那些从埃及出来而幸存的人，被那位后来称为\Person{耶稣（若苏厄）}（原名\Person{奥舍亚}）的人领进那地一样。\par

% structure: p0034 source=h2 target=subsection reason=relative-heading-rank; base=h2

\subsection{第一三三章　犹太人的心硬，基督徒为之祈祷}

\Emph{尤斯定：}如今，尽管这些以及所有其他如此出人意料、令人惊奇的事工曾在你们中间、在不同时期施行并被人看见，但先知们证明你们竟至于把自己的儿女献给魔鬼；此外，你们还敢对基督做出这样的事，而且仍然敢做；但愿为这一切，你们能从天主和祂的基督那里获得怜悯与救恩。因为天主预先知道你们会做这样的事，便藉先知依撒意亚向你们宣布了这诅咒：\Quote{祸哉，他们的灵魂！他们为自己策划了恶计，说：让我们捆绑义人，因为他使我们厌烦。因此他们必吃自己行为的果实。祸哉，恶人！恶必按他手所作的临到他。我的百姓啊，你们的压榨者将你们捡拾，剥削你们的人要管辖你们。我的百姓啊，称你们为有福的人使你们走迷，扰乱你们所行的路。但现在上主必协助祂的百姓施行审判，祂必与民间的长老和首领一同进入审判。你们为什么烧毁我的葡萄园？为什么穷人的掠夺物在你们家中？你们为什么欺压我的百姓，使卑微人的脸面蒙羞？}\BibleRef{依 3:9}此外，同一先知又以别的话说了同样的意思：\Quote{祸哉，那些以长绳牵引罪孽，以牛轭的索子牵引过犯的人！他们说：愿祂的速临临近，愿以色列圣者的计谋来临，好让我们知道。祸哉，那些称恶为善、称善为恶的人！那些以暗为光、以光为暗的人！那些以苦为甜、以甜为苦的人！祸哉，那些自以为有智慧、自视为精明的人！祸哉，你们中间那些饮酒的勇士，那些大有能力的壮士，搀和烈酒的人！他们因贿赂而称恶人为义，夺去义人的正义。因此，火炭必烧灭碎秸，火焰必烧尽碎屑，他们的根必如羊毛，他们的花必如尘土上腾。因为他们不愿遵守万军上主的法律，藐视以色列圣者上主的话语。所以万军上主大大发怒，伸手攻击他们，打击他们；祂因山岳而震怒，他们的尸首在路中如同粪土。虽然如此，他们仍不悔改，他们的手仍然高举。}\BibleRef{依 5:18}因为你们的手确实高举而行恶，因为你们杀害了基督，且不为此悔改；不但如此，你们还恨恶并杀害我们这些藉着祂而相信万有的天主和父的人，凡你们能做的，你们就做；你们不停地咒骂祂，也咒骂那些站在祂一边的人；但我们众人却为你们和所有人祈祷，正如我们的基督和主教导我们那样，祂命令我们甚至要为仇敌祈祷，爱那恨我们的人，为那咒骂我们的人祝福。\par

% structure: p0036 source=h2 target=subsection reason=relative-heading-rank; base=h2

\subsection{第一三四章　雅各伯的婚姻是教会的预像}

\Person{儒斯定}: 如果先知和祂自己的教训打动了你们，那么跟随天主比跟随你们那些轻率盲目的师傅更好，他们至今仍允许每人娶四五个妻子；如果有人看见美貌女子想要她，他们就引用雅各伯（称为以色列）和其他祖宗的作为，声称做这种事不算错；因为他们在这事上极度无知。因为，如我先前所说，这类行为中每一件都成就了某种重大奥秘的措置。因为在雅各伯的婚姻中，我要提一提成就了何种措置和预言，好使你们由此知道你们的师傅从未察看促使每件行为的神圣动机，只看那卑下腐败的情欲。因此请注意我所说的话。雅各伯的婚姻是基督将成就之事的预像。因为雅各伯同时娶两姐妹是不合法的。他为拉班的一个女儿服事；但在得到幼女的事上被欺骗，他就再服事七年。现在，肋阿是你们的民族和会堂；辣黑耳却是我们的教会。为这两种人及其中的仆人，基督至今仍在服事。因为当诺厄把第三子的后裔作为仆人赐给两个儿子时，如今反过来，基督来恢复他们中的自由子民和仆人，将所有遵守他诫命的人都赐予同样的尊荣。正如雅各伯所生的自由妇人和婢女的儿子都是儿子，地位同等。并且预先告知了每个人按等级和预知应成为什么。雅各伯为有斑点和许多斑点的羊服事拉班；基督也为人类各种各样、形态众多的种族服事，甚至至于十字架的奴役，藉十字架的血和奥迹赢得他们。肋阿视力弱；因为你们灵魂的眼睛极其软弱。辣黑耳偷了拉班的神像，并藏起来直到今日；我们却失去了我们的父神和物质神。雅各伯一直被他兄弟憎恨；我们现在和我们主自身也被你们和所有人憎恨，尽管我们本性上是兄弟。雅各伯被称为以色列；以色列已被证实为基督，他就是，并被称为耶稣。\par

% structure: p0038 source=h2 target=subsection reason=relative-heading-rank; base=h2

\subsection{第一三五章  基督是以色列的王，基督徒是以色列民族}

\Person{儒斯定}: 当圣经说：\BibleQuote{我是上主天主，以色列的圣者，使以色列的王为你所知，} \BibleRef{依 43:15} 你们岂不明白基督确实是永生的王吗？因为你们知道，依撒格的儿子雅各伯从未作王。因此圣经再次向我们解释说，雅各伯和以色列所指的王是谁：\BibleQuote{雅各伯是我的仆人，我必扶持祂；以色列是我的被选者，我的灵魂必接纳祂。我已将我的圣神赐给祂；祂必将审判带给外邦人。祂不呼喊，祂的声音不在外头被人听见。压伤的芦苇，祂不折断；将熄的灯芯，祂不吹灭，直到祂使审判得胜。祂要照耀，不被折断，直到祂在地上设立审判。外邦人必因祂的名而信赖。} \BibleRef{依 42:1} 那么，是先祖雅各伯，外邦人和你们将信赖他吗？还是基督呢？因此，正如基督是以色列和雅各伯，同样，我们这些从基督的腹中采掘出来的人，才是真正的以色列民族。但让我们更留意这句话：\BibleQuote{祂说：我要从雅各伯中领出后裔，从犹大中领出：他们必承受我的圣山；我的被选者和我的仆人必承受产业，住在那里；在灌木丛中必有羊群的圈，阿苛尔谷必成为寻求我之人的牛群安息之所。至于你们，那离弃我、忘记我的圣山、为魔鬼摆设筵席、为魔鬼斟满奠酒的，我要将你们交与刀剑。你们必倒在刀下；因为我呼唤，你们不听从，在我眼前行恶，拣选我所不喜悦的。} \BibleRef{依 65:9} 圣经的话就是这样；因此要明白，这里所说的雅各伯的后裔是别的，并非如可能认为的那样，是指你们民族。因为雅各伯的后裔不可能为雅各伯的子孙留出入口，天主也不可能接纳祂曾责备为不配承受产业的那些人，又应许给他们；但正如那里先知所说：\BibleQuote{雅各伯家啊，来！让我们在上主的光明中行走；因为祂已遣散祂的子民雅各伯家，因为他们的地，如同起初，充满了占卜和卜卦；} 同样，我们必须在这里观察到有两个犹大的后裔和两个种族，正如有两个雅各伯家：一个由血肉而生，另一个由信德和圣神而生。\par

% structure: p0040 source=h2 target=subsection reason=relative-heading-rank; base=h2

\subsection{第一三六章  犹太人因拒绝基督，拒绝了派遣祂的天主}

\Person{Justin}：因为你看祂现在如何对百姓说话，稍前说：\BibleQuote{如葡萄在串中找到，他们要说：不要毁坏它，因其中有祝福；我必因我仆人的缘故如此行：为他的缘故，我不将他们都灭绝。}此后祂又说：\BibleQuote{我必从雅各伯中领出后裔，从犹大中领出承受我山的人。}那么显然，若祂如此向他们发怒，并威胁只留下他们中极少的人，祂应许要领出另一些人，他们将住在祂的山上。但这些人正是祂所说要播种并生养的人。因为你们既不容祂召唤你们，也不听祂对你们说话，反在天主面前行了恶。但你们邪恶的极点在于：你们恨那义者，并杀死了祂；并且这样对待那些从祂领受了一切所有、且虔诚、正义、仁爱的人。因此，主说：\BibleQuote{祸哉，他们的灵魂！}\BibleRef{依 3:9}\BibleQuote{因为他们对自己图谋了恶计，说：\InnerQuote{让我们除掉义人，因为他令我们厌恶。}}的确，你们不惯于向巴耳献祭，如你们的祖先那样，也不在林中和高处为天上万军摆供饼；但你们没有接受天主的基督。因为不认识祂的人，就不认识天主的旨意；而侮辱和仇恨祂的人，就是侮辱和仇恨那派遣祂的。凡不信祂的，也不信先知们的宣告，那些先知曾宣讲并传报祂给众人。\par

% structure: p0042 source=h2 target=subsection reason=relative-heading-rank; base=h2

\subsection{第一三七章 他劝勉犹太人悔改}

\Person{Justin}：兄弟们，不要对被钉十字架者说恶言，不要轻蔑那使众人得医治的鞭伤，正如我们得医治一样。因为若被圣经说服，你们从心硬中受割损，那将是好的：不是那出于灌输给你们的教条的割损；因为那割损是作为记号赐下的，并非为正义的行为，正如圣经迫使你们［承认］的。因此，要同意，不要嘲笑天主子；不要听从法利塞人的教师，不要讥讽以色列的君王，一如你们会堂的领袖在祈祷后教导你们所做的：因为若触摸那不为天主喜悦的人\BibleRef{匝 2:8}，就如同触摸祂眼中的瞳人，何况触摸祂的爱者呢！而这就是祂，已经充分证明了。\par

他们保持沉默，我继续说：\par

朋友们，我现在引用《七十贤士译本》所译的圣经；因为我先前引用你们所拥有的版本时，我试探了你们［以了解］你们的态度如何。因为，提到那说\InnerQuote{祸哉他们！因为他们对自己图谋了恶计，说}\BibleRef{依 3:9}\BibleQuote{让我们除掉义人，因为他令我们厌恶；}的经文时（如《七十贤士译本》所译，我继续道）：\BibleQuote{让我们除掉义人，因为他令我们厌恶；}而在讨论开始时，我补充了你们译本所说的：\BibleQuote{让我们捆绑义人，因为他令我们厌恶。}但你们忙于别的事，似乎听了这些话却没有留意。但现在，因为白昼将尽，太阳快要下山，我要在我所说的话上再加一点，就此结束。我确实已经说过同样的话，但我认为再次加以考虑是适宜的。\par

% structure: p0046 source=h2 target=subsection reason=relative-heading-rank; base=h2

\subsection{第一三八章 诺厄是基督的预像，祂藉水、信德和木头（即十字架）重生我们}

\Person{Justin}：先生们，你们知道，天主曾在依撒意亚中对耶路撒冷说：\BibleQuote{我在诺厄的洪水时代拯救了你。}天主所说的这句话意指救人之奥迹在洪水中显明了。因为义人诺厄，连同洪水中其他凡人，即他的妻子、三个儿子和他们的妻子，共八人，是第八日的象征，基督在复活时显现于第八日，永远在权能上为首。因为基督是\OriginalTerm{一切受造物的首生者}{\GreekText{πρωτότοκος} \GreekText{πάσης} \GreekText{κτίσεως}}，又成为另一族类的首领，这族类是由祂自己藉水、信德和木头重生，其中含有\OriginalTerm{十字架的奥迹}{\GreekText{μυστήριον} \GreekText{τοῦ} \GreekText{σταυροῦ}}；正如诺厄藉木头得救，他带着全家乘木在水上漂流。因此，当先知说\BibleQuote{我在诺厄的日子里拯救了你}，正如我已提到的，他是对同样忠于天主并拥有同样记号的人说的。因为梅瑟手中有杖时，他带领你们的民族穿过大海。而你们相信这只是对你们民族或对土地说的。但全地，如圣经所说，都被淹没，水高过所有山岭十五肘；所以显然这不是对土地说的，而是对那服从祂的人说的：祂也曾为他们预备了耶路撒冷的安息之所，正如先前由洪水的一切象征所表明的；我的意思是，藉水、信德和木头，那些预先预备好、并悔改自己所犯之罪的人，将逃避天主即将来临的审判。\par

% structure: p0048 source=h2 target=subsection reason=relative-heading-rank; base=h2

\subsection{第一三九章 诺厄所宣的祝福和咒诅是对未来的预言}

\Emph{犹斯定：}因为在诺厄的日子，另一奥迹业已完成并预言，只是你们尚未察觉。那就是：诺厄祝福他两个儿子时所倾赐的祝福，以及对他孙子的诅咒。因为预言之神不会诅咒那个已与（他的弟兄们）一同蒙天主祝福的儿子。但由于嘲弄父亲赤身的那个儿子的罪罚将附着于那个儿子整个后裔，他便使诅咒始于他的儿子。于是，他所说的话预言了闪的子孙将占有客纳罕子孙的产业和住处；又预言耶斐特的子孙将夺取闪的子孙从客纳罕子孙手中夺得的产业，并掠夺闪的子孙，正如他们劫掠客纳罕的儿子们一样。请听这事是如何成就的。因为你们这些从闪所出的人，藉天主的旨意侵占了客纳罕子孙的领土，并占有了它。显然，耶斐特的子孙又藉天主的审判入侵了你们，从你们手中夺走你们的土地，并占有了它。正如经上所记载：\BibleQuote{诺厄醒了酒，知道他小儿对他作的事，就说：\InnerQuote{客纳罕应受咒骂，给兄弟当最下贱的奴隶。}又说：\InnerQuote{上主，闪的天主，应受赞美，客纳罕应作他的奴隶。愿天主扩展耶斐特，使他住在闪的帐幕里；客纳罕应作他的奴隶。}}\BibleRef{创 9:24}因此，正如两个民族受祝福——从闪所出的和从耶斐特所出的——并预定闪的子孙首先占有客纳罕的住处，而耶斐特的子孙被预言轮番接受同样的产业，且客纳罕这单一民族被交给这两个民族为奴；同样，基督已来到，带着从全能圣父所赐予的能力，召唤人们进入友谊、祝福、悔改和共居，并应许——正如已证实的——在此同一片土地上，将来所有圣人都将拥有一份产业。因此，所有相信基督、并在祂自己的话及祂先知的话中认识真理的人，无论何处，为奴的或自主的，都知道他们将与祂在那片土地上，并承受永恒不朽的福乐。\par

% structure: p0050 source=h2 target=subsection reason=relative-heading-rank; base=h2

\subsection{第一四〇章　在基督内众人皆得自由。犹太人因是亚巴郎的子孙而妄自期待救恩}

\Emph{犹斯定：}因此，如我先前所说，雅各伯本人是基督的预像，他娶了两位自由妻子的两个婢女，并由此生子，以预示基督同样会接纳耶斐特种族中所有客纳罕的后裔，如同自由人一样，并使这些孩子成为共同继承人。我们就是这样的人；但你们不能理解，因为你们不能饮于天主的活泉，却饮于破裂不能存水的蓄水池，正如圣经所说：\BibleQuote{他们离弃了我这活水的泉源，却为自己掘了蓄水池，破裂不能存水的蓄水池}\BibleRef{耶 2:13}。但它们是破裂不能存水的蓄水池，是你们的教师所挖掘的，正如圣经明确断言：\BibleQuote{他们恭敬我也是假的，因为他们所讲授的教义，是人的规律}\BibleRef{依 29:13}。此外，他们迷惑自己和你们，以为那永远的王国必定赐给那些按血肉属于亚巴郎的散居者，即使他们是罪人、不信者、悖逆天主的人——但圣经已证明并非如此。若是这样，依撒意亚就不会说：\BibleQuote{若不是万军的上主给我们留下些许残余，我们早已像索多玛一样，相似哈摩辣了}\BibleRef{依 1:9}。厄则克耳也说：\BibleQuote{即使诺厄、雅各伯和达尼尔为儿女祈祷，他们的请求也不得应允。但父亲不会因儿子而丧亡，儿子也不会因父亲而丧亡；各人因自己的罪而受罚，各人因自己的正义而得救}\BibleRef{则 18:20}。依撒意亚又说：\BibleQuote{他们必观看那些悖逆者的尸体：他们的虫永不死亡，他们的火永不熄灭；他们将成为一切血肉之躯的鉴戒}\BibleRef{依 66:24}。而且，我们的主，按照那派遣祂者——万有的父和主——的旨意，祂绝不会说：\BibleQuote{从东和西将有人来，与亚巴郎、依撒格和雅各伯一同坐席在天国里；但本国的子民，反要被驱逐到外面的黑暗中}\BibleRef{玛 8:11-12}。此外，我在前面已经证明，那些被预知为不义的人，无论是人还是天使，并不是由于天主的过错而成为邪恶的，而是各人因自己的过错成为他所显出的样子。\par

% structure: p0052 source=h2 target=subsection reason=relative-heading-rank; base=h2

\subsection{第一四一章　论人及天使的自由意志}

\Person{犹斯定}：但为免你们有借口说基督必须被钉十字架，而悖逆者必须是你们民族中的，且事情不能是别样，我预先简要地说：天主愿意人和天使随从祂的旨意，决意创造他们自由地行义；拥有理性，好使他们认识被谁创造，并通过谁他们原不存在，如今存在；且立一法律，使他们若行任何违背正当理性的事，应受祂审判；并且我们自己和天使，将因自己犯的罪而被定罪，除非我们预先悔改。但如果天主的圣言预言有些天使和人必受惩罚，那是因祂预知他们将永恒地[邪恶]，而非因天主如此创造他们。所以，如果他们悔改，凡愿悔改者都能从天主获得仁慈；圣经预言他们有福，说：\BibleQuote{凡主不归罪于其身的人是有福的}即，因悔改其罪，得以从天主获得赦免；并非如你们欺骗自己，以及某些在这事上像你们的人，他们说：即使他们是罪人，但既然认识天主，主将不将罪归咎于他们。我们以此达味的一次跌倒为证，它因他的傲慢而发生，在他如此悲伤哭泣时被赦免，如经上所记。但如果连这样一个人，在悔改之前没有获得赦免，且只有当这位伟大的君王、受傅者、先知悲伤并如此行为，那些不洁净和彻底堕落的人，若他们不哭泣、不悲伤、不悔改，怎能指望主不将罪归咎于他们呢？\par

这达味的一次跌倒，关于乌黎雅妻子的事，诸位，证明圣祖们有多妻，不是为了行奸淫，而是为借由他们实现某种安排和一切奥秘；因为，如果允许娶任何妻子，或任意娶多少妻子，随他所愿，正如你们民族的人在全地，无论他们旅居何处，或被派往何处，以婚姻之名取女子，达味就更有理由这样做了。\par

当我说完这话，最亲爱的\Person{马尔库斯·庞培}，我便结束了。\par

% structure: p0056 source=h2 target=subsection reason=relative-heading-rank; base=h2

\subsection{第一四二章 犹太人致谢并离开犹斯定}

\Person{特里丰}：\Emph{稍作停顿后}你们看到，我们并非有意前来讨论这些论点。我承认，我对这次交谈特别满意；我认为这些人也和我持相同意见。因为我们找到了超出期望、甚至超出可能预期的东西。如果我们可以更频繁地这样做，我们在探究圣经本身方面会得到很大帮助。但既然您即将离开，并且每日期望启航，不要犹豫，在您离去后记我们为友。\par

\Person{犹斯定}：就我而言，如果我留下，我也愿意每天这样做。但现在，既然我预期，赖天主的旨意和助佑，将要启航，我劝你们要在这个为你们自己救恩的极大挣扎中全力以赴，并热切地看重全能天主的基督胜过你们的老师。\par

此后，他们离开了我，祝我航行安全，免于一切不幸。我为他们祈祷，说：\Quote{诸位，我不能为你们祝愿比这更好的事，即：这样认识到智慧赐给每一个人，你们可能与我们持相同意见，并相信耶稣是天主的基督。}\par

\SourceRecord{Translated by Marcus Dods and George Reith.；Ante-Nicene Fathers；Vol. 1.；Edited by Alexander Roberts, James Donaldson, and A. Cleveland Coxe.；Buffalo, NY: Christian Literature Publishing Co.,；1885.；<http://www.newadvent.org/fathers/01289.htm>.}
